% Môn: Vật lý
% Lớp: 12
% Chương: Chương III. Từ trường
% Bài: L12C3 Bài 17. Máy phát điện xoay chiều
% Số câu: 324
% App đọc file này trực tiếp — sửa rồi Commit trên GitHub, bấm Đồng bộ GitHub .tex

% L12C3 Bài 17. Máy phát điện xoay chiều

% ===== Câu 1 =====
% ID: L12C3B17-01-DS
% Mức: TH
\dangbt{Bài toán thực tiễn về vận hành và hiệu suất máy phát điện xoay chiều}
\begin{ex}
% ID: L12C3B17-01-DS
% Mức: TH
Xét Bài toán thực tiễn về vận hành và hiệu suất máy phát điện xoay chiều (Tình huống 1). Hãy xác định tính đúng, sai của bốn phát biểu sau.
\choiceTF
{\True Hiệu suất máy phát được tính bằng công suất điện có ích chia cho công suất cơ đưa vào.}
{Hiệu suất máy phát bằng công suất hao phí chia công suất có ích.}
{\True Khi giải cần xác định đúng dữ kiện, phương chiều, điều kiện áp dụng và đơn vị.}
{Kết quả không phụ thuộc dữ kiện và có thể bỏ qua điều kiện.}
\loigiai{
\begin{itemchoice}
\item (Đúng) Hiệu suất máy phát được tính bằng công suất điện có ích chia cho công suất cơ đưa vào. (Vì): Hiệu suất máy phát được tính bằng công suất điện có ích chia cho công suất cơ đưa vào. b) Sai: Hiệu suất máy phát bằng công suất hao phí chia công suất có ích. c) Đúng: Khi giải cần xác định đúng dữ kiện, phương chiều, điều kiện áp dụng và đơn vị. d) Sai: Kết quả không phụ thuộc dữ kiện và có thể bỏ qua điều kiện. Kiến thức cốt lõi: Hiệu suất máy phát được tính bằng công suất điện có ích chia cho công suất cơ đưa vào
\item (Sai) Hiệu suất máy phát bằng công suất hao phí chia công suất có ích.
\item (Đúng) Khi giải cần xác định đúng dữ kiện, phương chiều, điều kiện áp dụng và đơn vị.
\item (Sai) Kết quả không phụ thuộc dữ kiện và có thể bỏ qua điều kiện.
\end{itemchoice}
}
\end{ex}

% ===== Câu 2 =====
% ID: L12C3B17-01-TL
% Mức: TH
\begin{bt}
% ID: L12C3B17-01-TL
% Mức: TH
Trình bày kiến thức cốt lõi của Bài toán thực tiễn về vận hành và hiệu suất máy phát điện xoay chiều và nêu điều kiện áp dụng.
\loigiai{
Cần nêu được: Hiệu suất máy phát được tính bằng công suất điện có ích chia cho công suất cơ đưa vào. Khi vận dụng phải xác định đúng dữ kiện, phương chiều nếu có, điều kiện áp dụng, hệ đơn vị và kiểm tra tính hợp lí. Nhận định “Hiệu suất máy phát bằng công suất hao phí chia công suất có ích.” sai.
}
\end{bt}

% ===== Câu 3 =====
% ID: L12C3B17-01-TLN
% Mức: TH
\begin{ex}
% ID: L12C3B17-01-TLN
% Mức: TH
Trong bốn phát biểu về Bài toán thực tiễn về vận hành và hiệu suất máy phát điện xoay chiều, có bao nhiêu phát biểu đúng? (Chỉ ghi một số nguyên.) (1) Hiệu suất máy phát được tính bằng công suất điện có ích chia cho công suất cơ đưa vào. (2) Hiệu suất máy phát bằng công suất hao phí chia công suất có ích. (3) Cần kiểm tra điều kiện áp dụng. (4) Có thể bỏ qua đơn vị.
\shortans{2}
\loigiai{
Đối chiếu với kiến thức: Hiệu suất máy phát được tính bằng công suất điện có ích chia cho công suất cơ đưa vào. Có 2 phát biểu đúng.
}
\end{ex}

% ===== Câu 4 =====
% ID: L12C3B17-01-TN
% Mức: NB
\begin{ex}
% ID: L12C3B17-01-TN
% Mức: NB
Phát biểu nào sau đây đúng về Bài toán thực tiễn về vận hành và hiệu suất máy phát điện xoay chiều?
\choice
{Hiệu suất máy phát bằng công suất hao phí chia công suất có ích.}
{Có thể bỏ qua phương, chiều và đơn vị trong mọi trường hợp.}
{Kết quả không phụ thuộc dữ kiện và điều kiện áp dụng.}
{\True Hiệu suất máy phát được tính bằng công suất điện có ích chia cho công suất cơ đưa vào.}
\loigiai{
Hiệu suất máy phát được tính bằng công suất điện có ích chia cho công suất cơ đưa vào. Vì vậy phương án $D$ đúng; các phương án còn lại sai.
}
\end{ex}

% ===== Câu 5 =====
% ID: L12C3B17-02-DS
% Mức: TH
\begin{ex}
% ID: L12C3B17-02-DS
% Mức: TH
Xét Bài toán thực tiễn về vận hành và hiệu suất máy phát điện xoay chiều (Tình huống 2). Hãy xác định tính đúng, sai của bốn phát biểu sau.
\choiceTF
{Hiệu suất máy phát bằng công suất hao phí chia công suất có ích.}
{\True Hiệu suất máy phát được tính bằng công suất điện có ích chia cho công suất cơ đưa vào.}
{Kết quả không phụ thuộc dữ kiện và có thể bỏ qua điều kiện.}
{\True Khi giải cần xác định đúng dữ kiện, phương chiều, điều kiện áp dụng và đơn vị.}
\loigiai{
\begin{itemchoice}
\item (Sai) Hiệu suất máy phát bằng công suất hao phí chia công suất có ích. (Vì): Hiệu suất máy phát bằng công suất hao phí chia công suất có ích. b) Đúng: Hiệu suất máy phát được tính bằng công suất điện có ích chia cho công suất cơ đưa vào. c) Sai: Kết quả không phụ thuộc dữ kiện và có thể bỏ qua điều kiện. d) Đúng: Khi giải cần xác định đúng dữ kiện, phương chiều, điều kiện áp dụng và đơn vị. Kiến thức cốt lõi: Hiệu suất máy phát được tính bằng công suất điện có ích chia cho công suất cơ đưa vào
\item (Đúng) Hiệu suất máy phát được tính bằng công suất điện có ích chia cho công suất cơ đưa vào.
\item (Sai) Kết quả không phụ thuộc dữ kiện và có thể bỏ qua điều kiện.
\item (Đúng) Khi giải cần xác định đúng dữ kiện, phương chiều, điều kiện áp dụng và đơn vị.
\end{itemchoice}
}
\end{ex}

% ===== Câu 6 =====
% ID: L12C3B17-02-TL
% Mức: TH
\begin{bt}
% ID: L12C3B17-02-TL
% Mức: TH
Giải thích vì sao nhận định sau đúng: “Hiệu suất máy phát được tính bằng công suất điện có ích chia cho công suất cơ đưa vào.”
\loigiai{
Cần nêu được: Hiệu suất máy phát được tính bằng công suất điện có ích chia cho công suất cơ đưa vào. Khi vận dụng phải xác định đúng dữ kiện, phương chiều nếu có, điều kiện áp dụng, hệ đơn vị và kiểm tra tính hợp lí. Nhận định “Hiệu suất máy phát bằng công suất hao phí chia công suất có ích.” sai.
}
\end{bt}

% ===== Câu 7 =====
% ID: L12C3B17-02-TLN
% Mức: TH
\begin{ex}
% ID: L12C3B17-02-TLN
% Mức: TH
Trong bốn phát biểu về Bài toán thực tiễn về vận hành và hiệu suất máy phát điện xoay chiều, có bao nhiêu phát biểu đúng? (Chỉ ghi một số nguyên.) (1) Hiệu suất máy phát bằng công suất hao phí chia công suất có ích. (2) Hiệu suất máy phát được tính bằng công suất điện có ích chia cho công suất cơ đưa vào. (3) Kết quả phải phù hợp ý nghĩa vật lí. (4) Mọi đại lượng đều là vô hướng.
\shortans{1}
\loigiai{
Đối chiếu với kiến thức: Hiệu suất máy phát được tính bằng công suất điện có ích chia cho công suất cơ đưa vào. Có 1 phát biểu đúng.
}
\end{ex}

% ===== Câu 8 =====
% ID: L12C3B17-02-TN
% Mức: NB
\begin{ex}
% ID: L12C3B17-02-TN
% Mức: NB
Trong nội dung Bài toán thực tiễn về vận hành và hiệu suất máy phát điện xoay chiều?
\choice
{\True Hiệu suất máy phát được tính bằng công suất điện có ích chia cho công suất cơ đưa vào.}
{Hiệu suất máy phát bằng công suất hao phí chia công suất có ích.}
{Có thể bỏ qua phương, chiều và đơn vị trong mọi trường hợp.}
{Kết quả không phụ thuộc dữ kiện và điều kiện áp dụng.}
\loigiai{
Hiệu suất máy phát được tính bằng công suất điện có ích chia cho công suất cơ đưa vào. Vì vậy phương án $A$ đúng; các phương án còn lại sai.
}
\end{ex}

% ===== Câu 9 =====
% ID: L12C3B17-03-DS
% Mức: VD
\begin{ex}
% ID: L12C3B17-03-DS
% Mức: VD
Xét Bài toán thực tiễn về vận hành và hiệu suất máy phát điện xoay chiều (Tình huống 3). Hãy xác định tính đúng, sai của bốn phát biểu sau.
\choiceTF
{\True Khi giải cần xác định đúng dữ kiện, phương chiều, điều kiện áp dụng và đơn vị.}
{Kết quả không phụ thuộc dữ kiện và có thể bỏ qua điều kiện.}
{\True Hiệu suất máy phát được tính bằng công suất điện có ích chia cho công suất cơ đưa vào.}
{Hiệu suất máy phát bằng công suất hao phí chia công suất có ích.}
\loigiai{
\begin{itemchoice}
\item (Đúng) Khi giải cần xác định đúng dữ kiện, phương chiều, điều kiện áp dụng và đơn vị. (Vì): Khi giải cần xác định đúng dữ kiện, phương chiều, điều kiện áp dụng và đơn vị. b) Sai: Kết quả không phụ thuộc dữ kiện và có thể bỏ qua điều kiện. c) Đúng: Hiệu suất máy phát được tính bằng công suất điện có ích chia cho công suất cơ đưa vào. d) Sai: Hiệu suất máy phát bằng công suất hao phí chia công suất có ích. Kiến thức cốt lõi: Hiệu suất máy phát được tính bằng công suất điện có ích chia cho công suất cơ đưa vào
\item (Sai) Kết quả không phụ thuộc dữ kiện và có thể bỏ qua điều kiện.
\item (Đúng) Hiệu suất máy phát được tính bằng công suất điện có ích chia cho công suất cơ đưa vào.
\item (Sai) Hiệu suất máy phát bằng công suất hao phí chia công suất có ích.
\end{itemchoice}
}
\end{ex}

% ===== Câu 10 =====
% ID: L12C3B17-03-TL
% Mức: VD
\begin{bt}
% ID: L12C3B17-03-TL
% Mức: VD
Phân tích sai lầm trong nhận định: “Hiệu suất máy phát bằng công suất hao phí chia công suất có ích.”
\loigiai{
Cần nêu được: Hiệu suất máy phát được tính bằng công suất điện có ích chia cho công suất cơ đưa vào. Khi vận dụng phải xác định đúng dữ kiện, phương chiều nếu có, điều kiện áp dụng, hệ đơn vị và kiểm tra tính hợp lí. Nhận định “Hiệu suất máy phát bằng công suất hao phí chia công suất có ích.” sai.
}
\end{bt}

% ===== Câu 11 =====
% ID: L12C3B17-03-TLN
% Mức: TH
\begin{ex}
% ID: L12C3B17-03-TLN
% Mức: TH
Trong bốn phát biểu về Bài toán thực tiễn về vận hành và hiệu suất máy phát điện xoay chiều, có bao nhiêu phát biểu đúng? (Chỉ ghi một số nguyên.) (1) Hiệu suất máy phát được tính bằng công suất điện có ích chia cho công suất cơ đưa vào. (2) Hiệu suất máy phát bằng công suất hao phí chia công suất có ích. (3) Cần kiểm tra điều kiện áp dụng. (4) Có thể bỏ qua đơn vị.
\shortans{2}
\loigiai{
Đối chiếu với kiến thức: Hiệu suất máy phát được tính bằng công suất điện có ích chia cho công suất cơ đưa vào. Có 2 phát biểu đúng.
}
\end{ex}

% ===== Câu 12 =====
% ID: L12C3B17-03-TN
% Mức: NB
\begin{ex}
% ID: L12C3B17-03-TN
% Mức: NB
Tình huống 3: chọn kết luận chính xác về Bài toán thực tiễn về vận hành và hiệu suất máy phát điện xoay chiều?
\choice
{Kết quả không phụ thuộc dữ kiện và điều kiện áp dụng.}
{\True Hiệu suất máy phát được tính bằng công suất điện có ích chia cho công suất cơ đưa vào.}
{Hiệu suất máy phát bằng công suất hao phí chia công suất có ích.}
{Có thể bỏ qua phương, chiều và đơn vị trong mọi trường hợp.}
\loigiai{
Hiệu suất máy phát được tính bằng công suất điện có ích chia cho công suất cơ đưa vào. Vì vậy phương án $B$ đúng; các phương án còn lại sai.
}
\end{ex}

% ===== Câu 13 =====
% ID: L12C3B17-04-DS
% Mức: VD
\begin{ex}
% ID: L12C3B17-04-DS
% Mức: VD
Xét Bài toán thực tiễn về vận hành và hiệu suất máy phát điện xoay chiều (Tình huống 4). Hãy xác định tính đúng, sai của bốn phát biểu sau.
\choiceTF
{Kết quả không phụ thuộc dữ kiện và có thể bỏ qua điều kiện.}
{\True Khi giải cần xác định đúng dữ kiện, phương chiều, điều kiện áp dụng và đơn vị.}
{Hiệu suất máy phát bằng công suất hao phí chia công suất có ích.}
{\True Hiệu suất máy phát được tính bằng công suất điện có ích chia cho công suất cơ đưa vào.}
\loigiai{
\begin{itemchoice}
\item (Sai) Kết quả không phụ thuộc dữ kiện và có thể bỏ qua điều kiện. (Vì): Kết quả không phụ thuộc dữ kiện và có thể bỏ qua điều kiện. b) Đúng: Khi giải cần xác định đúng dữ kiện, phương chiều, điều kiện áp dụng và đơn vị. c) Sai: Hiệu suất máy phát bằng công suất hao phí chia công suất có ích. d) Đúng: Hiệu suất máy phát được tính bằng công suất điện có ích chia cho công suất cơ đưa vào. Kiến thức cốt lõi: Hiệu suất máy phát được tính bằng công suất điện có ích chia cho công suất cơ đưa vào
\item (Đúng) Khi giải cần xác định đúng dữ kiện, phương chiều, điều kiện áp dụng và đơn vị.
\item (Sai) Hiệu suất máy phát bằng công suất hao phí chia công suất có ích.
\item (Đúng) Hiệu suất máy phát được tính bằng công suất điện có ích chia cho công suất cơ đưa vào.
\end{itemchoice}
}
\end{ex}

% ===== Câu 14 =====
% ID: L12C3B17-04-TL
% Mức: VD
\begin{bt}
% ID: L12C3B17-04-TL
% Mức: VD
Đề xuất các bước giải một bài toán thuộc dạng Bài toán thực tiễn về vận hành và hiệu suất máy phát điện xoay chiều.
\loigiai{
Cần nêu được: Hiệu suất máy phát được tính bằng công suất điện có ích chia cho công suất cơ đưa vào. Khi vận dụng phải xác định đúng dữ kiện, phương chiều nếu có, điều kiện áp dụng, hệ đơn vị và kiểm tra tính hợp lí. Nhận định “Hiệu suất máy phát bằng công suất hao phí chia công suất có ích.” sai.
}
\end{bt}

% ===== Câu 15 =====
% ID: L12C3B17-04-TLN
% Mức: VD
\begin{ex}
% ID: L12C3B17-04-TLN
% Mức: VD
Trong bốn phát biểu về Bài toán thực tiễn về vận hành và hiệu suất máy phát điện xoay chiều, có bao nhiêu phát biểu đúng? (Chỉ ghi một số nguyên.) (1) Hiệu suất máy phát bằng công suất hao phí chia công suất có ích. (2) Hiệu suất máy phát được tính bằng công suất điện có ích chia cho công suất cơ đưa vào. (3) Kết quả phải phù hợp ý nghĩa vật lí. (4) Mọi đại lượng đều là vô hướng.
\shortans{2}
\loigiai{
Đối chiếu với kiến thức: Hiệu suất máy phát được tính bằng công suất điện có ích chia cho công suất cơ đưa vào. Có 2 phát biểu đúng.
}
\end{ex}

% ===== Câu 16 =====
% ID: L12C3B17-04-TN
% Mức: TH
\begin{ex}
% ID: L12C3B17-04-TN
% Mức: TH
Nội dung nào mô tả đúng nhất Bài toán thực tiễn về vận hành và hiệu suất máy phát điện xoay chiều?
\choice
{Có thể bỏ qua phương, chiều và đơn vị trong mọi trường hợp.}
{Kết quả không phụ thuộc dữ kiện và điều kiện áp dụng.}
{\True Hiệu suất máy phát được tính bằng công suất điện có ích chia cho công suất cơ đưa vào.}
{Hiệu suất máy phát bằng công suất hao phí chia công suất có ích.}
\loigiai{
Hiệu suất máy phát được tính bằng công suất điện có ích chia cho công suất cơ đưa vào. Vì vậy phương án $C$ đúng; các phương án còn lại sai.
}
\end{ex}

% ===== Câu 17 =====
% ID: L12C3B17-05-DS
% Mức: VD
\begin{ex}
% ID: L12C3B17-05-DS
% Mức: VD
Xét Bài toán thực tiễn về vận hành và hiệu suất máy phát điện xoay chiều (Tình huống 5). Hãy xác định tính đúng, sai của bốn phát biểu sau.
\choiceTF
{\True Hiệu suất máy phát được tính bằng công suất điện có ích chia cho công suất cơ đưa vào.}
{Kết quả không phụ thuộc dữ kiện và có thể bỏ qua điều kiện.}
{Hiệu suất máy phát bằng công suất hao phí chia công suất có ích.}
{\True Khi giải cần xác định đúng dữ kiện, phương chiều, điều kiện áp dụng và đơn vị.}
\loigiai{
\begin{itemchoice}
\item (Đúng) Hiệu suất máy phát được tính bằng công suất điện có ích chia cho công suất cơ đưa vào. (Vì): Hiệu suất máy phát được tính bằng công suất điện có ích chia cho công suất cơ đưa vào. b) Sai: Kết quả không phụ thuộc dữ kiện và có thể bỏ qua điều kiện. c) Sai: Hiệu suất máy phát bằng công suất hao phí chia công suất có ích. d) Đúng: Khi giải cần xác định đúng dữ kiện, phương chiều, điều kiện áp dụng và đơn vị. Kiến thức cốt lõi: Hiệu suất máy phát được tính bằng công suất điện có ích chia cho công suất cơ đưa vào
\item (Sai) Kết quả không phụ thuộc dữ kiện và có thể bỏ qua điều kiện.
\item (Sai) Hiệu suất máy phát bằng công suất hao phí chia công suất có ích.
\item (Đúng) Khi giải cần xác định đúng dữ kiện, phương chiều, điều kiện áp dụng và đơn vị.
\end{itemchoice}
}
\end{ex}

% ===== Câu 18 =====
% ID: L12C3B17-05-TL
% Mức: VD
\begin{bt}
% ID: L12C3B17-05-TL
% Mức: VD
Nêu một tình huống thực tiễn liên quan đến Bài toán thực tiễn về vận hành và hiệu suất máy phát điện xoay chiều và giải thích bằng kiến thức Vật lí.
\loigiai{
Cần nêu được: Hiệu suất máy phát được tính bằng công suất điện có ích chia cho công suất cơ đưa vào. Khi vận dụng phải xác định đúng dữ kiện, phương chiều nếu có, điều kiện áp dụng, hệ đơn vị và kiểm tra tính hợp lí. Nhận định “Hiệu suất máy phát bằng công suất hao phí chia công suất có ích.” sai.
}
\end{bt}

% ===== Câu 19 =====
% ID: L12C3B17-05-TLN
% Mức: VD
\begin{ex}
% ID: L12C3B17-05-TLN
% Mức: VD
Trong bốn phát biểu về Bài toán thực tiễn về vận hành và hiệu suất máy phát điện xoay chiều, có bao nhiêu phát biểu đúng? (Chỉ ghi một số nguyên.) (1) Hiệu suất máy phát được tính bằng công suất điện có ích chia cho công suất cơ đưa vào. (2) Hiệu suất máy phát bằng công suất hao phí chia công suất có ích. (3) Cần kiểm tra điều kiện áp dụng. (4) Có thể bỏ qua đơn vị.
\shortans{2}
\loigiai{
Đối chiếu với kiến thức: Hiệu suất máy phát được tính bằng công suất điện có ích chia cho công suất cơ đưa vào. Có 2 phát biểu đúng.
}
\end{ex}

% ===== Câu 20 =====
% ID: L12C3B17-05-TN
% Mức: TH
\begin{ex}
% ID: L12C3B17-05-TN
% Mức: TH
Khi phân tích Bài toán thực tiễn về vận hành và hiệu suất máy phát điện xoay chiều?
\choice
{Hiệu suất máy phát bằng công suất hao phí chia công suất có ích.}
{Có thể bỏ qua phương, chiều và đơn vị trong mọi trường hợp.}
{Kết quả không phụ thuộc dữ kiện và điều kiện áp dụng.}
{\True Hiệu suất máy phát được tính bằng công suất điện có ích chia cho công suất cơ đưa vào.}
\loigiai{
Hiệu suất máy phát được tính bằng công suất điện có ích chia cho công suất cơ đưa vào. Vì vậy phương án $D$ đúng; các phương án còn lại sai.
}
\end{ex}

% ===== Câu 21 =====
% ID: L12C3B17-06-DS
% Mức: TH
\dangbt{Khai thác đồ thị từ thông và suất điện động theo thời gian}
\begin{ex}
% ID: L12C3B17-06-DS
% Mức: TH
Xét Khai thác đồ thị từ thông và suất điện động theo thời gian (Tình huống 1). Hãy xác định tính đúng, sai của bốn phát biểu sau.
\choiceTF
{\True Suất điện động lệch pha pi/2 so với từ thông; khi từ thông cực đại thì suất điện động bằng 0.}
{Từ thông và suất điện động luôn đạt cực đại đồng thời.}
{\True Khi giải cần xác định đúng dữ kiện, phương chiều, điều kiện áp dụng và đơn vị.}
{Kết quả không phụ thuộc dữ kiện và có thể bỏ qua điều kiện.}
\loigiai{
\begin{itemchoice}
\item (Đúng) Suất điện động lệch pha pi/2 so với từ thông; khi từ thông cực đại thì suất điện động bằng 0. (Vì): Suất điện động lệch pha pi/2 so với từ thông; khi từ thông cực đại thì suất điện động bằng 0. b) Sai: Từ thông và suất điện động luôn đạt cực đại đồng thời. c) Đúng: Khi giải cần xác định đúng dữ kiện, phương chiều, điều kiện áp dụng và đơn vị. d) Sai: Kết quả không phụ thuộc dữ kiện và có thể bỏ qua điều kiện. Kiến thức cốt lõi: Suất điện động lệch pha pi/2 so với từ thông; khi từ thông cực đại thì suất điện động bằng 0
\item (Sai) Từ thông và suất điện động luôn đạt cực đại đồng thời.
\item (Đúng) Khi giải cần xác định đúng dữ kiện, phương chiều, điều kiện áp dụng và đơn vị.
\item (Sai) Kết quả không phụ thuộc dữ kiện và có thể bỏ qua điều kiện.
\end{itemchoice}
}
\end{ex}

% ===== Câu 22 =====
% ID: L12C3B17-06-TL
% Mức: VDC
\dangbt{Bài toán thực tiễn về vận hành và hiệu suất máy phát điện xoay chiều}
\begin{bt}
% ID: L12C3B17-06-TL
% Mức: VDC
So sánh nhận định đúng “Hiệu suất máy phát được tính bằng công suất điện có ích chia cho công suất cơ đưa vào.” với nhận định sai “Hiệu suất máy phát bằng công suất hao phí chia công suất có ích.”; chỉ rõ căn cứ Vật lí.
\loigiai{
Cần nêu được: Hiệu suất máy phát được tính bằng công suất điện có ích chia cho công suất cơ đưa vào. Khi vận dụng phải xác định đúng dữ kiện, phương chiều nếu có, điều kiện áp dụng, hệ đơn vị và kiểm tra tính hợp lí. Nhận định “Hiệu suất máy phát bằng công suất hao phí chia công suất có ích.” sai.
}
\end{bt}

% ===== Câu 23 =====
% ID: L12C3B17-06-TLN
% Mức: VDC
\begin{ex}
% ID: L12C3B17-06-TLN
% Mức: VDC
Trong bốn phát biểu về Bài toán thực tiễn về vận hành và hiệu suất máy phát điện xoay chiều, có bao nhiêu phát biểu đúng? (Chỉ ghi một số nguyên.) (1) Khi giải cần xác định đúng dữ kiện, phương chiều, điều kiện áp dụng và đơn vị. (2) Hiệu suất máy phát bằng công suất hao phí chia công suất có ích. (3) Phải dùng đúng định luật cho tình huống. (4) Hiệu suất máy phát được tính bằng công suất điện có ích chia cho công suất cơ đưa vào.
\shortans{3}
\loigiai{
Đối chiếu với kiến thức: Hiệu suất máy phát được tính bằng công suất điện có ích chia cho công suất cơ đưa vào. Có 3 phát biểu đúng.
}
\end{ex}

% ===== Câu 24 =====
% ID: L12C3B17-06-TN
% Mức: TH
\begin{ex}
% ID: L12C3B17-06-TN
% Mức: TH
Một học sinh ôn tập Bài toán thực tiễn về vận hành và hiệu suất máy phát điện xoay chiều?
\choice
{\True Hiệu suất máy phát được tính bằng công suất điện có ích chia cho công suất cơ đưa vào.}
{Hiệu suất máy phát bằng công suất hao phí chia công suất có ích.}
{Có thể bỏ qua phương, chiều và đơn vị trong mọi trường hợp.}
{Kết quả không phụ thuộc dữ kiện và điều kiện áp dụng.}
\loigiai{
Hiệu suất máy phát được tính bằng công suất điện có ích chia cho công suất cơ đưa vào. Vì vậy phương án $A$ đúng; các phương án còn lại sai.
}
\end{ex}

% ===== Câu 25 =====
% ID: L12C3B17-07-DS
% Mức: TH
\dangbt{Khai thác đồ thị từ thông và suất điện động theo thời gian}
\begin{ex}
% ID: L12C3B17-07-DS
% Mức: TH
Xét Khai thác đồ thị từ thông và suất điện động theo thời gian (Tình huống 2). Hãy xác định tính đúng, sai của bốn phát biểu sau.
\choiceTF
{Từ thông và suất điện động luôn đạt cực đại đồng thời.}
{\True Suất điện động lệch pha pi/2 so với từ thông; khi từ thông cực đại thì suất điện động bằng 0.}
{Kết quả không phụ thuộc dữ kiện và có thể bỏ qua điều kiện.}
{\True Khi giải cần xác định đúng dữ kiện, phương chiều, điều kiện áp dụng và đơn vị.}
\loigiai{
\begin{itemchoice}
\item (Sai) Từ thông và suất điện động luôn đạt cực đại đồng thời. (Vì): Từ thông và suất điện động luôn đạt cực đại đồng thời. b) Đúng: Suất điện động lệch pha pi/2 so với từ thông; khi từ thông cực đại thì suất điện động bằng 0. c) Sai: Kết quả không phụ thuộc dữ kiện và có thể bỏ qua điều kiện. d) Đúng: Khi giải cần xác định đúng dữ kiện, phương chiều, điều kiện áp dụng và đơn vị. Kiến thức cốt lõi: Suất điện động lệch pha pi/2 so với từ thông; khi từ thông cực đại thì suất điện động bằng 0
\item (Đúng) Suất điện động lệch pha pi/2 so với từ thông; khi từ thông cực đại thì suất điện động bằng 0.
\item (Sai) Kết quả không phụ thuộc dữ kiện và có thể bỏ qua điều kiện.
\item (Đúng) Khi giải cần xác định đúng dữ kiện, phương chiều, điều kiện áp dụng và đơn vị.
\end{itemchoice}
}
\end{ex}

% ===== Câu 26 =====
% ID: L12C3B17-07-TL
% Mức: TH
\begin{bt}
% ID: L12C3B17-07-TL
% Mức: TH
Trình bày kiến thức cốt lõi của Khai thác đồ thị từ thông và suất điện động theo thời gian và nêu điều kiện áp dụng.
\loigiai{
Cần nêu được: Suất điện động lệch pha pi/2 so với từ thông; khi từ thông cực đại thì suất điện động bằng 0. Khi vận dụng phải xác định đúng dữ kiện, phương chiều nếu có, điều kiện áp dụng, hệ đơn vị và kiểm tra tính hợp lí. Nhận định “Từ thông và suất điện động luôn đạt cực đại đồng thời.” sai.
}
\end{bt}

% ===== Câu 27 =====
% ID: L12C3B17-07-TLN
% Mức: TH
\begin{ex}
% ID: L12C3B17-07-TLN
% Mức: TH
Trong bốn phát biểu về Khai thác đồ thị từ thông và suất điện động theo thời gian, có bao nhiêu phát biểu đúng? (Chỉ ghi một số nguyên.) (1) Suất điện động lệch pha pi/2 so với từ thông; khi từ thông cực đại thì suất điện động bằng 0. (2) Từ thông và suất điện động luôn đạt cực đại đồng thời. (3) Cần kiểm tra điều kiện áp dụng. (4) Có thể bỏ qua đơn vị.
\shortans{2}
\loigiai{
Đối chiếu với kiến thức: Suất điện động lệch pha pi/2 so với từ thông; khi từ thông cực đại thì suất điện động bằng 0. Có 2 phát biểu đúng.
}
\end{ex}

% ===== Câu 28 =====
% ID: L12C3B17-07-TN
% Mức: TH
\dangbt{Bài toán thực tiễn về vận hành và hiệu suất máy phát điện xoay chiều}
\begin{ex}
% ID: L12C3B17-07-TN
% Mức: TH
Chọn phát biểu không trái với kiến thức về Bài toán thực tiễn về vận hành và hiệu suất máy phát điện xoay chiều?
\choice
{Kết quả không phụ thuộc dữ kiện và điều kiện áp dụng.}
{\True Hiệu suất máy phát được tính bằng công suất điện có ích chia cho công suất cơ đưa vào.}
{Hiệu suất máy phát bằng công suất hao phí chia công suất có ích.}
{Có thể bỏ qua phương, chiều và đơn vị trong mọi trường hợp.}
\loigiai{
Hiệu suất máy phát được tính bằng công suất điện có ích chia cho công suất cơ đưa vào. Vì vậy phương án $B$ đúng; các phương án còn lại sai.
}
\end{ex}

% ===== Câu 29 =====
% ID: L12C3B17-08-DS
% Mức: VD
\dangbt{Khai thác đồ thị từ thông và suất điện động theo thời gian}
\begin{ex}
% ID: L12C3B17-08-DS
% Mức: VD
Xét Khai thác đồ thị từ thông và suất điện động theo thời gian (Tình huống 3). Hãy xác định tính đúng, sai của bốn phát biểu sau.
\choiceTF
{\True Khi giải cần xác định đúng dữ kiện, phương chiều, điều kiện áp dụng và đơn vị.}
{Kết quả không phụ thuộc dữ kiện và có thể bỏ qua điều kiện.}
{\True Suất điện động lệch pha pi/2 so với từ thông; khi từ thông cực đại thì suất điện động bằng 0.}
{Từ thông và suất điện động luôn đạt cực đại đồng thời.}
\loigiai{
\begin{itemchoice}
\item (Đúng) Khi giải cần xác định đúng dữ kiện, phương chiều, điều kiện áp dụng và đơn vị. (Vì): Khi giải cần xác định đúng dữ kiện, phương chiều, điều kiện áp dụng và đơn vị. b) Sai: Kết quả không phụ thuộc dữ kiện và có thể bỏ qua điều kiện. c) Đúng: Suất điện động lệch pha pi/2 so với từ thông; khi từ thông cực đại thì suất điện động bằng 0. d) Sai: Từ thông và suất điện động luôn đạt cực đại đồng thời. Kiến thức cốt lõi: Suất điện động lệch pha pi/2 so với từ thông; khi từ thông cực đại thì suất điện động bằng 0
\item (Sai) Kết quả không phụ thuộc dữ kiện và có thể bỏ qua điều kiện.
\item (Đúng) Suất điện động lệch pha pi/2 so với từ thông; khi từ thông cực đại thì suất điện động bằng 0.
\item (Sai) Từ thông và suất điện động luôn đạt cực đại đồng thời.
\end{itemchoice}
}
\end{ex}

% ===== Câu 30 =====
% ID: L12C3B17-08-TL
% Mức: TH
\begin{bt}
% ID: L12C3B17-08-TL
% Mức: TH
Giải thích vì sao nhận định sau đúng: “Suất điện động lệch pha pi/2 so với từ thông; khi từ thông cực đại thì suất điện động bằng 0.”
\loigiai{
Cần nêu được: Suất điện động lệch pha pi/2 so với từ thông; khi từ thông cực đại thì suất điện động bằng 0. Khi vận dụng phải xác định đúng dữ kiện, phương chiều nếu có, điều kiện áp dụng, hệ đơn vị và kiểm tra tính hợp lí. Nhận định “Từ thông và suất điện động luôn đạt cực đại đồng thời.” sai.
}
\end{bt}

% ===== Câu 31 =====
% ID: L12C3B17-08-TLN
% Mức: TH
\begin{ex}
% ID: L12C3B17-08-TLN
% Mức: TH
Trong bốn phát biểu về Khai thác đồ thị từ thông và suất điện động theo thời gian, có bao nhiêu phát biểu đúng? (Chỉ ghi một số nguyên.) (1) Từ thông và suất điện động luôn đạt cực đại đồng thời. (2) Suất điện động lệch pha pi/2 so với từ thông; khi từ thông cực đại thì suất điện động bằng 0. (3) Kết quả phải phù hợp ý nghĩa vật lí. (4) Mọi đại lượng đều là vô hướng.
\shortans{1}
\loigiai{
Đối chiếu với kiến thức: Suất điện động lệch pha pi/2 so với từ thông; khi từ thông cực đại thì suất điện động bằng 0. Có 1 phát biểu đúng.
}
\end{ex}

% ===== Câu 32 =====
% ID: L12C3B17-08-TN
% Mức: VD
\dangbt{Bài toán thực tiễn về vận hành và hiệu suất máy phát điện xoay chiều}
\begin{ex}
% ID: L12C3B17-08-TN
% Mức: VD
Cơ sở Vật lí đúng dùng để giải Bài toán thực tiễn về vận hành và hiệu suất máy phát điện xoay chiều?
\choice
{Có thể bỏ qua phương, chiều và đơn vị trong mọi trường hợp.}
{Kết quả không phụ thuộc dữ kiện và điều kiện áp dụng.}
{\True Hiệu suất máy phát được tính bằng công suất điện có ích chia cho công suất cơ đưa vào.}
{Hiệu suất máy phát bằng công suất hao phí chia công suất có ích.}
\loigiai{
Hiệu suất máy phát được tính bằng công suất điện có ích chia cho công suất cơ đưa vào. Vì vậy phương án $C$ đúng; các phương án còn lại sai.
}
\end{ex}

% ===== Câu 33 =====
% ID: L12C3B17-09-DS
% Mức: VD
\dangbt{Khai thác đồ thị từ thông và suất điện động theo thời gian}
\begin{ex}
% ID: L12C3B17-09-DS
% Mức: VD
Xét Khai thác đồ thị từ thông và suất điện động theo thời gian (Tình huống 4). Hãy xác định tính đúng, sai của bốn phát biểu sau.
\choiceTF
{Kết quả không phụ thuộc dữ kiện và có thể bỏ qua điều kiện.}
{\True Khi giải cần xác định đúng dữ kiện, phương chiều, điều kiện áp dụng và đơn vị.}
{Từ thông và suất điện động luôn đạt cực đại đồng thời.}
{\True Suất điện động lệch pha pi/2 so với từ thông; khi từ thông cực đại thì suất điện động bằng 0.}
\loigiai{
\begin{itemchoice}
\item (Sai) Kết quả không phụ thuộc dữ kiện và có thể bỏ qua điều kiện. (Vì): Kết quả không phụ thuộc dữ kiện và có thể bỏ qua điều kiện. b) Đúng: Khi giải cần xác định đúng dữ kiện, phương chiều, điều kiện áp dụng và đơn vị. c) Sai: Từ thông và suất điện động luôn đạt cực đại đồng thời. d) Đúng: Suất điện động lệch pha pi/2 so với từ thông; khi từ thông cực đại thì suất điện động bằng 0. Kiến thức cốt lõi: Suất điện động lệch pha pi/2 so với từ thông; khi từ thông cực đại thì suất điện động bằng 0
\item (Đúng) Khi giải cần xác định đúng dữ kiện, phương chiều, điều kiện áp dụng và đơn vị.
\item (Sai) Từ thông và suất điện động luôn đạt cực đại đồng thời.
\item (Đúng) Suất điện động lệch pha pi/2 so với từ thông; khi từ thông cực đại thì suất điện động bằng 0.
\end{itemchoice}
}
\end{ex}

% ===== Câu 34 =====
% ID: L12C3B17-09-TL
% Mức: VD
\begin{bt}
% ID: L12C3B17-09-TL
% Mức: VD
Phân tích sai lầm trong nhận định: “Từ thông và suất điện động luôn đạt cực đại đồng thời.”
\loigiai{
Cần nêu được: Suất điện động lệch pha pi/2 so với từ thông; khi từ thông cực đại thì suất điện động bằng 0. Khi vận dụng phải xác định đúng dữ kiện, phương chiều nếu có, điều kiện áp dụng, hệ đơn vị và kiểm tra tính hợp lí. Nhận định “Từ thông và suất điện động luôn đạt cực đại đồng thời.” sai.
}
\end{bt}

% ===== Câu 35 =====
% ID: L12C3B17-09-TLN
% Mức: TH
\begin{ex}
% ID: L12C3B17-09-TLN
% Mức: TH
Trong bốn phát biểu về Khai thác đồ thị từ thông và suất điện động theo thời gian, có bao nhiêu phát biểu đúng? (Chỉ ghi một số nguyên.) (1) Suất điện động lệch pha pi/2 so với từ thông; khi từ thông cực đại thì suất điện động bằng 0. (2) Từ thông và suất điện động luôn đạt cực đại đồng thời. (3) Cần kiểm tra điều kiện áp dụng. (4) Có thể bỏ qua đơn vị.
\shortans{2}
\loigiai{
Đối chiếu với kiến thức: Suất điện động lệch pha pi/2 so với từ thông; khi từ thông cực đại thì suất điện động bằng 0. Có 2 phát biểu đúng.
}
\end{ex}

% ===== Câu 36 =====
% ID: L12C3B17-09-TN
% Mức: VD
\dangbt{Bài toán thực tiễn về vận hành và hiệu suất máy phát điện xoay chiều}
\begin{ex}
% ID: L12C3B17-09-TN
% Mức: VD
Trong một bài thực hành về Bài toán thực tiễn về vận hành và hiệu suất máy phát điện xoay chiều?
\choice
{Hiệu suất máy phát bằng công suất hao phí chia công suất có ích.}
{Có thể bỏ qua phương, chiều và đơn vị trong mọi trường hợp.}
{Kết quả không phụ thuộc dữ kiện và điều kiện áp dụng.}
{\True Hiệu suất máy phát được tính bằng công suất điện có ích chia cho công suất cơ đưa vào.}
\loigiai{
Hiệu suất máy phát được tính bằng công suất điện có ích chia cho công suất cơ đưa vào. Vì vậy phương án $D$ đúng; các phương án còn lại sai.
}
\end{ex}

% ===== Câu 37 =====
% ID: L12C3B17-10-DS
% Mức: VD
\dangbt{Khai thác đồ thị từ thông và suất điện động theo thời gian}
\begin{ex}
% ID: L12C3B17-10-DS
% Mức: VD
Xét Khai thác đồ thị từ thông và suất điện động theo thời gian (Tình huống 5). Hãy xác định tính đúng, sai của bốn phát biểu sau.
\choiceTF
{\True Suất điện động lệch pha pi/2 so với từ thông; khi từ thông cực đại thì suất điện động bằng 0.}
{Kết quả không phụ thuộc dữ kiện và có thể bỏ qua điều kiện.}
{Từ thông và suất điện động luôn đạt cực đại đồng thời.}
{\True Khi giải cần xác định đúng dữ kiện, phương chiều, điều kiện áp dụng và đơn vị.}
\loigiai{
\begin{itemchoice}
\item (Đúng) Suất điện động lệch pha pi/2 so với từ thông; khi từ thông cực đại thì suất điện động bằng 0. (Vì): Suất điện động lệch pha pi/2 so với từ thông; khi từ thông cực đại thì suất điện động bằng 0. b) Sai: Kết quả không phụ thuộc dữ kiện và có thể bỏ qua điều kiện. c) Sai: Từ thông và suất điện động luôn đạt cực đại đồng thời. d) Đúng: Khi giải cần xác định đúng dữ kiện, phương chiều, điều kiện áp dụng và đơn vị. Kiến thức cốt lõi: Suất điện động lệch pha pi/2 so với từ thông; khi từ thông cực đại thì suất điện động bằng 0
\item (Sai) Kết quả không phụ thuộc dữ kiện và có thể bỏ qua điều kiện.
\item (Sai) Từ thông và suất điện động luôn đạt cực đại đồng thời.
\item (Đúng) Khi giải cần xác định đúng dữ kiện, phương chiều, điều kiện áp dụng và đơn vị.
\end{itemchoice}
}
\end{ex}

% ===== Câu 38 =====
% ID: L12C3B17-10-TL
% Mức: VD
\begin{bt}
% ID: L12C3B17-10-TL
% Mức: VD
Đề xuất các bước giải một bài toán thuộc dạng Khai thác đồ thị từ thông và suất điện động theo thời gian.
\loigiai{
Cần nêu được: Suất điện động lệch pha pi/2 so với từ thông; khi từ thông cực đại thì suất điện động bằng 0. Khi vận dụng phải xác định đúng dữ kiện, phương chiều nếu có, điều kiện áp dụng, hệ đơn vị và kiểm tra tính hợp lí. Nhận định “Từ thông và suất điện động luôn đạt cực đại đồng thời.” sai.
}
\end{bt}

% ===== Câu 39 =====
% ID: L12C3B17-10-TLN
% Mức: VD
\begin{ex}
% ID: L12C3B17-10-TLN
% Mức: VD
Trong bốn phát biểu về Khai thác đồ thị từ thông và suất điện động theo thời gian, có bao nhiêu phát biểu đúng? (Chỉ ghi một số nguyên.) (1) Từ thông và suất điện động luôn đạt cực đại đồng thời. (2) Suất điện động lệch pha pi/2 so với từ thông; khi từ thông cực đại thì suất điện động bằng 0. (3) Kết quả phải phù hợp ý nghĩa vật lí. (4) Mọi đại lượng đều là vô hướng.
\shortans{2}
\loigiai{
Đối chiếu với kiến thức: Suất điện động lệch pha pi/2 so với từ thông; khi từ thông cực đại thì suất điện động bằng 0. Có 2 phát biểu đúng.
}
\end{ex}

% ===== Câu 40 =====
% ID: L12C3B17-10-TN
% Mức: VD
\dangbt{Bài toán thực tiễn về vận hành và hiệu suất máy phát điện xoay chiều}
\begin{ex}
% ID: L12C3B17-10-TN
% Mức: VD
Kết luận đúng khi vận dụng Bài toán thực tiễn về vận hành và hiệu suất máy phát điện xoay chiều?
\choice
{\True Hiệu suất máy phát được tính bằng công suất điện có ích chia cho công suất cơ đưa vào.}
{Hiệu suất máy phát bằng công suất hao phí chia công suất có ích.}
{Có thể bỏ qua phương, chiều và đơn vị trong mọi trường hợp.}
{Kết quả không phụ thuộc dữ kiện và điều kiện áp dụng.}
\loigiai{
Hiệu suất máy phát được tính bằng công suất điện có ích chia cho công suất cơ đưa vào. Vì vậy phương án $A$ đúng; các phương án còn lại sai.
}
\end{ex}

% ===== Câu 41 =====
% ID: L12C3B17-100-TN
% Mức: VD
\dangbt{Xác định biên độ, tần số, chu kì và pha của suất điện động xoay chiều}
\begin{ex}
% ID: L12C3B17-100-TN
% Mức: VD
Kết luận đúng khi vận dụng Xác định biên độ, tần số, chu kì và pha của suất điện động xoay chiều?
\choice
{\True Biên độ suất điện động là E0 = NBS omega và tần số f = omega/(2pi).}
{Biên độ suất điện động không phụ thuộc tốc độ quay.}
{Có thể bỏ qua phương, chiều và đơn vị trong mọi trường hợp.}
{Kết quả không phụ thuộc dữ kiện và điều kiện áp dụng.}
\loigiai{
Biên độ suất điện động là E0 = NBS omega và tần số f = omega/(2pi). Vì vậy phương án $A$ đúng; các phương án còn lại sai.
}
\end{ex}

% ===== Câu 42 =====
% ID: L12C3B17-101-TN
% Mức: NB
\dangbt{Khai thác đồ thị từ thông và suất điện động theo thời gian}
\begin{ex}
% ID: L12C3B17-101-TN
% Mức: NB
Phát biểu nào sau đây đúng về Khai thác đồ thị từ thông và suất điện động theo thời gian?
\choice
{Từ thông và suất điện động luôn đạt cực đại đồng thời.}
{Có thể bỏ qua phương, chiều và đơn vị trong mọi trường hợp.}
{Kết quả không phụ thuộc dữ kiện và điều kiện áp dụng.}
{\True Suất điện động lệch pha pi/2 so với từ thông; khi từ thông cực đại thì suất điện động bằng 0.}
\loigiai{
Suất điện động lệch pha pi/2 so với từ thông; khi từ thông cực đại thì suất điện động bằng 0. Vì vậy phương án $D$ đúng; các phương án còn lại sai.
}
\end{ex}

% ===== Câu 43 =====
% ID: L12C3B17-102-TN
% Mức: NB
\begin{ex}
% ID: L12C3B17-102-TN
% Mức: NB
Trong nội dung Khai thác đồ thị từ thông và suất điện động theo thời gian?
\choice
{\True Suất điện động lệch pha pi/2 so với từ thông; khi từ thông cực đại thì suất điện động bằng 0.}
{Từ thông và suất điện động luôn đạt cực đại đồng thời.}
{Có thể bỏ qua phương, chiều và đơn vị trong mọi trường hợp.}
{Kết quả không phụ thuộc dữ kiện và điều kiện áp dụng.}
\loigiai{
Suất điện động lệch pha pi/2 so với từ thông; khi từ thông cực đại thì suất điện động bằng 0. Vì vậy phương án $A$ đúng; các phương án còn lại sai.
}
\end{ex}

% ===== Câu 44 =====
% ID: L12C3B17-103-TN
% Mức: NB
\begin{ex}
% ID: L12C3B17-103-TN
% Mức: NB
Tình huống 3: chọn kết luận chính xác về Khai thác đồ thị từ thông và suất điện động theo thời gian?
\choice
{Kết quả không phụ thuộc dữ kiện và điều kiện áp dụng.}
{\True Suất điện động lệch pha pi/2 so với từ thông; khi từ thông cực đại thì suất điện động bằng 0.}
{Từ thông và suất điện động luôn đạt cực đại đồng thời.}
{Có thể bỏ qua phương, chiều và đơn vị trong mọi trường hợp.}
\loigiai{
Suất điện động lệch pha pi/2 so với từ thông; khi từ thông cực đại thì suất điện động bằng 0. Vì vậy phương án $B$ đúng; các phương án còn lại sai.
}
\end{ex}

% ===== Câu 45 =====
% ID: L12C3B17-104-TN
% Mức: TH
\begin{ex}
% ID: L12C3B17-104-TN
% Mức: TH
Nội dung nào mô tả đúng nhất Khai thác đồ thị từ thông và suất điện động theo thời gian?
\choice
{Có thể bỏ qua phương, chiều và đơn vị trong mọi trường hợp.}
{Kết quả không phụ thuộc dữ kiện và điều kiện áp dụng.}
{\True Suất điện động lệch pha pi/2 so với từ thông; khi từ thông cực đại thì suất điện động bằng 0.}
{Từ thông và suất điện động luôn đạt cực đại đồng thời.}
\loigiai{
Suất điện động lệch pha pi/2 so với từ thông; khi từ thông cực đại thì suất điện động bằng 0. Vì vậy phương án $C$ đúng; các phương án còn lại sai.
}
\end{ex}

% ===== Câu 46 =====
% ID: L12C3B17-105-TN
% Mức: TH
\begin{ex}
% ID: L12C3B17-105-TN
% Mức: TH
Khi phân tích Khai thác đồ thị từ thông và suất điện động theo thời gian?
\choice
{Từ thông và suất điện động luôn đạt cực đại đồng thời.}
{Có thể bỏ qua phương, chiều và đơn vị trong mọi trường hợp.}
{Kết quả không phụ thuộc dữ kiện và điều kiện áp dụng.}
{\True Suất điện động lệch pha pi/2 so với từ thông; khi từ thông cực đại thì suất điện động bằng 0.}
\loigiai{
Suất điện động lệch pha pi/2 so với từ thông; khi từ thông cực đại thì suất điện động bằng 0. Vì vậy phương án $D$ đúng; các phương án còn lại sai.
}
\end{ex}

% ===== Câu 47 =====
% ID: L12C3B17-106-TN
% Mức: TH
\begin{ex}
% ID: L12C3B17-106-TN
% Mức: TH
Một học sinh ôn tập Khai thác đồ thị từ thông và suất điện động theo thời gian?
\choice
{\True Suất điện động lệch pha pi/2 so với từ thông; khi từ thông cực đại thì suất điện động bằng 0.}
{Từ thông và suất điện động luôn đạt cực đại đồng thời.}
{Có thể bỏ qua phương, chiều và đơn vị trong mọi trường hợp.}
{Kết quả không phụ thuộc dữ kiện và điều kiện áp dụng.}
\loigiai{
Suất điện động lệch pha pi/2 so với từ thông; khi từ thông cực đại thì suất điện động bằng 0. Vì vậy phương án $A$ đúng; các phương án còn lại sai.
}
\end{ex}

% ===== Câu 48 =====
% ID: L12C3B17-107-TN
% Mức: TH
\begin{ex}
% ID: L12C3B17-107-TN
% Mức: TH
Chọn phát biểu không trái với kiến thức về Khai thác đồ thị từ thông và suất điện động theo thời gian?
\choice
{Kết quả không phụ thuộc dữ kiện và điều kiện áp dụng.}
{\True Suất điện động lệch pha pi/2 so với từ thông; khi từ thông cực đại thì suất điện động bằng 0.}
{Từ thông và suất điện động luôn đạt cực đại đồng thời.}
{Có thể bỏ qua phương, chiều và đơn vị trong mọi trường hợp.}
\loigiai{
Suất điện động lệch pha pi/2 so với từ thông; khi từ thông cực đại thì suất điện động bằng 0. Vì vậy phương án $B$ đúng; các phương án còn lại sai.
}
\end{ex}

% ===== Câu 49 =====
% ID: L12C3B17-108-TN
% Mức: VD
\begin{ex}
% ID: L12C3B17-108-TN
% Mức: VD
Cơ sở Vật lí đúng dùng để giải Khai thác đồ thị từ thông và suất điện động theo thời gian?
\choice
{Có thể bỏ qua phương, chiều và đơn vị trong mọi trường hợp.}
{Kết quả không phụ thuộc dữ kiện và điều kiện áp dụng.}
{\True Suất điện động lệch pha pi/2 so với từ thông; khi từ thông cực đại thì suất điện động bằng 0.}
{Từ thông và suất điện động luôn đạt cực đại đồng thời.}
\loigiai{
Suất điện động lệch pha pi/2 so với từ thông; khi từ thông cực đại thì suất điện động bằng 0. Vì vậy phương án $C$ đúng; các phương án còn lại sai.
}
\end{ex}

% ===== Câu 50 =====
% ID: L12C3B17-109-TN
% Mức: VD
\begin{ex}
% ID: L12C3B17-109-TN
% Mức: VD
Trong một bài thực hành về Khai thác đồ thị từ thông và suất điện động theo thời gian?
\choice
{Từ thông và suất điện động luôn đạt cực đại đồng thời.}
{Có thể bỏ qua phương, chiều và đơn vị trong mọi trường hợp.}
{Kết quả không phụ thuộc dữ kiện và điều kiện áp dụng.}
{\True Suất điện động lệch pha pi/2 so với từ thông; khi từ thông cực đại thì suất điện động bằng 0.}
\loigiai{
Suất điện động lệch pha pi/2 so với từ thông; khi từ thông cực đại thì suất điện động bằng 0. Vì vậy phương án $D$ đúng; các phương án còn lại sai.
}
\end{ex}

% ===== Câu 51 =====
% ID: L12C3B17-11-DS
% Mức: TH
\dangbt{Nhận biết cấu tạo và nguyên lí của máy phát điện xoay chiều}
\begin{ex}
% ID: L12C3B17-11-DS
% Mức: TH
Xét Nhận biết cấu tạo và nguyên lí của máy phát điện xoay chiều (Tình huống 1). Hãy xác định tính đúng, sai của bốn phát biểu sau.
\choiceTF
{\True Máy phát điện xoay chiều hoạt động dựa trên hiện tượng cảm ứng điện từ khi từ thông qua cuộn dây biến thiên.}
{Máy phát điện xoay chiều biến điện năng thành cơ năng nhờ hiệu ứng nhiệt.}
{\True Khi giải cần xác định đúng dữ kiện, phương chiều, điều kiện áp dụng và đơn vị.}
{Kết quả không phụ thuộc dữ kiện và có thể bỏ qua điều kiện.}
\loigiai{
\begin{itemchoice}
\item (Đúng) Máy phát điện xoay chiều hoạt động dựa trên hiện tượng cảm ứng điện từ khi từ thông qua cuộn dây biến thiên. (Vì): Máy phát điện xoay chiều hoạt động dựa trên hiện tượng cảm ứng điện từ khi từ thông qua cuộn dây biến thiên. b) Sai: Máy phát điện xoay chiều biến điện năng thành cơ năng nhờ hiệu ứng nhiệt. c) Đúng: Khi giải cần xác định đúng dữ kiện, phương chiều, điều kiện áp dụng và đơn vị. d) Sai: Kết quả không phụ thuộc dữ kiện và có thể bỏ qua điều kiện. Kiến thức cốt lõi: Máy phát điện xoay chiều hoạt động dựa trên hiện tượng cảm ứng điện từ khi từ thông qua cuộn dây biến thiên
\item (Sai) Máy phát điện xoay chiều biến điện năng thành cơ năng nhờ hiệu ứng nhiệt.
\item (Đúng) Khi giải cần xác định đúng dữ kiện, phương chiều, điều kiện áp dụng và đơn vị.
\item (Sai) Kết quả không phụ thuộc dữ kiện và có thể bỏ qua điều kiện.
\end{itemchoice}
}
\end{ex}

% ===== Câu 52 =====
% ID: L12C3B17-11-TL
% Mức: VD
\dangbt{Khai thác đồ thị từ thông và suất điện động theo thời gian}
\begin{bt}
% ID: L12C3B17-11-TL
% Mức: VD
Nêu một tình huống thực tiễn liên quan đến Khai thác đồ thị từ thông và suất điện động theo thời gian và giải thích bằng kiến thức Vật lí.
\loigiai{
Cần nêu được: Suất điện động lệch pha pi/2 so với từ thông; khi từ thông cực đại thì suất điện động bằng 0. Khi vận dụng phải xác định đúng dữ kiện, phương chiều nếu có, điều kiện áp dụng, hệ đơn vị và kiểm tra tính hợp lí. Nhận định “Từ thông và suất điện động luôn đạt cực đại đồng thời.” sai.
}
\end{bt}

% ===== Câu 53 =====
% ID: L12C3B17-11-TLN
% Mức: VD
\begin{ex}
% ID: L12C3B17-11-TLN
% Mức: VD
Trong bốn phát biểu về Khai thác đồ thị từ thông và suất điện động theo thời gian, có bao nhiêu phát biểu đúng? (Chỉ ghi một số nguyên.) (1) Suất điện động lệch pha pi/2 so với từ thông; khi từ thông cực đại thì suất điện động bằng 0. (2) Từ thông và suất điện động luôn đạt cực đại đồng thời. (3) Cần kiểm tra điều kiện áp dụng. (4) Có thể bỏ qua đơn vị.
\shortans{2}
\loigiai{
Đối chiếu với kiến thức: Suất điện động lệch pha pi/2 so với từ thông; khi từ thông cực đại thì suất điện động bằng 0. Có 2 phát biểu đúng.
}
\end{ex}

% ===== Câu 54 =====
% ID: L12C3B17-11-TN
% Mức: NB
\begin{ex}
% ID: L12C3B17-11-TN
% Mức: NB
Phát biểu nào sau đây đúng về Khai thác đồ thị từ thông và suất điện động theo thời gian?
\choice
{Từ thông và suất điện động luôn đạt cực đại đồng thời.}
{Có thể bỏ qua phương, chiều và đơn vị trong mọi trường hợp.}
{Kết quả không phụ thuộc dữ kiện và điều kiện áp dụng.}
{\True Suất điện động lệch pha pi/2 so với từ thông; khi từ thông cực đại thì suất điện động bằng 0.}
\loigiai{
Suất điện động lệch pha pi/2 so với từ thông; khi từ thông cực đại thì suất điện động bằng 0. Vì vậy phương án $D$ đúng; các phương án còn lại sai.
}
\end{ex}

% ===== Câu 55 =====
% ID: L12C3B17-110-TN
% Mức: VD
\begin{ex}
% ID: L12C3B17-110-TN
% Mức: VD
Kết luận đúng khi vận dụng Khai thác đồ thị từ thông và suất điện động theo thời gian?
\choice
{\True Suất điện động lệch pha pi/2 so với từ thông; khi từ thông cực đại thì suất điện động bằng 0.}
{Từ thông và suất điện động luôn đạt cực đại đồng thời.}
{Có thể bỏ qua phương, chiều và đơn vị trong mọi trường hợp.}
{Kết quả không phụ thuộc dữ kiện và điều kiện áp dụng.}
\loigiai{
Suất điện động lệch pha pi/2 so với từ thông; khi từ thông cực đại thì suất điện động bằng 0. Vì vậy phương án $A$ đúng; các phương án còn lại sai.
}
\end{ex}

% ===== Câu 56 =====
% ID: L12C3B17-111-TN
% Mức: NB
\dangbt{Bài toán thực tiễn về vận hành và hiệu suất máy phát điện xoay chiều}
\begin{ex}
% ID: L12C3B17-111-TN
% Mức: NB
Phát biểu nào sau đây đúng về Bài toán thực tiễn về vận hành và hiệu suất máy phát điện xoay chiều?
\choice
{Hiệu suất máy phát bằng công suất hao phí chia công suất có ích.}
{Có thể bỏ qua phương, chiều và đơn vị trong mọi trường hợp.}
{Kết quả không phụ thuộc dữ kiện và điều kiện áp dụng.}
{\True Hiệu suất máy phát được tính bằng công suất điện có ích chia cho công suất cơ đưa vào.}
\loigiai{
Hiệu suất máy phát được tính bằng công suất điện có ích chia cho công suất cơ đưa vào. Vì vậy phương án $D$ đúng; các phương án còn lại sai.
}
\end{ex}

% ===== Câu 57 =====
% ID: L12C3B17-112-TN
% Mức: NB
\begin{ex}
% ID: L12C3B17-112-TN
% Mức: NB
Trong nội dung Bài toán thực tiễn về vận hành và hiệu suất máy phát điện xoay chiều?
\choice
{\True Hiệu suất máy phát được tính bằng công suất điện có ích chia cho công suất cơ đưa vào.}
{Hiệu suất máy phát bằng công suất hao phí chia công suất có ích.}
{Có thể bỏ qua phương, chiều và đơn vị trong mọi trường hợp.}
{Kết quả không phụ thuộc dữ kiện và điều kiện áp dụng.}
\loigiai{
Hiệu suất máy phát được tính bằng công suất điện có ích chia cho công suất cơ đưa vào. Vì vậy phương án $A$ đúng; các phương án còn lại sai.
}
\end{ex}

% ===== Câu 58 =====
% ID: L12C3B17-113-TN
% Mức: NB
\begin{ex}
% ID: L12C3B17-113-TN
% Mức: NB
Tình huống 3: chọn kết luận chính xác về Bài toán thực tiễn về vận hành và hiệu suất máy phát điện xoay chiều?
\choice
{Kết quả không phụ thuộc dữ kiện và điều kiện áp dụng.}
{\True Hiệu suất máy phát được tính bằng công suất điện có ích chia cho công suất cơ đưa vào.}
{Hiệu suất máy phát bằng công suất hao phí chia công suất có ích.}
{Có thể bỏ qua phương, chiều và đơn vị trong mọi trường hợp.}
\loigiai{
Hiệu suất máy phát được tính bằng công suất điện có ích chia cho công suất cơ đưa vào. Vì vậy phương án $B$ đúng; các phương án còn lại sai.
}
\end{ex}

% ===== Câu 59 =====
% ID: L12C3B17-114-TN
% Mức: TH
\begin{ex}
% ID: L12C3B17-114-TN
% Mức: TH
Nội dung nào mô tả đúng nhất Bài toán thực tiễn về vận hành và hiệu suất máy phát điện xoay chiều?
\choice
{Có thể bỏ qua phương, chiều và đơn vị trong mọi trường hợp.}
{Kết quả không phụ thuộc dữ kiện và điều kiện áp dụng.}
{\True Hiệu suất máy phát được tính bằng công suất điện có ích chia cho công suất cơ đưa vào.}
{Hiệu suất máy phát bằng công suất hao phí chia công suất có ích.}
\loigiai{
Hiệu suất máy phát được tính bằng công suất điện có ích chia cho công suất cơ đưa vào. Vì vậy phương án $C$ đúng; các phương án còn lại sai.
}
\end{ex}

% ===== Câu 60 =====
% ID: L12C3B17-115-TN
% Mức: TH
\begin{ex}
% ID: L12C3B17-115-TN
% Mức: TH
Khi phân tích Bài toán thực tiễn về vận hành và hiệu suất máy phát điện xoay chiều?
\choice
{Hiệu suất máy phát bằng công suất hao phí chia công suất có ích.}
{Có thể bỏ qua phương, chiều và đơn vị trong mọi trường hợp.}
{Kết quả không phụ thuộc dữ kiện và điều kiện áp dụng.}
{\True Hiệu suất máy phát được tính bằng công suất điện có ích chia cho công suất cơ đưa vào.}
\loigiai{
Hiệu suất máy phát được tính bằng công suất điện có ích chia cho công suất cơ đưa vào. Vì vậy phương án $D$ đúng; các phương án còn lại sai.
}
\end{ex}

% ===== Câu 61 =====
% ID: L12C3B17-116-TN
% Mức: TH
\begin{ex}
% ID: L12C3B17-116-TN
% Mức: TH
Một học sinh ôn tập Bài toán thực tiễn về vận hành và hiệu suất máy phát điện xoay chiều?
\choice
{\True Hiệu suất máy phát được tính bằng công suất điện có ích chia cho công suất cơ đưa vào.}
{Hiệu suất máy phát bằng công suất hao phí chia công suất có ích.}
{Có thể bỏ qua phương, chiều và đơn vị trong mọi trường hợp.}
{Kết quả không phụ thuộc dữ kiện và điều kiện áp dụng.}
\loigiai{
Hiệu suất máy phát được tính bằng công suất điện có ích chia cho công suất cơ đưa vào. Vì vậy phương án $A$ đúng; các phương án còn lại sai.
}
\end{ex}

% ===== Câu 62 =====
% ID: L12C3B17-117-TN
% Mức: TH
\begin{ex}
% ID: L12C3B17-117-TN
% Mức: TH
Chọn phát biểu không trái với kiến thức về Bài toán thực tiễn về vận hành và hiệu suất máy phát điện xoay chiều?
\choice
{Kết quả không phụ thuộc dữ kiện và điều kiện áp dụng.}
{\True Hiệu suất máy phát được tính bằng công suất điện có ích chia cho công suất cơ đưa vào.}
{Hiệu suất máy phát bằng công suất hao phí chia công suất có ích.}
{Có thể bỏ qua phương, chiều và đơn vị trong mọi trường hợp.}
\loigiai{
Hiệu suất máy phát được tính bằng công suất điện có ích chia cho công suất cơ đưa vào. Vì vậy phương án $B$ đúng; các phương án còn lại sai.
}
\end{ex}

% ===== Câu 63 =====
% ID: L12C3B17-118-TN
% Mức: VD
\begin{ex}
% ID: L12C3B17-118-TN
% Mức: VD
Cơ sở Vật lí đúng dùng để giải Bài toán thực tiễn về vận hành và hiệu suất máy phát điện xoay chiều?
\choice
{Có thể bỏ qua phương, chiều và đơn vị trong mọi trường hợp.}
{Kết quả không phụ thuộc dữ kiện và điều kiện áp dụng.}
{\True Hiệu suất máy phát được tính bằng công suất điện có ích chia cho công suất cơ đưa vào.}
{Hiệu suất máy phát bằng công suất hao phí chia công suất có ích.}
\loigiai{
Hiệu suất máy phát được tính bằng công suất điện có ích chia cho công suất cơ đưa vào. Vì vậy phương án $C$ đúng; các phương án còn lại sai.
}
\end{ex}

% ===== Câu 64 =====
% ID: L12C3B17-119-TN
% Mức: VD
\begin{ex}
% ID: L12C3B17-119-TN
% Mức: VD
Trong một bài thực hành về Bài toán thực tiễn về vận hành và hiệu suất máy phát điện xoay chiều?
\choice
{Hiệu suất máy phát bằng công suất hao phí chia công suất có ích.}
{Có thể bỏ qua phương, chiều và đơn vị trong mọi trường hợp.}
{Kết quả không phụ thuộc dữ kiện và điều kiện áp dụng.}
{\True Hiệu suất máy phát được tính bằng công suất điện có ích chia cho công suất cơ đưa vào.}
\loigiai{
Hiệu suất máy phát được tính bằng công suất điện có ích chia cho công suất cơ đưa vào. Vì vậy phương án $D$ đúng; các phương án còn lại sai.
}
\end{ex}

% ===== Câu 65 =====
% ID: L12C3B17-12-DS
% Mức: TH
\dangbt{Nhận biết cấu tạo và nguyên lí của máy phát điện xoay chiều}
\begin{ex}
% ID: L12C3B17-12-DS
% Mức: TH
Xét Nhận biết cấu tạo và nguyên lí của máy phát điện xoay chiều (Tình huống 2). Hãy xác định tính đúng, sai của bốn phát biểu sau.
\choiceTF
{Máy phát điện xoay chiều biến điện năng thành cơ năng nhờ hiệu ứng nhiệt.}
{\True Máy phát điện xoay chiều hoạt động dựa trên hiện tượng cảm ứng điện từ khi từ thông qua cuộn dây biến thiên.}
{Kết quả không phụ thuộc dữ kiện và có thể bỏ qua điều kiện.}
{\True Khi giải cần xác định đúng dữ kiện, phương chiều, điều kiện áp dụng và đơn vị.}
\loigiai{
\begin{itemchoice}
\item (Sai) Máy phát điện xoay chiều biến điện năng thành cơ năng nhờ hiệu ứng nhiệt. (Vì): Máy phát điện xoay chiều biến điện năng thành cơ năng nhờ hiệu ứng nhiệt. b) Đúng: Máy phát điện xoay chiều hoạt động dựa trên hiện tượng cảm ứng điện từ khi từ thông qua cuộn dây biến thiên. c) Sai: Kết quả không phụ thuộc dữ kiện và có thể bỏ qua điều kiện. d) Đúng: Khi giải cần xác định đúng dữ kiện, phương chiều, điều kiện áp dụng và đơn vị. Kiến thức cốt lõi: Máy phát điện xoay chiều hoạt động dựa trên hiện tượng cảm ứng điện từ khi từ thông qua cuộn dây biến thiên
\item (Đúng) Máy phát điện xoay chiều hoạt động dựa trên hiện tượng cảm ứng điện từ khi từ thông qua cuộn dây biến thiên.
\item (Sai) Kết quả không phụ thuộc dữ kiện và có thể bỏ qua điều kiện.
\item (Đúng) Khi giải cần xác định đúng dữ kiện, phương chiều, điều kiện áp dụng và đơn vị.
\end{itemchoice}
}
\end{ex}

% ===== Câu 66 =====
% ID: L12C3B17-12-TL
% Mức: VDC
\dangbt{Khai thác đồ thị từ thông và suất điện động theo thời gian}
\begin{bt}
% ID: L12C3B17-12-TL
% Mức: VDC
So sánh nhận định đúng “Suất điện động lệch pha pi/2 so với từ thông; khi từ thông cực đại thì suất điện động bằng 0.” với nhận định sai “Từ thông và suất điện động luôn đạt cực đại đồng thời.”; chỉ rõ căn cứ Vật lí.
\loigiai{
Cần nêu được: Suất điện động lệch pha pi/2 so với từ thông; khi từ thông cực đại thì suất điện động bằng 0. Khi vận dụng phải xác định đúng dữ kiện, phương chiều nếu có, điều kiện áp dụng, hệ đơn vị và kiểm tra tính hợp lí. Nhận định “Từ thông và suất điện động luôn đạt cực đại đồng thời.” sai.
}
\end{bt}

% ===== Câu 67 =====
% ID: L12C3B17-12-TLN
% Mức: VDC
\begin{ex}
% ID: L12C3B17-12-TLN
% Mức: VDC
Trong bốn phát biểu về Khai thác đồ thị từ thông và suất điện động theo thời gian, có bao nhiêu phát biểu đúng? (Chỉ ghi một số nguyên.) (1) Khi giải cần xác định đúng dữ kiện, phương chiều, điều kiện áp dụng và đơn vị. (2) Từ thông và suất điện động luôn đạt cực đại đồng thời. (3) Phải dùng đúng định luật cho tình huống. (4) Suất điện động lệch pha pi/2 so với từ thông; khi từ thông cực đại thì suất điện động bằng 0.
\shortans{3}
\loigiai{
Đối chiếu với kiến thức: Suất điện động lệch pha pi/2 so với từ thông; khi từ thông cực đại thì suất điện động bằng 0. Có 3 phát biểu đúng.
}
\end{ex}

% ===== Câu 68 =====
% ID: L12C3B17-12-TN
% Mức: NB
\begin{ex}
% ID: L12C3B17-12-TN
% Mức: NB
Trong nội dung Khai thác đồ thị từ thông và suất điện động theo thời gian?
\choice
{\True Suất điện động lệch pha pi/2 so với từ thông; khi từ thông cực đại thì suất điện động bằng 0.}
{Từ thông và suất điện động luôn đạt cực đại đồng thời.}
{Có thể bỏ qua phương, chiều và đơn vị trong mọi trường hợp.}
{Kết quả không phụ thuộc dữ kiện và điều kiện áp dụng.}
\loigiai{
Suất điện động lệch pha pi/2 so với từ thông; khi từ thông cực đại thì suất điện động bằng 0. Vì vậy phương án $A$ đúng; các phương án còn lại sai.
}
\end{ex}

% ===== Câu 69 =====
% ID: L12C3B17-120-TN
% Mức: VD
\dangbt{Bài toán thực tiễn về vận hành và hiệu suất máy phát điện xoay chiều}
\begin{ex}
% ID: L12C3B17-120-TN
% Mức: VD
Kết luận đúng khi vận dụng Bài toán thực tiễn về vận hành và hiệu suất máy phát điện xoay chiều?
\choice
{\True Hiệu suất máy phát được tính bằng công suất điện có ích chia cho công suất cơ đưa vào.}
{Hiệu suất máy phát bằng công suất hao phí chia công suất có ích.}
{Có thể bỏ qua phương, chiều và đơn vị trong mọi trường hợp.}
{Kết quả không phụ thuộc dữ kiện và điều kiện áp dụng.}
\loigiai{
Hiệu suất máy phát được tính bằng công suất điện có ích chia cho công suất cơ đưa vào. Vì vậy phương án $A$ đúng; các phương án còn lại sai.
}
\end{ex}

% ===== Câu 70 =====
% ID: L12C3B17-13-DS
% Mức: VD
\dangbt{Nhận biết cấu tạo và nguyên lí của máy phát điện xoay chiều}
\begin{ex}
% ID: L12C3B17-13-DS
% Mức: VD
Xét Nhận biết cấu tạo và nguyên lí của máy phát điện xoay chiều (Tình huống 3). Hãy xác định tính đúng, sai của bốn phát biểu sau.
\choiceTF
{\True Khi giải cần xác định đúng dữ kiện, phương chiều, điều kiện áp dụng và đơn vị.}
{Kết quả không phụ thuộc dữ kiện và có thể bỏ qua điều kiện.}
{\True Máy phát điện xoay chiều hoạt động dựa trên hiện tượng cảm ứng điện từ khi từ thông qua cuộn dây biến thiên.}
{Máy phát điện xoay chiều biến điện năng thành cơ năng nhờ hiệu ứng nhiệt.}
\loigiai{
\begin{itemchoice}
\item (Đúng) Khi giải cần xác định đúng dữ kiện, phương chiều, điều kiện áp dụng và đơn vị. (Vì): Khi giải cần xác định đúng dữ kiện, phương chiều, điều kiện áp dụng và đơn vị. b) Sai: Kết quả không phụ thuộc dữ kiện và có thể bỏ qua điều kiện. c) Đúng: Máy phát điện xoay chiều hoạt động dựa trên hiện tượng cảm ứng điện từ khi từ thông qua cuộn dây biến thiên. d) Sai: Máy phát điện xoay chiều biến điện năng thành cơ năng nhờ hiệu ứng nhiệt. Kiến thức cốt lõi: Máy phát điện xoay chiều hoạt động dựa trên hiện tượng cảm ứng điện từ khi từ thông qua cuộn dây biến thiên
\item (Sai) Kết quả không phụ thuộc dữ kiện và có thể bỏ qua điều kiện.
\item (Đúng) Máy phát điện xoay chiều hoạt động dựa trên hiện tượng cảm ứng điện từ khi từ thông qua cuộn dây biến thiên.
\item (Sai) Máy phát điện xoay chiều biến điện năng thành cơ năng nhờ hiệu ứng nhiệt.
\end{itemchoice}
}
\end{ex}

% ===== Câu 71 =====
% ID: L12C3B17-13-TL
% Mức: TH
\begin{bt}
% ID: L12C3B17-13-TL
% Mức: TH
Trình bày kiến thức cốt lõi của Nhận biết cấu tạo và nguyên lí của máy phát điện xoay chiều và nêu điều kiện áp dụng.
\loigiai{
Cần nêu được: Máy phát điện xoay chiều hoạt động dựa trên hiện tượng cảm ứng điện từ khi từ thông qua cuộn dây biến thiên. Khi vận dụng phải xác định đúng dữ kiện, phương chiều nếu có, điều kiện áp dụng, hệ đơn vị và kiểm tra tính hợp lí. Nhận định “Máy phát điện xoay chiều biến điện năng thành cơ năng nhờ hiệu ứng nhiệt.” sai.
}
\end{bt}

% ===== Câu 72 =====
% ID: L12C3B17-13-TLN
% Mức: TH
\begin{ex}
% ID: L12C3B17-13-TLN
% Mức: TH
Trong bốn phát biểu về Nhận biết cấu tạo và nguyên lí của máy phát điện xoay chiều, có bao nhiêu phát biểu đúng? (Chỉ ghi một số nguyên.) (1) Máy phát điện xoay chiều hoạt động dựa trên hiện tượng cảm ứng điện từ khi từ thông qua cuộn dây biến thiên. (2) Máy phát điện xoay chiều biến điện năng thành cơ năng nhờ hiệu ứng nhiệt. (3) Cần kiểm tra điều kiện áp dụng. (4) Có thể bỏ qua đơn vị.
\shortans{2}
\loigiai{
Đối chiếu với kiến thức: Máy phát điện xoay chiều hoạt động dựa trên hiện tượng cảm ứng điện từ khi từ thông qua cuộn dây biến thiên. Có 2 phát biểu đúng.
}
\end{ex}

% ===== Câu 73 =====
% ID: L12C3B17-13-TN
% Mức: NB
\dangbt{Khai thác đồ thị từ thông và suất điện động theo thời gian}
\begin{ex}
% ID: L12C3B17-13-TN
% Mức: NB
Tình huống 3: chọn kết luận chính xác về Khai thác đồ thị từ thông và suất điện động theo thời gian?
\choice
{Kết quả không phụ thuộc dữ kiện và điều kiện áp dụng.}
{\True Suất điện động lệch pha pi/2 so với từ thông; khi từ thông cực đại thì suất điện động bằng 0.}
{Từ thông và suất điện động luôn đạt cực đại đồng thời.}
{Có thể bỏ qua phương, chiều và đơn vị trong mọi trường hợp.}
\loigiai{
Suất điện động lệch pha pi/2 so với từ thông; khi từ thông cực đại thì suất điện động bằng 0. Vì vậy phương án $B$ đúng; các phương án còn lại sai.
}
\end{ex}

% ===== Câu 74 =====
% ID: L12C3B17-14-DS
% Mức: VD
\dangbt{Nhận biết cấu tạo và nguyên lí của máy phát điện xoay chiều}
\begin{ex}
% ID: L12C3B17-14-DS
% Mức: VD
Xét Nhận biết cấu tạo và nguyên lí của máy phát điện xoay chiều (Tình huống 4). Hãy xác định tính đúng, sai của bốn phát biểu sau.
\choiceTF
{Kết quả không phụ thuộc dữ kiện và có thể bỏ qua điều kiện.}
{\True Khi giải cần xác định đúng dữ kiện, phương chiều, điều kiện áp dụng và đơn vị.}
{Máy phát điện xoay chiều biến điện năng thành cơ năng nhờ hiệu ứng nhiệt.}
{\True Máy phát điện xoay chiều hoạt động dựa trên hiện tượng cảm ứng điện từ khi từ thông qua cuộn dây biến thiên.}
\loigiai{
\begin{itemchoice}
\item (Sai) Kết quả không phụ thuộc dữ kiện và có thể bỏ qua điều kiện. (Vì): Kết quả không phụ thuộc dữ kiện và có thể bỏ qua điều kiện. b) Đúng: Khi giải cần xác định đúng dữ kiện, phương chiều, điều kiện áp dụng và đơn vị. c) Sai: Máy phát điện xoay chiều biến điện năng thành cơ năng nhờ hiệu ứng nhiệt. d) Đúng: Máy phát điện xoay chiều hoạt động dựa trên hiện tượng cảm ứng điện từ khi từ thông qua cuộn dây biến thiên. Kiến thức cốt lõi: Máy phát điện xoay chiều hoạt động dựa trên hiện tượng cảm ứng điện từ khi từ thông qua cuộn dây biến thiên
\item (Đúng) Khi giải cần xác định đúng dữ kiện, phương chiều, điều kiện áp dụng và đơn vị.
\item (Sai) Máy phát điện xoay chiều biến điện năng thành cơ năng nhờ hiệu ứng nhiệt.
\item (Đúng) Máy phát điện xoay chiều hoạt động dựa trên hiện tượng cảm ứng điện từ khi từ thông qua cuộn dây biến thiên.
\end{itemchoice}
}
\end{ex}

% ===== Câu 75 =====
% ID: L12C3B17-14-TL
% Mức: TH
\begin{bt}
% ID: L12C3B17-14-TL
% Mức: TH
Giải thích vì sao nhận định sau đúng: “Máy phát điện xoay chiều hoạt động dựa trên hiện tượng cảm ứng điện từ khi từ thông qua cuộn dây biến thiên.”
\loigiai{
Cần nêu được: Máy phát điện xoay chiều hoạt động dựa trên hiện tượng cảm ứng điện từ khi từ thông qua cuộn dây biến thiên. Khi vận dụng phải xác định đúng dữ kiện, phương chiều nếu có, điều kiện áp dụng, hệ đơn vị và kiểm tra tính hợp lí. Nhận định “Máy phát điện xoay chiều biến điện năng thành cơ năng nhờ hiệu ứng nhiệt.” sai.
}
\end{bt}

% ===== Câu 76 =====
% ID: L12C3B17-14-TLN
% Mức: TH
\begin{ex}
% ID: L12C3B17-14-TLN
% Mức: TH
Trong bốn phát biểu về Nhận biết cấu tạo và nguyên lí của máy phát điện xoay chiều, có bao nhiêu phát biểu đúng? (Chỉ ghi một số nguyên.) (1) Máy phát điện xoay chiều biến điện năng thành cơ năng nhờ hiệu ứng nhiệt. (2) Máy phát điện xoay chiều hoạt động dựa trên hiện tượng cảm ứng điện từ khi từ thông qua cuộn dây biến thiên. (3) Kết quả phải phù hợp ý nghĩa vật lí. (4) Mọi đại lượng đều là vô hướng.
\shortans{1}
\loigiai{
Đối chiếu với kiến thức: Máy phát điện xoay chiều hoạt động dựa trên hiện tượng cảm ứng điện từ khi từ thông qua cuộn dây biến thiên. Có 1 phát biểu đúng.
}
\end{ex}

% ===== Câu 77 =====
% ID: L12C3B17-14-TN
% Mức: TH
\dangbt{Khai thác đồ thị từ thông và suất điện động theo thời gian}
\begin{ex}
% ID: L12C3B17-14-TN
% Mức: TH
Nội dung nào mô tả đúng nhất Khai thác đồ thị từ thông và suất điện động theo thời gian?
\choice
{Có thể bỏ qua phương, chiều và đơn vị trong mọi trường hợp.}
{Kết quả không phụ thuộc dữ kiện và điều kiện áp dụng.}
{\True Suất điện động lệch pha pi/2 so với từ thông; khi từ thông cực đại thì suất điện động bằng 0.}
{Từ thông và suất điện động luôn đạt cực đại đồng thời.}
\loigiai{
Suất điện động lệch pha pi/2 so với từ thông; khi từ thông cực đại thì suất điện động bằng 0. Vì vậy phương án $C$ đúng; các phương án còn lại sai.
}
\end{ex}

% ===== Câu 78 =====
% ID: L12C3B17-15-DS
% Mức: VD
\dangbt{Nhận biết cấu tạo và nguyên lí của máy phát điện xoay chiều}
\begin{ex}
% ID: L12C3B17-15-DS
% Mức: VD
Xét Nhận biết cấu tạo và nguyên lí của máy phát điện xoay chiều (Tình huống 5). Hãy xác định tính đúng, sai của bốn phát biểu sau.
\choiceTF
{\True Máy phát điện xoay chiều hoạt động dựa trên hiện tượng cảm ứng điện từ khi từ thông qua cuộn dây biến thiên.}
{Kết quả không phụ thuộc dữ kiện và có thể bỏ qua điều kiện.}
{Máy phát điện xoay chiều biến điện năng thành cơ năng nhờ hiệu ứng nhiệt.}
{\True Khi giải cần xác định đúng dữ kiện, phương chiều, điều kiện áp dụng và đơn vị.}
\loigiai{
\begin{itemchoice}
\item (Đúng) Máy phát điện xoay chiều hoạt động dựa trên hiện tượng cảm ứng điện từ khi từ thông qua cuộn dây biến thiên. (Vì): Máy phát điện xoay chiều hoạt động dựa trên hiện tượng cảm ứng điện từ khi từ thông qua cuộn dây biến thiên. b) Sai: Kết quả không phụ thuộc dữ kiện và có thể bỏ qua điều kiện. c) Sai: Máy phát điện xoay chiều biến điện năng thành cơ năng nhờ hiệu ứng nhiệt. d) Đúng: Khi giải cần xác định đúng dữ kiện, phương chiều, điều kiện áp dụng và đơn vị. Kiến thức cốt lõi: Máy phát điện xoay chiều hoạt động dựa trên hiện tượng cảm ứng điện từ khi từ thông qua cuộn dây biến thiên
\item (Sai) Kết quả không phụ thuộc dữ kiện và có thể bỏ qua điều kiện.
\item (Sai) Máy phát điện xoay chiều biến điện năng thành cơ năng nhờ hiệu ứng nhiệt.
\item (Đúng) Khi giải cần xác định đúng dữ kiện, phương chiều, điều kiện áp dụng và đơn vị.
\end{itemchoice}
}
\end{ex}

% ===== Câu 79 =====
% ID: L12C3B17-15-TL
% Mức: VD
\begin{bt}
% ID: L12C3B17-15-TL
% Mức: VD
Phân tích sai lầm trong nhận định: “Máy phát điện xoay chiều biến điện năng thành cơ năng nhờ hiệu ứng nhiệt.”
\loigiai{
Cần nêu được: Máy phát điện xoay chiều hoạt động dựa trên hiện tượng cảm ứng điện từ khi từ thông qua cuộn dây biến thiên. Khi vận dụng phải xác định đúng dữ kiện, phương chiều nếu có, điều kiện áp dụng, hệ đơn vị và kiểm tra tính hợp lí. Nhận định “Máy phát điện xoay chiều biến điện năng thành cơ năng nhờ hiệu ứng nhiệt.” sai.
}
\end{bt}

% ===== Câu 80 =====
% ID: L12C3B17-15-TLN
% Mức: TH
\begin{ex}
% ID: L12C3B17-15-TLN
% Mức: TH
Trong bốn phát biểu về Nhận biết cấu tạo và nguyên lí của máy phát điện xoay chiều, có bao nhiêu phát biểu đúng? (Chỉ ghi một số nguyên.) (1) Máy phát điện xoay chiều hoạt động dựa trên hiện tượng cảm ứng điện từ khi từ thông qua cuộn dây biến thiên. (2) Máy phát điện xoay chiều biến điện năng thành cơ năng nhờ hiệu ứng nhiệt. (3) Cần kiểm tra điều kiện áp dụng. (4) Có thể bỏ qua đơn vị.
\shortans{2}
\loigiai{
Đối chiếu với kiến thức: Máy phát điện xoay chiều hoạt động dựa trên hiện tượng cảm ứng điện từ khi từ thông qua cuộn dây biến thiên. Có 2 phát biểu đúng.
}
\end{ex}

% ===== Câu 81 =====
% ID: L12C3B17-15-TN
% Mức: TH
\dangbt{Khai thác đồ thị từ thông và suất điện động theo thời gian}
\begin{ex}
% ID: L12C3B17-15-TN
% Mức: TH
Khi phân tích Khai thác đồ thị từ thông và suất điện động theo thời gian?
\choice
{Từ thông và suất điện động luôn đạt cực đại đồng thời.}
{Có thể bỏ qua phương, chiều và đơn vị trong mọi trường hợp.}
{Kết quả không phụ thuộc dữ kiện và điều kiện áp dụng.}
{\True Suất điện động lệch pha pi/2 so với từ thông; khi từ thông cực đại thì suất điện động bằng 0.}
\loigiai{
Suất điện động lệch pha pi/2 so với từ thông; khi từ thông cực đại thì suất điện động bằng 0. Vì vậy phương án $D$ đúng; các phương án còn lại sai.
}
\end{ex}

% ===== Câu 82 =====
% ID: L12C3B17-16-DS
% Mức: TH
\dangbt{Thiết lập biểu thức suất điện động xoay chiều do máy phát tạo ra}
\begin{ex}
% ID: L12C3B17-16-DS
% Mức: TH
Xét Thiết lập biểu thức suất điện động xoay chiều do máy phát tạo ra (Tình huống 1). Hãy xác định tính đúng, sai của bốn phát biểu sau.
\choiceTF
{\True Nếu Phi = Phi0 cos(omega t + phi) thì e = omega Phi0 sin(omega t + phi).}
{Suất điện động luôn cùng pha và bằng đúng từ thông.}
{\True Khi giải cần xác định đúng dữ kiện, phương chiều, điều kiện áp dụng và đơn vị.}
{Kết quả không phụ thuộc dữ kiện và có thể bỏ qua điều kiện.}
\loigiai{
\begin{itemchoice}
\item (Đúng) Nếu Phi = Phi0 cos(omega t + phi) thì e = omega Phi0 sin(omega t + phi). (Vì): Nếu Phi = Phi0 cos(omega t + phi) thì e = omega Phi0 sin(omega t + phi). b) Sai: Suất điện động luôn cùng pha và bằng đúng từ thông. c) Đúng: Khi giải cần xác định đúng dữ kiện, phương chiều, điều kiện áp dụng và đơn vị. d) Sai: Kết quả không phụ thuộc dữ kiện và có thể bỏ qua điều kiện. Kiến thức cốt lõi: Nếu Phi = Phi0 cos(omega t + phi) thì e = omega Phi0 sin(omega t + phi
\item (Sai) Suất điện động luôn cùng pha và bằng đúng từ thông.
\item (Đúng) Khi giải cần xác định đúng dữ kiện, phương chiều, điều kiện áp dụng và đơn vị.
\item (Sai) Kết quả không phụ thuộc dữ kiện và có thể bỏ qua điều kiện.
\end{itemchoice}
}
\end{ex}

% ===== Câu 83 =====
% ID: L12C3B17-16-TL
% Mức: VD
\dangbt{Nhận biết cấu tạo và nguyên lí của máy phát điện xoay chiều}
\begin{bt}
% ID: L12C3B17-16-TL
% Mức: VD
Đề xuất các bước giải một bài toán thuộc dạng Nhận biết cấu tạo và nguyên lí của máy phát điện xoay chiều.
\loigiai{
Cần nêu được: Máy phát điện xoay chiều hoạt động dựa trên hiện tượng cảm ứng điện từ khi từ thông qua cuộn dây biến thiên. Khi vận dụng phải xác định đúng dữ kiện, phương chiều nếu có, điều kiện áp dụng, hệ đơn vị và kiểm tra tính hợp lí. Nhận định “Máy phát điện xoay chiều biến điện năng thành cơ năng nhờ hiệu ứng nhiệt.” sai.
}
\end{bt}

% ===== Câu 84 =====
% ID: L12C3B17-16-TLN
% Mức: VD
\begin{ex}
% ID: L12C3B17-16-TLN
% Mức: VD
Trong bốn phát biểu về Nhận biết cấu tạo và nguyên lí của máy phát điện xoay chiều, có bao nhiêu phát biểu đúng? (Chỉ ghi một số nguyên.) (1) Máy phát điện xoay chiều biến điện năng thành cơ năng nhờ hiệu ứng nhiệt. (2) Máy phát điện xoay chiều hoạt động dựa trên hiện tượng cảm ứng điện từ khi từ thông qua cuộn dây biến thiên. (3) Kết quả phải phù hợp ý nghĩa vật lí. (4) Mọi đại lượng đều là vô hướng.
\shortans{2}
\loigiai{
Đối chiếu với kiến thức: Máy phát điện xoay chiều hoạt động dựa trên hiện tượng cảm ứng điện từ khi từ thông qua cuộn dây biến thiên. Có 2 phát biểu đúng.
}
\end{ex}

% ===== Câu 85 =====
% ID: L12C3B17-16-TN
% Mức: TH
\dangbt{Khai thác đồ thị từ thông và suất điện động theo thời gian}
\begin{ex}
% ID: L12C3B17-16-TN
% Mức: TH
Một học sinh ôn tập Khai thác đồ thị từ thông và suất điện động theo thời gian?
\choice
{\True Suất điện động lệch pha pi/2 so với từ thông; khi từ thông cực đại thì suất điện động bằng 0.}
{Từ thông và suất điện động luôn đạt cực đại đồng thời.}
{Có thể bỏ qua phương, chiều và đơn vị trong mọi trường hợp.}
{Kết quả không phụ thuộc dữ kiện và điều kiện áp dụng.}
\loigiai{
Suất điện động lệch pha pi/2 so với từ thông; khi từ thông cực đại thì suất điện động bằng 0. Vì vậy phương án $A$ đúng; các phương án còn lại sai.
}
\end{ex}

% ===== Câu 86 =====
% ID: L12C3B17-17-DS
% Mức: TH
\dangbt{Thiết lập biểu thức suất điện động xoay chiều do máy phát tạo ra}
\begin{ex}
% ID: L12C3B17-17-DS
% Mức: TH
Xét Thiết lập biểu thức suất điện động xoay chiều do máy phát tạo ra (Tình huống 2). Hãy xác định tính đúng, sai của bốn phát biểu sau.
\choiceTF
{Suất điện động luôn cùng pha và bằng đúng từ thông.}
{\True Nếu Phi = Phi0 cos(omega t + phi) thì e = omega Phi0 sin(omega t + phi).}
{Kết quả không phụ thuộc dữ kiện và có thể bỏ qua điều kiện.}
{\True Khi giải cần xác định đúng dữ kiện, phương chiều, điều kiện áp dụng và đơn vị.}
\loigiai{
\begin{itemchoice}
\item (Sai) Suất điện động luôn cùng pha và bằng đúng từ thông. (Vì): uất điện động luôn cùng pha và bằng đúng từ thông. b) Đúng: Nếu Phi = Phi0 cos(omega t + phi) thì e = omega Phi0 sin(omega t + phi). c) Sai: Kết quả không phụ thuộc dữ kiện và có thể bỏ qua điều kiện. d) Đúng: Khi giải cần xác định đúng dữ kiện, phương chiều, điều kiện áp dụng và đơn vị. Kiến thức cốt lõi: Nếu Phi = Phi0 cos(omega t + phi) thì e = omega Phi0 sin(omega t + phi
\item (Đúng) Nếu Phi = Phi0 cos(omega t + phi) thì e = omega Phi0 sin(omega t + phi).
\item (Sai) Kết quả không phụ thuộc dữ kiện và có thể bỏ qua điều kiện.
\item (Đúng) Khi giải cần xác định đúng dữ kiện, phương chiều, điều kiện áp dụng và đơn vị.
\end{itemchoice}
}
\end{ex}

% ===== Câu 87 =====
% ID: L12C3B17-17-TL
% Mức: VD
\dangbt{Nhận biết cấu tạo và nguyên lí của máy phát điện xoay chiều}
\begin{bt}
% ID: L12C3B17-17-TL
% Mức: VD
Nêu một tình huống thực tiễn liên quan đến Nhận biết cấu tạo và nguyên lí của máy phát điện xoay chiều và giải thích bằng kiến thức Vật lí.
\loigiai{
Cần nêu được: Máy phát điện xoay chiều hoạt động dựa trên hiện tượng cảm ứng điện từ khi từ thông qua cuộn dây biến thiên. Khi vận dụng phải xác định đúng dữ kiện, phương chiều nếu có, điều kiện áp dụng, hệ đơn vị và kiểm tra tính hợp lí. Nhận định “Máy phát điện xoay chiều biến điện năng thành cơ năng nhờ hiệu ứng nhiệt.” sai.
}
\end{bt}

% ===== Câu 88 =====
% ID: L12C3B17-17-TLN
% Mức: VD
\begin{ex}
% ID: L12C3B17-17-TLN
% Mức: VD
Trong bốn phát biểu về Nhận biết cấu tạo và nguyên lí của máy phát điện xoay chiều, có bao nhiêu phát biểu đúng? (Chỉ ghi một số nguyên.) (1) Máy phát điện xoay chiều hoạt động dựa trên hiện tượng cảm ứng điện từ khi từ thông qua cuộn dây biến thiên. (2) Máy phát điện xoay chiều biến điện năng thành cơ năng nhờ hiệu ứng nhiệt. (3) Cần kiểm tra điều kiện áp dụng. (4) Có thể bỏ qua đơn vị.
\shortans{2}
\loigiai{
Đối chiếu với kiến thức: Máy phát điện xoay chiều hoạt động dựa trên hiện tượng cảm ứng điện từ khi từ thông qua cuộn dây biến thiên. Có 2 phát biểu đúng.
}
\end{ex}

% ===== Câu 89 =====
% ID: L12C3B17-17-TN
% Mức: TH
\dangbt{Khai thác đồ thị từ thông và suất điện động theo thời gian}
\begin{ex}
% ID: L12C3B17-17-TN
% Mức: TH
Chọn phát biểu không trái với kiến thức về Khai thác đồ thị từ thông và suất điện động theo thời gian?
\choice
{Kết quả không phụ thuộc dữ kiện và điều kiện áp dụng.}
{\True Suất điện động lệch pha pi/2 so với từ thông; khi từ thông cực đại thì suất điện động bằng 0.}
{Từ thông và suất điện động luôn đạt cực đại đồng thời.}
{Có thể bỏ qua phương, chiều và đơn vị trong mọi trường hợp.}
\loigiai{
Suất điện động lệch pha pi/2 so với từ thông; khi từ thông cực đại thì suất điện động bằng 0. Vì vậy phương án $B$ đúng; các phương án còn lại sai.
}
\end{ex}

% ===== Câu 90 =====
% ID: L12C3B17-18-DS
% Mức: VD
\dangbt{Thiết lập biểu thức suất điện động xoay chiều do máy phát tạo ra}
\begin{ex}
% ID: L12C3B17-18-DS
% Mức: VD
Xét Thiết lập biểu thức suất điện động xoay chiều do máy phát tạo ra (Tình huống 3). Hãy xác định tính đúng, sai của bốn phát biểu sau.
\choiceTF
{\True Khi giải cần xác định đúng dữ kiện, phương chiều, điều kiện áp dụng và đơn vị.}
{Kết quả không phụ thuộc dữ kiện và có thể bỏ qua điều kiện.}
{\True Nếu Phi = Phi0 cos(omega t + phi) thì e = omega Phi0 sin(omega t + phi).}
{Suất điện động luôn cùng pha và bằng đúng từ thông.}
\loigiai{
\begin{itemchoice}
\item (Đúng) Khi giải cần xác định đúng dữ kiện, phương chiều, điều kiện áp dụng và đơn vị. (Vì): Khi giải cần xác định đúng dữ kiện, phương chiều, điều kiện áp dụng và đơn vị. b) Sai: Kết quả không phụ thuộc dữ kiện và có thể bỏ qua điều kiện. c) Đúng: Nếu Phi = Phi0 cos(omega t + phi) thì e = omega Phi0 sin(omega t + phi). d) Sai: Suất điện động luôn cùng pha và bằng đúng từ thông. Kiến thức cốt lõi: Nếu Phi = Phi0 cos(omega t + phi) thì e = omega Phi0 sin(omega t + phi
\item (Sai) Kết quả không phụ thuộc dữ kiện và có thể bỏ qua điều kiện.
\item (Đúng) Nếu Phi = Phi0 cos(omega t + phi) thì e = omega Phi0 sin(omega t + phi).
\item (Sai) Suất điện động luôn cùng pha và bằng đúng từ thông.
\end{itemchoice}
}
\end{ex}

% ===== Câu 91 =====
% ID: L12C3B17-18-TL
% Mức: VDC
\dangbt{Nhận biết cấu tạo và nguyên lí của máy phát điện xoay chiều}
\begin{bt}
% ID: L12C3B17-18-TL
% Mức: VDC
So sánh nhận định đúng “Máy phát điện xoay chiều hoạt động dựa trên hiện tượng cảm ứng điện từ khi từ thông qua cuộn dây biến thiên.” với nhận định sai “Máy phát điện xoay chiều biến điện năng thành cơ năng nhờ hiệu ứng nhiệt.”; chỉ rõ căn cứ Vật lí.
\loigiai{
Cần nêu được: Máy phát điện xoay chiều hoạt động dựa trên hiện tượng cảm ứng điện từ khi từ thông qua cuộn dây biến thiên. Khi vận dụng phải xác định đúng dữ kiện, phương chiều nếu có, điều kiện áp dụng, hệ đơn vị và kiểm tra tính hợp lí. Nhận định “Máy phát điện xoay chiều biến điện năng thành cơ năng nhờ hiệu ứng nhiệt.” sai.
}
\end{bt}

% ===== Câu 92 =====
% ID: L12C3B17-18-TLN
% Mức: VDC
\begin{ex}
% ID: L12C3B17-18-TLN
% Mức: VDC
Trong bốn phát biểu về Nhận biết cấu tạo và nguyên lí của máy phát điện xoay chiều, có bao nhiêu phát biểu đúng? (Chỉ ghi một số nguyên.) (1) Khi giải cần xác định đúng dữ kiện, phương chiều, điều kiện áp dụng và đơn vị. (2) Máy phát điện xoay chiều biến điện năng thành cơ năng nhờ hiệu ứng nhiệt. (3) Phải dùng đúng định luật cho tình huống. (4) Máy phát điện xoay chiều hoạt động dựa trên hiện tượng cảm ứng điện từ khi từ thông qua cuộn dây biến thiên.
\shortans{3}
\loigiai{
Đối chiếu với kiến thức: Máy phát điện xoay chiều hoạt động dựa trên hiện tượng cảm ứng điện từ khi từ thông qua cuộn dây biến thiên. Có 3 phát biểu đúng.
}
\end{ex}

% ===== Câu 93 =====
% ID: L12C3B17-18-TN
% Mức: VD
\dangbt{Khai thác đồ thị từ thông và suất điện động theo thời gian}
\begin{ex}
% ID: L12C3B17-18-TN
% Mức: VD
Cơ sở Vật lí đúng dùng để giải Khai thác đồ thị từ thông và suất điện động theo thời gian?
\choice
{Có thể bỏ qua phương, chiều và đơn vị trong mọi trường hợp.}
{Kết quả không phụ thuộc dữ kiện và điều kiện áp dụng.}
{\True Suất điện động lệch pha pi/2 so với từ thông; khi từ thông cực đại thì suất điện động bằng 0.}
{Từ thông và suất điện động luôn đạt cực đại đồng thời.}
\loigiai{
Suất điện động lệch pha pi/2 so với từ thông; khi từ thông cực đại thì suất điện động bằng 0. Vì vậy phương án $C$ đúng; các phương án còn lại sai.
}
\end{ex}

% ===== Câu 94 =====
% ID: L12C3B17-19-DS
% Mức: VD
\dangbt{Thiết lập biểu thức suất điện động xoay chiều do máy phát tạo ra}
\begin{ex}
% ID: L12C3B17-19-DS
% Mức: VD
Xét Thiết lập biểu thức suất điện động xoay chiều do máy phát tạo ra (Tình huống 4). Hãy xác định tính đúng, sai của bốn phát biểu sau.
\choiceTF
{Kết quả không phụ thuộc dữ kiện và có thể bỏ qua điều kiện.}
{\True Khi giải cần xác định đúng dữ kiện, phương chiều, điều kiện áp dụng và đơn vị.}
{Suất điện động luôn cùng pha và bằng đúng từ thông.}
{\True Nếu Phi = Phi0 cos(omega t + phi) thì e = omega Phi0 sin(omega t + phi).}
\loigiai{
\begin{itemchoice}
\item (Sai) Kết quả không phụ thuộc dữ kiện và có thể bỏ qua điều kiện. (Vì): Kết quả không phụ thuộc dữ kiện và có thể bỏ qua điều kiện. b) Đúng: Khi giải cần xác định đúng dữ kiện, phương chiều, điều kiện áp dụng và đơn vị. c) Sai: Suất điện động luôn cùng pha và bằng đúng từ thông. d) Đúng: Nếu Phi = Phi0 cos(omega t + phi) thì e = omega Phi0 sin(omega t + phi). Kiến thức cốt lõi: Nếu Phi = Phi0 cos(omega t + phi) thì e = omega Phi0 sin(omega t + phi
\item (Đúng) Khi giải cần xác định đúng dữ kiện, phương chiều, điều kiện áp dụng và đơn vị.
\item (Sai) Suất điện động luôn cùng pha và bằng đúng từ thông.
\item (Đúng) Nếu Phi = Phi0 cos(omega t + phi) thì e = omega Phi0 sin(omega t + phi).
\end{itemchoice}
}
\end{ex}

% ===== Câu 95 =====
% ID: L12C3B17-19-TL
% Mức: TH
\begin{bt}
% ID: L12C3B17-19-TL
% Mức: TH
Trình bày kiến thức cốt lõi của Thiết lập biểu thức suất điện động xoay chiều do máy phát tạo ra và nêu điều kiện áp dụng.
\loigiai{
Cần nêu được: Nếu Phi = Phi0 cos(omega t + phi) thì e = omega Phi0 sin(omega t + phi). Khi vận dụng phải xác định đúng dữ kiện, phương chiều nếu có, điều kiện áp dụng, hệ đơn vị và kiểm tra tính hợp lí. Nhận định “Suất điện động luôn cùng pha và bằng đúng từ thông.” sai.
}
\end{bt}

% ===== Câu 96 =====
% ID: L12C3B17-19-TLN
% Mức: TH
\begin{ex}
% ID: L12C3B17-19-TLN
% Mức: TH
Trong bốn phát biểu về Thiết lập biểu thức suất điện động xoay chiều do máy phát tạo ra, có bao nhiêu phát biểu đúng? (Chỉ ghi một số nguyên.) (1) Nếu Phi = Phi0 cos(omega t + phi) thì e = omega Phi0 sin(omega t + phi). (2) Suất điện động luôn cùng pha và bằng đúng từ thông. (3) Cần kiểm tra điều kiện áp dụng. (4) Có thể bỏ qua đơn vị.
\shortans{2}
\loigiai{
Đối chiếu với kiến thức: Nếu Phi = Phi0 cos(omega t + phi) thì e = omega Phi0 sin(omega t + phi). Có 2 phát biểu đúng.
}
\end{ex}

% ===== Câu 97 =====
% ID: L12C3B17-19-TN
% Mức: VD
\dangbt{Khai thác đồ thị từ thông và suất điện động theo thời gian}
\begin{ex}
% ID: L12C3B17-19-TN
% Mức: VD
Trong một bài thực hành về Khai thác đồ thị từ thông và suất điện động theo thời gian?
\choice
{Từ thông và suất điện động luôn đạt cực đại đồng thời.}
{Có thể bỏ qua phương, chiều và đơn vị trong mọi trường hợp.}
{Kết quả không phụ thuộc dữ kiện và điều kiện áp dụng.}
{\True Suất điện động lệch pha pi/2 so với từ thông; khi từ thông cực đại thì suất điện động bằng 0.}
\loigiai{
Suất điện động lệch pha pi/2 so với từ thông; khi từ thông cực đại thì suất điện động bằng 0. Vì vậy phương án $D$ đúng; các phương án còn lại sai.
}
\end{ex}

% ===== Câu 98 =====
% ID: L12C3B17-20-DS
% Mức: VD
\dangbt{Thiết lập biểu thức suất điện động xoay chiều do máy phát tạo ra}
\begin{ex}
% ID: L12C3B17-20-DS
% Mức: VD
Xét Thiết lập biểu thức suất điện động xoay chiều do máy phát tạo ra (Tình huống 5). Hãy xác định tính đúng, sai của bốn phát biểu sau.
\choiceTF
{\True Nếu Phi = Phi0 cos(omega t + phi) thì e = omega Phi0 sin(omega t + phi).}
{Kết quả không phụ thuộc dữ kiện và có thể bỏ qua điều kiện.}
{Suất điện động luôn cùng pha và bằng đúng từ thông.}
{\True Khi giải cần xác định đúng dữ kiện, phương chiều, điều kiện áp dụng và đơn vị.}
\loigiai{
\begin{itemchoice}
\item (Đúng) Nếu Phi = Phi0 cos(omega t + phi) thì e = omega Phi0 sin(omega t + phi). (Vì): Nếu Phi = Phi0 cos(omega t + phi) thì e = omega Phi0 sin(omega t + phi). b) Sai: Kết quả không phụ thuộc dữ kiện và có thể bỏ qua điều kiện. c) Sai: Suất điện động luôn cùng pha và bằng đúng từ thông. d) Đúng: Khi giải cần xác định đúng dữ kiện, phương chiều, điều kiện áp dụng và đơn vị. Kiến thức cốt lõi: Nếu Phi = Phi0 cos(omega t + phi) thì e = omega Phi0 sin(omega t + phi
\item (Sai) Kết quả không phụ thuộc dữ kiện và có thể bỏ qua điều kiện.
\item (Sai) Suất điện động luôn cùng pha và bằng đúng từ thông.
\item (Đúng) Khi giải cần xác định đúng dữ kiện, phương chiều, điều kiện áp dụng và đơn vị.
\end{itemchoice}
}
\end{ex}

% ===== Câu 99 =====
% ID: L12C3B17-20-TL
% Mức: TH
\begin{bt}
% ID: L12C3B17-20-TL
% Mức: TH
Giải thích vì sao nhận định sau đúng: “Nếu Phi = Phi0 cos(omega t + phi) thì e = omega Phi0 sin(omega t + phi).”
\loigiai{
Cần nêu được: Nếu Phi = Phi0 cos(omega t + phi) thì e = omega Phi0 sin(omega t + phi). Khi vận dụng phải xác định đúng dữ kiện, phương chiều nếu có, điều kiện áp dụng, hệ đơn vị và kiểm tra tính hợp lí. Nhận định “Suất điện động luôn cùng pha và bằng đúng từ thông.” sai.
}
\end{bt}

% ===== Câu 100 =====
% ID: L12C3B17-20-TLN
% Mức: TH
\begin{ex}
% ID: L12C3B17-20-TLN
% Mức: TH
Trong bốn phát biểu về Thiết lập biểu thức suất điện động xoay chiều do máy phát tạo ra, có bao nhiêu phát biểu đúng? (Chỉ ghi một số nguyên.) (1) Suất điện động luôn cùng pha và bằng đúng từ thông. (2) Nếu Phi = Phi0 cos(omega t + phi) thì e = omega Phi0 sin(omega t + phi). (3) Kết quả phải phù hợp ý nghĩa vật lí. (4) Mọi đại lượng đều là vô hướng.
\shortans{1}
\loigiai{
Đối chiếu với kiến thức: Nếu Phi = Phi0 cos(omega t + phi) thì e = omega Phi0 sin(omega t + phi). Có 1 phát biểu đúng.
}
\end{ex}

% ===== Câu 101 =====
% ID: L12C3B17-20-TN
% Mức: VD
\dangbt{Khai thác đồ thị từ thông và suất điện động theo thời gian}
\begin{ex}
% ID: L12C3B17-20-TN
% Mức: VD
Kết luận đúng khi vận dụng Khai thác đồ thị từ thông và suất điện động theo thời gian?
\choice
{\True Suất điện động lệch pha pi/2 so với từ thông; khi từ thông cực đại thì suất điện động bằng 0.}
{Từ thông và suất điện động luôn đạt cực đại đồng thời.}
{Có thể bỏ qua phương, chiều và đơn vị trong mọi trường hợp.}
{Kết quả không phụ thuộc dữ kiện và điều kiện áp dụng.}
\loigiai{
Suất điện động lệch pha pi/2 so với từ thông; khi từ thông cực đại thì suất điện động bằng 0. Vì vậy phương án $A$ đúng; các phương án còn lại sai.
}
\end{ex}

% ===== Câu 102 =====
% ID: L12C3B17-21-DS
% Mức: TH
\dangbt{Tính từ thông qua khung dây quay trong từ trường}
\begin{ex}
% ID: L12C3B17-21-DS
% Mức: TH
Xét Tính từ thông qua khung dây quay trong từ trường (Tình huống 1). Hãy xác định tính đúng, sai của bốn phát biểu sau.
\choiceTF
{\True Khung $N$ vòng quay đều có từ thông Phi = NBS cos(omega t + phi).}
{Từ thông qua khung quay luôn không đổi và bằng NBS.}
{\True Khi giải cần xác định đúng dữ kiện, phương chiều, điều kiện áp dụng và đơn vị.}
{Kết quả không phụ thuộc dữ kiện và có thể bỏ qua điều kiện.}
\loigiai{
\begin{itemchoice}
\item (Đúng) Khung $N$ vòng quay đều có từ thông Phi = NBS cos(omega t + phi). (Vì): Khung $N$ vòng quay đều có từ thông Phi = NBS cos(omega t + phi). b) Sai: Từ thông qua khung quay luôn không đổi và bằng NBS. c) Đúng: Khi giải cần xác định đúng dữ kiện, phương chiều, điều kiện áp dụng và đơn vị. d) Sai: Kết quả không phụ thuộc dữ kiện và có thể bỏ qua điều kiện. Kiến thức cốt lõi: Khung $N$ vòng quay đều có từ thông Phi = NBS cos(omega t + phi
\item (Sai) Từ thông qua khung quay luôn không đổi và bằng NBS.
\item (Đúng) Khi giải cần xác định đúng dữ kiện, phương chiều, điều kiện áp dụng và đơn vị.
\item (Sai) Kết quả không phụ thuộc dữ kiện và có thể bỏ qua điều kiện.
\end{itemchoice}
}
\end{ex}

% ===== Câu 103 =====
% ID: L12C3B17-21-TL
% Mức: VD
\dangbt{Thiết lập biểu thức suất điện động xoay chiều do máy phát tạo ra}
\begin{bt}
% ID: L12C3B17-21-TL
% Mức: VD
Phân tích sai lầm trong nhận định: “Suất điện động luôn cùng pha và bằng đúng từ thông.”
\loigiai{
Cần nêu được: Nếu Phi = Phi0 cos(omega t + phi) thì e = omega Phi0 sin(omega t + phi). Khi vận dụng phải xác định đúng dữ kiện, phương chiều nếu có, điều kiện áp dụng, hệ đơn vị và kiểm tra tính hợp lí. Nhận định “Suất điện động luôn cùng pha và bằng đúng từ thông.” sai.
}
\end{bt}

% ===== Câu 104 =====
% ID: L12C3B17-21-TLN
% Mức: TH
\begin{ex}
% ID: L12C3B17-21-TLN
% Mức: TH
Trong bốn phát biểu về Thiết lập biểu thức suất điện động xoay chiều do máy phát tạo ra, có bao nhiêu phát biểu đúng? (Chỉ ghi một số nguyên.) (1) Nếu Phi = Phi0 cos(omega t + phi) thì e = omega Phi0 sin(omega t + phi). (2) Suất điện động luôn cùng pha và bằng đúng từ thông. (3) Cần kiểm tra điều kiện áp dụng. (4) Có thể bỏ qua đơn vị.
\shortans{2}
\loigiai{
Đối chiếu với kiến thức: Nếu Phi = Phi0 cos(omega t + phi) thì e = omega Phi0 sin(omega t + phi). Có 2 phát biểu đúng.
}
\end{ex}

% ===== Câu 105 =====
% ID: L12C3B17-21-TN
% Mức: NB
\dangbt{Nhận biết cấu tạo và nguyên lí của máy phát điện xoay chiều}
\begin{ex}
% ID: L12C3B17-21-TN
% Mức: NB
Phát biểu nào sau đây đúng về Nhận biết cấu tạo và nguyên lí của máy phát điện xoay chiều?
\choice
{Máy phát điện xoay chiều biến điện năng thành cơ năng nhờ hiệu ứng nhiệt.}
{Có thể bỏ qua phương, chiều và đơn vị trong mọi trường hợp.}
{Kết quả không phụ thuộc dữ kiện và điều kiện áp dụng.}
{\True Máy phát điện xoay chiều hoạt động dựa trên hiện tượng cảm ứng điện từ khi từ thông qua cuộn dây biến thiên.}
\loigiai{
Máy phát điện xoay chiều hoạt động dựa trên hiện tượng cảm ứng điện từ khi từ thông qua cuộn dây biến thiên. Vì vậy phương án $D$ đúng; các phương án còn lại sai.
}
\end{ex}

% ===== Câu 106 =====
% ID: L12C3B17-22-DS
% Mức: TH
\dangbt{Tính từ thông qua khung dây quay trong từ trường}
\begin{ex}
% ID: L12C3B17-22-DS
% Mức: TH
Xét Tính từ thông qua khung dây quay trong từ trường (Tình huống 2). Hãy xác định tính đúng, sai của bốn phát biểu sau.
\choiceTF
{Từ thông qua khung quay luôn không đổi và bằng NBS.}
{\True Khung $N$ vòng quay đều có từ thông Phi = NBS cos(omega t + phi).}
{Kết quả không phụ thuộc dữ kiện và có thể bỏ qua điều kiện.}
{\True Khi giải cần xác định đúng dữ kiện, phương chiều, điều kiện áp dụng và đơn vị.}
\loigiai{
\begin{itemchoice}
\item (Sai) Từ thông qua khung quay luôn không đổi và bằng NBS. (Vì): Từ thông qua khung quay luôn không đổi và bằng NBS. b) Đúng: Khung $N$ vòng quay đều có từ thông Phi = NBS cos(omega t + phi). c) Sai: Kết quả không phụ thuộc dữ kiện và có thể bỏ qua điều kiện. d) Đúng: Khi giải cần xác định đúng dữ kiện, phương chiều, điều kiện áp dụng và đơn vị. Kiến thức cốt lõi: Khung $N$ vòng quay đều có từ thông Phi = NBS cos(omega t + phi
\item (Đúng) Khung $N$ vòng quay đều có từ thông Phi = NBS cos(omega t + phi).
\item (Sai) Kết quả không phụ thuộc dữ kiện và có thể bỏ qua điều kiện.
\item (Đúng) Khi giải cần xác định đúng dữ kiện, phương chiều, điều kiện áp dụng và đơn vị.
\end{itemchoice}
}
\end{ex}

% ===== Câu 107 =====
% ID: L12C3B17-22-TL
% Mức: VD
\dangbt{Thiết lập biểu thức suất điện động xoay chiều do máy phát tạo ra}
\begin{bt}
% ID: L12C3B17-22-TL
% Mức: VD
Đề xuất các bước giải một bài toán thuộc dạng Thiết lập biểu thức suất điện động xoay chiều do máy phát tạo ra.
\loigiai{
Cần nêu được: Nếu Phi = Phi0 cos(omega t + phi) thì e = omega Phi0 sin(omega t + phi). Khi vận dụng phải xác định đúng dữ kiện, phương chiều nếu có, điều kiện áp dụng, hệ đơn vị và kiểm tra tính hợp lí. Nhận định “Suất điện động luôn cùng pha và bằng đúng từ thông.” sai.
}
\end{bt}

% ===== Câu 108 =====
% ID: L12C3B17-22-TLN
% Mức: VD
\begin{ex}
% ID: L12C3B17-22-TLN
% Mức: VD
Trong bốn phát biểu về Thiết lập biểu thức suất điện động xoay chiều do máy phát tạo ra, có bao nhiêu phát biểu đúng? (Chỉ ghi một số nguyên.) (1) Suất điện động luôn cùng pha và bằng đúng từ thông. (2) Nếu Phi = Phi0 cos(omega t + phi) thì e = omega Phi0 sin(omega t + phi). (3) Kết quả phải phù hợp ý nghĩa vật lí. (4) Mọi đại lượng đều là vô hướng.
\shortans{2}
\loigiai{
Đối chiếu với kiến thức: Nếu Phi = Phi0 cos(omega t + phi) thì e = omega Phi0 sin(omega t + phi). Có 2 phát biểu đúng.
}
\end{ex}

% ===== Câu 109 =====
% ID: L12C3B17-22-TN
% Mức: NB
\dangbt{Nhận biết cấu tạo và nguyên lí của máy phát điện xoay chiều}
\begin{ex}
% ID: L12C3B17-22-TN
% Mức: NB
Trong nội dung Nhận biết cấu tạo và nguyên lí của máy phát điện xoay chiều?
\choice
{\True Máy phát điện xoay chiều hoạt động dựa trên hiện tượng cảm ứng điện từ khi từ thông qua cuộn dây biến thiên.}
{Máy phát điện xoay chiều biến điện năng thành cơ năng nhờ hiệu ứng nhiệt.}
{Có thể bỏ qua phương, chiều và đơn vị trong mọi trường hợp.}
{Kết quả không phụ thuộc dữ kiện và điều kiện áp dụng.}
\loigiai{
Máy phát điện xoay chiều hoạt động dựa trên hiện tượng cảm ứng điện từ khi từ thông qua cuộn dây biến thiên. Vì vậy phương án $A$ đúng; các phương án còn lại sai.
}
\end{ex}

% ===== Câu 110 =====
% ID: L12C3B17-23-DS
% Mức: VD
\dangbt{Tính từ thông qua khung dây quay trong từ trường}
\begin{ex}
% ID: L12C3B17-23-DS
% Mức: VD
Xét Tính từ thông qua khung dây quay trong từ trường (Tình huống 3). Hãy xác định tính đúng, sai của bốn phát biểu sau.
\choiceTF
{\True Khi giải cần xác định đúng dữ kiện, phương chiều, điều kiện áp dụng và đơn vị.}
{Kết quả không phụ thuộc dữ kiện và có thể bỏ qua điều kiện.}
{\True Khung $N$ vòng quay đều có từ thông Phi = NBS cos(omega t + phi).}
{Từ thông qua khung quay luôn không đổi và bằng NBS.}
\loigiai{
\begin{itemchoice}
\item (Đúng) Khi giải cần xác định đúng dữ kiện, phương chiều, điều kiện áp dụng và đơn vị. (Vì): Khi giải cần xác định đúng dữ kiện, phương chiều, điều kiện áp dụng và đơn vị. b) Sai: Kết quả không phụ thuộc dữ kiện và có thể bỏ qua điều kiện. c) Đúng: Khung $N$ vòng quay đều có từ thông Phi = NBS cos(omega t + phi). d) Sai: Từ thông qua khung quay luôn không đổi và bằng NBS. Kiến thức cốt lõi: Khung $N$ vòng quay đều có từ thông Phi = NBS cos(omega t + phi
\item (Sai) Kết quả không phụ thuộc dữ kiện và có thể bỏ qua điều kiện.
\item (Đúng) Khung $N$ vòng quay đều có từ thông Phi = NBS cos(omega t + phi).
\item (Sai) Từ thông qua khung quay luôn không đổi và bằng NBS.
\end{itemchoice}
}
\end{ex}

% ===== Câu 111 =====
% ID: L12C3B17-23-TL
% Mức: VD
\dangbt{Thiết lập biểu thức suất điện động xoay chiều do máy phát tạo ra}
\begin{bt}
% ID: L12C3B17-23-TL
% Mức: VD
Nêu một tình huống thực tiễn liên quan đến Thiết lập biểu thức suất điện động xoay chiều do máy phát tạo ra và giải thích bằng kiến thức Vật lí.
\loigiai{
Cần nêu được: Nếu Phi = Phi0 cos(omega t + phi) thì e = omega Phi0 sin(omega t + phi). Khi vận dụng phải xác định đúng dữ kiện, phương chiều nếu có, điều kiện áp dụng, hệ đơn vị và kiểm tra tính hợp lí. Nhận định “Suất điện động luôn cùng pha và bằng đúng từ thông.” sai.
}
\end{bt}

% ===== Câu 112 =====
% ID: L12C3B17-23-TLN
% Mức: VD
\begin{ex}
% ID: L12C3B17-23-TLN
% Mức: VD
Trong bốn phát biểu về Thiết lập biểu thức suất điện động xoay chiều do máy phát tạo ra, có bao nhiêu phát biểu đúng? (Chỉ ghi một số nguyên.) (1) Nếu Phi = Phi0 cos(omega t + phi) thì e = omega Phi0 sin(omega t + phi). (2) Suất điện động luôn cùng pha và bằng đúng từ thông. (3) Cần kiểm tra điều kiện áp dụng. (4) Có thể bỏ qua đơn vị.
\shortans{2}
\loigiai{
Đối chiếu với kiến thức: Nếu Phi = Phi0 cos(omega t + phi) thì e = omega Phi0 sin(omega t + phi). Có 2 phát biểu đúng.
}
\end{ex}

% ===== Câu 113 =====
% ID: L12C3B17-23-TN
% Mức: NB
\dangbt{Nhận biết cấu tạo và nguyên lí của máy phát điện xoay chiều}
\begin{ex}
% ID: L12C3B17-23-TN
% Mức: NB
Tình huống 3: chọn kết luận chính xác về Nhận biết cấu tạo và nguyên lí của máy phát điện xoay chiều?
\choice
{Kết quả không phụ thuộc dữ kiện và điều kiện áp dụng.}
{\True Máy phát điện xoay chiều hoạt động dựa trên hiện tượng cảm ứng điện từ khi từ thông qua cuộn dây biến thiên.}
{Máy phát điện xoay chiều biến điện năng thành cơ năng nhờ hiệu ứng nhiệt.}
{Có thể bỏ qua phương, chiều và đơn vị trong mọi trường hợp.}
\loigiai{
Máy phát điện xoay chiều hoạt động dựa trên hiện tượng cảm ứng điện từ khi từ thông qua cuộn dây biến thiên. Vì vậy phương án $B$ đúng; các phương án còn lại sai.
}
\end{ex}

% ===== Câu 114 =====
% ID: L12C3B17-24-DS
% Mức: VD
\dangbt{Tính từ thông qua khung dây quay trong từ trường}
\begin{ex}
% ID: L12C3B17-24-DS
% Mức: VD
Xét Tính từ thông qua khung dây quay trong từ trường (Tình huống 4). Hãy xác định tính đúng, sai của bốn phát biểu sau.
\choiceTF
{Kết quả không phụ thuộc dữ kiện và có thể bỏ qua điều kiện.}
{\True Khi giải cần xác định đúng dữ kiện, phương chiều, điều kiện áp dụng và đơn vị.}
{Từ thông qua khung quay luôn không đổi và bằng NBS.}
{\True Khung $N$ vòng quay đều có từ thông Phi = NBS cos(omega t + phi).}
\loigiai{
\begin{itemchoice}
\item (Sai) Kết quả không phụ thuộc dữ kiện và có thể bỏ qua điều kiện. (Vì): Kết quả không phụ thuộc dữ kiện và có thể bỏ qua điều kiện. b) Đúng: Khi giải cần xác định đúng dữ kiện, phương chiều, điều kiện áp dụng và đơn vị. c) Sai: Từ thông qua khung quay luôn không đổi và bằng NBS. d) Đúng: Khung $N$ vòng quay đều có từ thông Phi = NBS cos(omega t + phi). Kiến thức cốt lõi: Khung $N$ vòng quay đều có từ thông Phi = NBS cos(omega t + phi
\item (Đúng) Khi giải cần xác định đúng dữ kiện, phương chiều, điều kiện áp dụng và đơn vị.
\item (Sai) Từ thông qua khung quay luôn không đổi và bằng NBS.
\item (Đúng) Khung $N$ vòng quay đều có từ thông Phi = NBS cos(omega t + phi).
\end{itemchoice}
}
\end{ex}

% ===== Câu 115 =====
% ID: L12C3B17-24-TL
% Mức: VDC
\dangbt{Thiết lập biểu thức suất điện động xoay chiều do máy phát tạo ra}
\begin{bt}
% ID: L12C3B17-24-TL
% Mức: VDC
So sánh nhận định đúng “Nếu Phi = Phi0 cos(omega t + phi) thì e = omega Phi0 sin(omega t + phi).” với nhận định sai “Suất điện động luôn cùng pha và bằng đúng từ thông.”; chỉ rõ căn cứ Vật lí.
\loigiai{
Cần nêu được: Nếu Phi = Phi0 cos(omega t + phi) thì e = omega Phi0 sin(omega t + phi). Khi vận dụng phải xác định đúng dữ kiện, phương chiều nếu có, điều kiện áp dụng, hệ đơn vị và kiểm tra tính hợp lí. Nhận định “Suất điện động luôn cùng pha và bằng đúng từ thông.” sai.
}
\end{bt}

% ===== Câu 116 =====
% ID: L12C3B17-24-TLN
% Mức: VDC
\begin{ex}
% ID: L12C3B17-24-TLN
% Mức: VDC
Trong bốn phát biểu về Thiết lập biểu thức suất điện động xoay chiều do máy phát tạo ra, có bao nhiêu phát biểu đúng? (Chỉ ghi một số nguyên.) (1) Khi giải cần xác định đúng dữ kiện, phương chiều, điều kiện áp dụng và đơn vị. (2) Suất điện động luôn cùng pha và bằng đúng từ thông. (3) Phải dùng đúng định luật cho tình huống. (4) Nếu Phi = Phi0 cos(omega t + phi) thì e = omega Phi0 sin(omega t + phi).
\shortans{3}
\loigiai{
Đối chiếu với kiến thức: Nếu Phi = Phi0 cos(omega t + phi) thì e = omega Phi0 sin(omega t + phi). Có 3 phát biểu đúng.
}
\end{ex}

% ===== Câu 117 =====
% ID: L12C3B17-24-TN
% Mức: TH
\dangbt{Nhận biết cấu tạo và nguyên lí của máy phát điện xoay chiều}
\begin{ex}
% ID: L12C3B17-24-TN
% Mức: TH
Nội dung nào mô tả đúng nhất Nhận biết cấu tạo và nguyên lí của máy phát điện xoay chiều?
\choice
{Có thể bỏ qua phương, chiều và đơn vị trong mọi trường hợp.}
{Kết quả không phụ thuộc dữ kiện và điều kiện áp dụng.}
{\True Máy phát điện xoay chiều hoạt động dựa trên hiện tượng cảm ứng điện từ khi từ thông qua cuộn dây biến thiên.}
{Máy phát điện xoay chiều biến điện năng thành cơ năng nhờ hiệu ứng nhiệt.}
\loigiai{
Máy phát điện xoay chiều hoạt động dựa trên hiện tượng cảm ứng điện từ khi từ thông qua cuộn dây biến thiên. Vì vậy phương án $C$ đúng; các phương án còn lại sai.
}
\end{ex}

% ===== Câu 118 =====
% ID: L12C3B17-25-DS
% Mức: VD
\dangbt{Tính từ thông qua khung dây quay trong từ trường}
\begin{ex}
% ID: L12C3B17-25-DS
% Mức: VD
Xét Tính từ thông qua khung dây quay trong từ trường (Tình huống 5). Hãy xác định tính đúng, sai của bốn phát biểu sau.
\choiceTF
{\True Khung $N$ vòng quay đều có từ thông Phi = NBS cos(omega t + phi).}
{Kết quả không phụ thuộc dữ kiện và có thể bỏ qua điều kiện.}
{Từ thông qua khung quay luôn không đổi và bằng NBS.}
{\True Khi giải cần xác định đúng dữ kiện, phương chiều, điều kiện áp dụng và đơn vị.}
\loigiai{
\begin{itemchoice}
\item (Đúng) Khung $N$ vòng quay đều có từ thông Phi = NBS cos(omega t + phi). (Vì): Khung $N$ vòng quay đều có từ thông Phi = NBS cos(omega t + phi). b) Sai: Kết quả không phụ thuộc dữ kiện và có thể bỏ qua điều kiện. c) Sai: Từ thông qua khung quay luôn không đổi và bằng NBS. d) Đúng: Khi giải cần xác định đúng dữ kiện, phương chiều, điều kiện áp dụng và đơn vị. Kiến thức cốt lõi: Khung $N$ vòng quay đều có từ thông Phi = NBS cos(omega t + phi
\item (Sai) Kết quả không phụ thuộc dữ kiện và có thể bỏ qua điều kiện.
\item (Sai) Từ thông qua khung quay luôn không đổi và bằng NBS.
\item (Đúng) Khi giải cần xác định đúng dữ kiện, phương chiều, điều kiện áp dụng và đơn vị.
\end{itemchoice}
}
\end{ex}

% ===== Câu 119 =====
% ID: L12C3B17-25-TL
% Mức: TH
\begin{bt}
% ID: L12C3B17-25-TL
% Mức: TH
Trình bày kiến thức cốt lõi của Tính từ thông qua khung dây quay trong từ trường và nêu điều kiện áp dụng.
\loigiai{
Cần nêu được: Khung $N$ vòng quay đều có từ thông Phi = NBS cos(omega t + phi). Khi vận dụng phải xác định đúng dữ kiện, phương chiều nếu có, điều kiện áp dụng, hệ đơn vị và kiểm tra tính hợp lí. Nhận định “Từ thông qua khung quay luôn không đổi và bằng NBS.” sai.
}
\end{bt}

% ===== Câu 120 =====
% ID: L12C3B17-25-TLN
% Mức: TH
\begin{ex}
% ID: L12C3B17-25-TLN
% Mức: TH
Trong bốn phát biểu về Tính từ thông qua khung dây quay trong từ trường, có bao nhiêu phát biểu đúng? (Chỉ ghi một số nguyên.) (1) Khung $N$ vòng quay đều có từ thông Phi = NBS cos(omega t + phi). (2) Từ thông qua khung quay luôn không đổi và bằng NBS. (3) Cần kiểm tra điều kiện áp dụng. (4) Có thể bỏ qua đơn vị.
\shortans{2}
\loigiai{
Đối chiếu với kiến thức: Khung $N$ vòng quay đều có từ thông Phi = NBS cos(omega t + phi). Có 2 phát biểu đúng.
}
\end{ex}

% ===== Câu 121 =====
% ID: L12C3B17-25-TN
% Mức: TH
\dangbt{Nhận biết cấu tạo và nguyên lí của máy phát điện xoay chiều}
\begin{ex}
% ID: L12C3B17-25-TN
% Mức: TH
Khi phân tích Nhận biết cấu tạo và nguyên lí của máy phát điện xoay chiều?
\choice
{Máy phát điện xoay chiều biến điện năng thành cơ năng nhờ hiệu ứng nhiệt.}
{Có thể bỏ qua phương, chiều và đơn vị trong mọi trường hợp.}
{Kết quả không phụ thuộc dữ kiện và điều kiện áp dụng.}
{\True Máy phát điện xoay chiều hoạt động dựa trên hiện tượng cảm ứng điện từ khi từ thông qua cuộn dây biến thiên.}
\loigiai{
Máy phát điện xoay chiều hoạt động dựa trên hiện tượng cảm ứng điện từ khi từ thông qua cuộn dây biến thiên. Vì vậy phương án $D$ đúng; các phương án còn lại sai.
}
\end{ex}

% ===== Câu 122 =====
% ID: L12C3B17-26-DS
% Mức: TH
\dangbt{Xác định biên độ, tần số, chu kì và pha của suất điện động xoay chiều}
\begin{ex}
% ID: L12C3B17-26-DS
% Mức: TH
Xét Xác định biên độ, tần số, chu kì và pha của suất điện động xoay chiều (Tình huống 1). Hãy xác định tính đúng, sai của bốn phát biểu sau.
\choiceTF
{\True Biên độ suất điện động là E0 = NBS omega và tần số f = omega/(2pi).}
{Biên độ suất điện động không phụ thuộc tốc độ quay.}
{\True Khi giải cần xác định đúng dữ kiện, phương chiều, điều kiện áp dụng và đơn vị.}
{Kết quả không phụ thuộc dữ kiện và có thể bỏ qua điều kiện.}
\loigiai{
\begin{itemchoice}
\item (Đúng) Biên độ suất điện động là E0 = NBS omega và tần số f = omega/(2pi). (Vì): Biên độ suất điện động là E0 = NBS omega và tần số f = omega/(2pi). b) Sai: Biên độ suất điện động không phụ thuộc tốc độ quay. c) Đúng: Khi giải cần xác định đúng dữ kiện, phương chiều, điều kiện áp dụng và đơn vị. d) Sai: Kết quả không phụ thuộc dữ kiện và có thể bỏ qua điều kiện. Kiến thức cốt lõi: Biên độ suất điện động là E0 = NBS omega và tần số f = omega/(2pi
\item (Sai) Biên độ suất điện động không phụ thuộc tốc độ quay.
\item (Đúng) Khi giải cần xác định đúng dữ kiện, phương chiều, điều kiện áp dụng và đơn vị.
\item (Sai) Kết quả không phụ thuộc dữ kiện và có thể bỏ qua điều kiện.
\end{itemchoice}
}
\end{ex}

% ===== Câu 123 =====
% ID: L12C3B17-26-TL
% Mức: TH
\dangbt{Tính từ thông qua khung dây quay trong từ trường}
\begin{bt}
% ID: L12C3B17-26-TL
% Mức: TH
Giải thích vì sao nhận định sau đúng: “Khung $N$ vòng quay đều có từ thông Phi = NBS cos(omega t + phi).”
\loigiai{
Cần nêu được: Khung $N$ vòng quay đều có từ thông Phi = NBS cos(omega t + phi). Khi vận dụng phải xác định đúng dữ kiện, phương chiều nếu có, điều kiện áp dụng, hệ đơn vị và kiểm tra tính hợp lí. Nhận định “Từ thông qua khung quay luôn không đổi và bằng NBS.” sai.
}
\end{bt}

% ===== Câu 124 =====
% ID: L12C3B17-26-TLN
% Mức: TH
\begin{ex}
% ID: L12C3B17-26-TLN
% Mức: TH
Trong bốn phát biểu về Tính từ thông qua khung dây quay trong từ trường, có bao nhiêu phát biểu đúng? (Chỉ ghi một số nguyên.) (1) Từ thông qua khung quay luôn không đổi và bằng NBS. (2) Khung $N$ vòng quay đều có từ thông Phi = NBS cos(omega t + phi). (3) Kết quả phải phù hợp ý nghĩa vật lí. (4) Mọi đại lượng đều là vô hướng.
\shortans{1}
\loigiai{
Đối chiếu với kiến thức: Khung $N$ vòng quay đều có từ thông Phi = NBS cos(omega t + phi). Có 1 phát biểu đúng.
}
\end{ex}

% ===== Câu 125 =====
% ID: L12C3B17-26-TN
% Mức: TH
\dangbt{Nhận biết cấu tạo và nguyên lí của máy phát điện xoay chiều}
\begin{ex}
% ID: L12C3B17-26-TN
% Mức: TH
Một học sinh ôn tập Nhận biết cấu tạo và nguyên lí của máy phát điện xoay chiều?
\choice
{\True Máy phát điện xoay chiều hoạt động dựa trên hiện tượng cảm ứng điện từ khi từ thông qua cuộn dây biến thiên.}
{Máy phát điện xoay chiều biến điện năng thành cơ năng nhờ hiệu ứng nhiệt.}
{Có thể bỏ qua phương, chiều và đơn vị trong mọi trường hợp.}
{Kết quả không phụ thuộc dữ kiện và điều kiện áp dụng.}
\loigiai{
Máy phát điện xoay chiều hoạt động dựa trên hiện tượng cảm ứng điện từ khi từ thông qua cuộn dây biến thiên. Vì vậy phương án $A$ đúng; các phương án còn lại sai.
}
\end{ex}

% ===== Câu 126 =====
% ID: L12C3B17-27-DS
% Mức: TH
\dangbt{Xác định biên độ, tần số, chu kì và pha của suất điện động xoay chiều}
\begin{ex}
% ID: L12C3B17-27-DS
% Mức: TH
Xét Xác định biên độ, tần số, chu kì và pha của suất điện động xoay chiều (Tình huống 2). Hãy xác định tính đúng, sai của bốn phát biểu sau.
\choiceTF
{Biên độ suất điện động không phụ thuộc tốc độ quay.}
{\True Biên độ suất điện động là E0 = NBS omega và tần số f = omega/(2pi).}
{Kết quả không phụ thuộc dữ kiện và có thể bỏ qua điều kiện.}
{\True Khi giải cần xác định đúng dữ kiện, phương chiều, điều kiện áp dụng và đơn vị.}
\loigiai{
\begin{itemchoice}
\item (Sai) Biên độ suất điện động không phụ thuộc tốc độ quay. (Vì): Biên độ suất điện động không phụ thuộc tốc độ quay. b) Đúng: Biên độ suất điện động là E0 = NBS omega và tần số f = omega/(2pi). c) Sai: Kết quả không phụ thuộc dữ kiện và có thể bỏ qua điều kiện. d) Đúng: Khi giải cần xác định đúng dữ kiện, phương chiều, điều kiện áp dụng và đơn vị. Kiến thức cốt lõi: Biên độ suất điện động là E0 = NBS omega và tần số f = omega/(2pi
\item (Đúng) Biên độ suất điện động là E0 = NBS omega và tần số f = omega/(2pi).
\item (Sai) Kết quả không phụ thuộc dữ kiện và có thể bỏ qua điều kiện.
\item (Đúng) Khi giải cần xác định đúng dữ kiện, phương chiều, điều kiện áp dụng và đơn vị.
\end{itemchoice}
}
\end{ex}

% ===== Câu 127 =====
% ID: L12C3B17-27-TL
% Mức: VD
\dangbt{Tính từ thông qua khung dây quay trong từ trường}
\begin{bt}
% ID: L12C3B17-27-TL
% Mức: VD
Phân tích sai lầm trong nhận định: “Từ thông qua khung quay luôn không đổi và bằng NBS.”
\loigiai{
Cần nêu được: Khung $N$ vòng quay đều có từ thông Phi = NBS cos(omega t + phi). Khi vận dụng phải xác định đúng dữ kiện, phương chiều nếu có, điều kiện áp dụng, hệ đơn vị và kiểm tra tính hợp lí. Nhận định “Từ thông qua khung quay luôn không đổi và bằng NBS.” sai.
}
\end{bt}

% ===== Câu 128 =====
% ID: L12C3B17-27-TLN
% Mức: TH
\begin{ex}
% ID: L12C3B17-27-TLN
% Mức: TH
Trong bốn phát biểu về Tính từ thông qua khung dây quay trong từ trường, có bao nhiêu phát biểu đúng? (Chỉ ghi một số nguyên.) (1) Khung $N$ vòng quay đều có từ thông Phi = NBS cos(omega t + phi). (2) Từ thông qua khung quay luôn không đổi và bằng NBS. (3) Cần kiểm tra điều kiện áp dụng. (4) Có thể bỏ qua đơn vị.
\shortans{2}
\loigiai{
Đối chiếu với kiến thức: Khung $N$ vòng quay đều có từ thông Phi = NBS cos(omega t + phi). Có 2 phát biểu đúng.
}
\end{ex}

% ===== Câu 129 =====
% ID: L12C3B17-27-TN
% Mức: TH
\dangbt{Nhận biết cấu tạo và nguyên lí của máy phát điện xoay chiều}
\begin{ex}
% ID: L12C3B17-27-TN
% Mức: TH
Chọn phát biểu không trái với kiến thức về Nhận biết cấu tạo và nguyên lí của máy phát điện xoay chiều?
\choice
{Kết quả không phụ thuộc dữ kiện và điều kiện áp dụng.}
{\True Máy phát điện xoay chiều hoạt động dựa trên hiện tượng cảm ứng điện từ khi từ thông qua cuộn dây biến thiên.}
{Máy phát điện xoay chiều biến điện năng thành cơ năng nhờ hiệu ứng nhiệt.}
{Có thể bỏ qua phương, chiều và đơn vị trong mọi trường hợp.}
\loigiai{
Máy phát điện xoay chiều hoạt động dựa trên hiện tượng cảm ứng điện từ khi từ thông qua cuộn dây biến thiên. Vì vậy phương án $B$ đúng; các phương án còn lại sai.
}
\end{ex}

% ===== Câu 130 =====
% ID: L12C3B17-28-DS
% Mức: VD
\dangbt{Xác định biên độ, tần số, chu kì và pha của suất điện động xoay chiều}
\begin{ex}
% ID: L12C3B17-28-DS
% Mức: VD
Xét Xác định biên độ, tần số, chu kì và pha của suất điện động xoay chiều (Tình huống 3). Hãy xác định tính đúng, sai của bốn phát biểu sau.
\choiceTF
{\True Khi giải cần xác định đúng dữ kiện, phương chiều, điều kiện áp dụng và đơn vị.}
{Kết quả không phụ thuộc dữ kiện và có thể bỏ qua điều kiện.}
{\True Biên độ suất điện động là E0 = NBS omega và tần số f = omega/(2pi).}
{Biên độ suất điện động không phụ thuộc tốc độ quay.}
\loigiai{
\begin{itemchoice}
\item (Đúng) Khi giải cần xác định đúng dữ kiện, phương chiều, điều kiện áp dụng và đơn vị. (Vì): Khi giải cần xác định đúng dữ kiện, phương chiều, điều kiện áp dụng và đơn vị. b) Sai: Kết quả không phụ thuộc dữ kiện và có thể bỏ qua điều kiện. c) Đúng: Biên độ suất điện động là E0 = NBS omega và tần số f = omega/(2pi). d) Sai: Biên độ suất điện động không phụ thuộc tốc độ quay. Kiến thức cốt lõi: Biên độ suất điện động là E0 = NBS omega và tần số f = omega/(2pi
\item (Sai) Kết quả không phụ thuộc dữ kiện và có thể bỏ qua điều kiện.
\item (Đúng) Biên độ suất điện động là E0 = NBS omega và tần số f = omega/(2pi).
\item (Sai) Biên độ suất điện động không phụ thuộc tốc độ quay.
\end{itemchoice}
}
\end{ex}

% ===== Câu 131 =====
% ID: L12C3B17-28-TL
% Mức: VD
\dangbt{Tính từ thông qua khung dây quay trong từ trường}
\begin{bt}
% ID: L12C3B17-28-TL
% Mức: VD
Đề xuất các bước giải một bài toán thuộc dạng Tính từ thông qua khung dây quay trong từ trường.
\loigiai{
Cần nêu được: Khung $N$ vòng quay đều có từ thông Phi = NBS cos(omega t + phi). Khi vận dụng phải xác định đúng dữ kiện, phương chiều nếu có, điều kiện áp dụng, hệ đơn vị và kiểm tra tính hợp lí. Nhận định “Từ thông qua khung quay luôn không đổi và bằng NBS.” sai.
}
\end{bt}

% ===== Câu 132 =====
% ID: L12C3B17-28-TLN
% Mức: VD
\begin{ex}
% ID: L12C3B17-28-TLN
% Mức: VD
Trong bốn phát biểu về Tính từ thông qua khung dây quay trong từ trường, có bao nhiêu phát biểu đúng? (Chỉ ghi một số nguyên.) (1) Từ thông qua khung quay luôn không đổi và bằng NBS. (2) Khung $N$ vòng quay đều có từ thông Phi = NBS cos(omega t + phi). (3) Kết quả phải phù hợp ý nghĩa vật lí. (4) Mọi đại lượng đều là vô hướng.
\shortans{2}
\loigiai{
Đối chiếu với kiến thức: Khung $N$ vòng quay đều có từ thông Phi = NBS cos(omega t + phi). Có 2 phát biểu đúng.
}
\end{ex}

% ===== Câu 133 =====
% ID: L12C3B17-28-TN
% Mức: VD
\dangbt{Nhận biết cấu tạo và nguyên lí của máy phát điện xoay chiều}
\begin{ex}
% ID: L12C3B17-28-TN
% Mức: VD
Cơ sở Vật lí đúng dùng để giải Nhận biết cấu tạo và nguyên lí của máy phát điện xoay chiều?
\choice
{Có thể bỏ qua phương, chiều và đơn vị trong mọi trường hợp.}
{Kết quả không phụ thuộc dữ kiện và điều kiện áp dụng.}
{\True Máy phát điện xoay chiều hoạt động dựa trên hiện tượng cảm ứng điện từ khi từ thông qua cuộn dây biến thiên.}
{Máy phát điện xoay chiều biến điện năng thành cơ năng nhờ hiệu ứng nhiệt.}
\loigiai{
Máy phát điện xoay chiều hoạt động dựa trên hiện tượng cảm ứng điện từ khi từ thông qua cuộn dây biến thiên. Vì vậy phương án $C$ đúng; các phương án còn lại sai.
}
\end{ex}

% ===== Câu 134 =====
% ID: L12C3B17-29-DS
% Mức: VD
\dangbt{Xác định biên độ, tần số, chu kì và pha của suất điện động xoay chiều}
\begin{ex}
% ID: L12C3B17-29-DS
% Mức: VD
Xét Xác định biên độ, tần số, chu kì và pha của suất điện động xoay chiều (Tình huống 4). Hãy xác định tính đúng, sai của bốn phát biểu sau.
\choiceTF
{Kết quả không phụ thuộc dữ kiện và có thể bỏ qua điều kiện.}
{\True Khi giải cần xác định đúng dữ kiện, phương chiều, điều kiện áp dụng và đơn vị.}
{Biên độ suất điện động không phụ thuộc tốc độ quay.}
{\True Biên độ suất điện động là E0 = NBS omega và tần số f = omega/(2pi).}
\loigiai{
\begin{itemchoice}
\item (Sai) Kết quả không phụ thuộc dữ kiện và có thể bỏ qua điều kiện. (Vì): Kết quả không phụ thuộc dữ kiện và có thể bỏ qua điều kiện. b) Đúng: Khi giải cần xác định đúng dữ kiện, phương chiều, điều kiện áp dụng và đơn vị. c) Sai: Biên độ suất điện động không phụ thuộc tốc độ quay. d) Đúng: Biên độ suất điện động là E0 = NBS omega và tần số f = omega/(2pi). Kiến thức cốt lõi: Biên độ suất điện động là E0 = NBS omega và tần số f = omega/(2pi
\item (Đúng) Khi giải cần xác định đúng dữ kiện, phương chiều, điều kiện áp dụng và đơn vị.
\item (Sai) Biên độ suất điện động không phụ thuộc tốc độ quay.
\item (Đúng) Biên độ suất điện động là E0 = NBS omega và tần số f = omega/(2pi).
\end{itemchoice}
}
\end{ex}

% ===== Câu 135 =====
% ID: L12C3B17-29-TL
% Mức: VD
\dangbt{Tính từ thông qua khung dây quay trong từ trường}
\begin{bt}
% ID: L12C3B17-29-TL
% Mức: VD
Nêu một tình huống thực tiễn liên quan đến Tính từ thông qua khung dây quay trong từ trường và giải thích bằng kiến thức Vật lí.
\loigiai{
Cần nêu được: Khung $N$ vòng quay đều có từ thông Phi = NBS cos(omega t + phi). Khi vận dụng phải xác định đúng dữ kiện, phương chiều nếu có, điều kiện áp dụng, hệ đơn vị và kiểm tra tính hợp lí. Nhận định “Từ thông qua khung quay luôn không đổi và bằng NBS.” sai.
}
\end{bt}

% ===== Câu 136 =====
% ID: L12C3B17-29-TLN
% Mức: VD
\begin{ex}
% ID: L12C3B17-29-TLN
% Mức: VD
Trong bốn phát biểu về Tính từ thông qua khung dây quay trong từ trường, có bao nhiêu phát biểu đúng? (Chỉ ghi một số nguyên.) (1) Khung $N$ vòng quay đều có từ thông Phi = NBS cos(omega t + phi). (2) Từ thông qua khung quay luôn không đổi và bằng NBS. (3) Cần kiểm tra điều kiện áp dụng. (4) Có thể bỏ qua đơn vị.
\shortans{2}
\loigiai{
Đối chiếu với kiến thức: Khung $N$ vòng quay đều có từ thông Phi = NBS cos(omega t + phi). Có 2 phát biểu đúng.
}
\end{ex}

% ===== Câu 137 =====
% ID: L12C3B17-29-TN
% Mức: VD
\dangbt{Nhận biết cấu tạo và nguyên lí của máy phát điện xoay chiều}
\begin{ex}
% ID: L12C3B17-29-TN
% Mức: VD
Trong một bài thực hành về Nhận biết cấu tạo và nguyên lí của máy phát điện xoay chiều?
\choice
{Máy phát điện xoay chiều biến điện năng thành cơ năng nhờ hiệu ứng nhiệt.}
{Có thể bỏ qua phương, chiều và đơn vị trong mọi trường hợp.}
{Kết quả không phụ thuộc dữ kiện và điều kiện áp dụng.}
{\True Máy phát điện xoay chiều hoạt động dựa trên hiện tượng cảm ứng điện từ khi từ thông qua cuộn dây biến thiên.}
\loigiai{
Máy phát điện xoay chiều hoạt động dựa trên hiện tượng cảm ứng điện từ khi từ thông qua cuộn dây biến thiên. Vì vậy phương án $D$ đúng; các phương án còn lại sai.
}
\end{ex}

% ===== Câu 138 =====
% ID: L12C3B17-30-DS
% Mức: VD
\dangbt{Xác định biên độ, tần số, chu kì và pha của suất điện động xoay chiều}
\begin{ex}
% ID: L12C3B17-30-DS
% Mức: VD
Xét Xác định biên độ, tần số, chu kì và pha của suất điện động xoay chiều (Tình huống 5). Hãy xác định tính đúng, sai của bốn phát biểu sau.
\choiceTF
{\True Biên độ suất điện động là E0 = NBS omega và tần số f = omega/(2pi).}
{Kết quả không phụ thuộc dữ kiện và có thể bỏ qua điều kiện.}
{Biên độ suất điện động không phụ thuộc tốc độ quay.}
{\True Khi giải cần xác định đúng dữ kiện, phương chiều, điều kiện áp dụng và đơn vị.}
\loigiai{
\begin{itemchoice}
\item (Đúng) Biên độ suất điện động là E0 = NBS omega và tần số f = omega/(2pi). (Vì): Biên độ suất điện động là E0 = NBS omega và tần số f = omega/(2pi). b) Sai: Kết quả không phụ thuộc dữ kiện và có thể bỏ qua điều kiện. c) Sai: Biên độ suất điện động không phụ thuộc tốc độ quay. d) Đúng: Khi giải cần xác định đúng dữ kiện, phương chiều, điều kiện áp dụng và đơn vị. Kiến thức cốt lõi: Biên độ suất điện động là E0 = NBS omega và tần số f = omega/(2pi
\item (Sai) Kết quả không phụ thuộc dữ kiện và có thể bỏ qua điều kiện.
\item (Sai) Biên độ suất điện động không phụ thuộc tốc độ quay.
\item (Đúng) Khi giải cần xác định đúng dữ kiện, phương chiều, điều kiện áp dụng và đơn vị.
\end{itemchoice}
}
\end{ex}

% ===== Câu 139 =====
% ID: L12C3B17-30-TL
% Mức: VDC
\dangbt{Tính từ thông qua khung dây quay trong từ trường}
\begin{bt}
% ID: L12C3B17-30-TL
% Mức: VDC
So sánh nhận định đúng “Khung $N$ vòng quay đều có từ thông Phi = NBS cos(omega t + phi).” với nhận định sai “Từ thông qua khung quay luôn không đổi và bằng NBS.”; chỉ rõ căn cứ Vật lí.
\loigiai{
Cần nêu được: Khung $N$ vòng quay đều có từ thông Phi = NBS cos(omega t + phi). Khi vận dụng phải xác định đúng dữ kiện, phương chiều nếu có, điều kiện áp dụng, hệ đơn vị và kiểm tra tính hợp lí. Nhận định “Từ thông qua khung quay luôn không đổi và bằng NBS.” sai.
}
\end{bt}

% ===== Câu 140 =====
% ID: L12C3B17-30-TLN
% Mức: VDC
\begin{ex}
% ID: L12C3B17-30-TLN
% Mức: VDC
Trong bốn phát biểu về Tính từ thông qua khung dây quay trong từ trường, có bao nhiêu phát biểu đúng? (Chỉ ghi một số nguyên.) (1) Khi giải cần xác định đúng dữ kiện, phương chiều, điều kiện áp dụng và đơn vị. (2) Từ thông qua khung quay luôn không đổi và bằng NBS. (3) Phải dùng đúng định luật cho tình huống. (4) Khung $N$ vòng quay đều có từ thông Phi = NBS cos(omega t + phi).
\shortans{3}
\loigiai{
Đối chiếu với kiến thức: Khung $N$ vòng quay đều có từ thông Phi = NBS cos(omega t + phi). Có 3 phát biểu đúng.
}
\end{ex}

% ===== Câu 141 =====
% ID: L12C3B17-30-TN
% Mức: VD
\dangbt{Nhận biết cấu tạo và nguyên lí của máy phát điện xoay chiều}
\begin{ex}
% ID: L12C3B17-30-TN
% Mức: VD
Kết luận đúng khi vận dụng Nhận biết cấu tạo và nguyên lí của máy phát điện xoay chiều?
\choice
{\True Máy phát điện xoay chiều hoạt động dựa trên hiện tượng cảm ứng điện từ khi từ thông qua cuộn dây biến thiên.}
{Máy phát điện xoay chiều biến điện năng thành cơ năng nhờ hiệu ứng nhiệt.}
{Có thể bỏ qua phương, chiều và đơn vị trong mọi trường hợp.}
{Kết quả không phụ thuộc dữ kiện và điều kiện áp dụng.}
\loigiai{
Máy phát điện xoay chiều hoạt động dựa trên hiện tượng cảm ứng điện từ khi từ thông qua cuộn dây biến thiên. Vì vậy phương án $A$ đúng; các phương án còn lại sai.
}
\end{ex}

% ===== Câu 142 =====
% ID: L12C3B17-31-DS
% Mức: TH
\begin{ex}
% ID: L12C3B17-31-DS
% Mức: TH
Xét Nhận biết cấu tạo và nguyên lí của máy phát điện xoay chiều (Tình huống 1). Hãy xác định tính đúng, sai của bốn phát biểu sau.
\choiceTF
{\True Máy phát điện xoay chiều hoạt động dựa trên hiện tượng cảm ứng điện từ khi từ thông qua cuộn dây biến thiên.}
{Máy phát điện xoay chiều biến điện năng thành cơ năng nhờ hiệu ứng nhiệt.}
{\True Khi giải cần xác định đúng dữ kiện, phương chiều, điều kiện áp dụng và đơn vị.}
{Kết quả không phụ thuộc dữ kiện và có thể bỏ qua điều kiện.}
\loigiai{
\begin{itemchoice}
\item (Đúng) Máy phát điện xoay chiều hoạt động dựa trên hiện tượng cảm ứng điện từ khi từ thông qua cuộn dây biến thiên. (Vì): Máy phát điện xoay chiều hoạt động dựa trên hiện tượng cảm ứng điện từ khi từ thông qua cuộn dây biến thiên. b) Sai: Máy phát điện xoay chiều biến điện năng thành cơ năng nhờ hiệu ứng nhiệt. c) Đúng: Khi giải cần xác định đúng dữ kiện, phương chiều, điều kiện áp dụng và đơn vị. d) Sai: Kết quả không phụ thuộc dữ kiện và có thể bỏ qua điều kiện. Kiến thức cốt lõi: Máy phát điện xoay chiều hoạt động dựa trên hiện tượng cảm ứng điện từ khi từ thông qua cuộn dây biến thiên
\item (Sai) Máy phát điện xoay chiều biến điện năng thành cơ năng nhờ hiệu ứng nhiệt.
\item (Đúng) Khi giải cần xác định đúng dữ kiện, phương chiều, điều kiện áp dụng và đơn vị.
\item (Sai) Kết quả không phụ thuộc dữ kiện và có thể bỏ qua điều kiện.
\end{itemchoice}
}
\end{ex}

% ===== Câu 143 =====
% ID: L12C3B17-31-TL
% Mức: TH
\dangbt{Xác định biên độ, tần số, chu kì và pha của suất điện động xoay chiều}
\begin{bt}
% ID: L12C3B17-31-TL
% Mức: TH
Trình bày kiến thức cốt lõi của Xác định biên độ, tần số, chu kì và pha của suất điện động xoay chiều và nêu điều kiện áp dụng.
\loigiai{
Cần nêu được: Biên độ suất điện động là E0 = NBS omega và tần số f = omega/(2pi). Khi vận dụng phải xác định đúng dữ kiện, phương chiều nếu có, điều kiện áp dụng, hệ đơn vị và kiểm tra tính hợp lí. Nhận định “Biên độ suất điện động không phụ thuộc tốc độ quay.” sai.
}
\end{bt}

% ===== Câu 144 =====
% ID: L12C3B17-31-TLN
% Mức: TH
\begin{ex}
% ID: L12C3B17-31-TLN
% Mức: TH
Trong bốn phát biểu về Xác định biên độ, tần số, chu kì và pha của suất điện động xoay chiều, có bao nhiêu phát biểu đúng? (Chỉ ghi một số nguyên.) (1) Biên độ suất điện động là E0 = NBS omega và tần số f = omega/(2pi). (2) Biên độ suất điện động không phụ thuộc tốc độ quay. (3) Cần kiểm tra điều kiện áp dụng. (4) Có thể bỏ qua đơn vị.
\shortans{2}
\loigiai{
Đối chiếu với kiến thức: Biên độ suất điện động là E0 = NBS omega và tần số f = omega/(2pi). Có 2 phát biểu đúng.
}
\end{ex}

% ===== Câu 145 =====
% ID: L12C3B17-31-TN
% Mức: NB
\dangbt{Thiết lập biểu thức suất điện động xoay chiều do máy phát tạo ra}
\begin{ex}
% ID: L12C3B17-31-TN
% Mức: NB
Phát biểu nào sau đây đúng về Thiết lập biểu thức suất điện động xoay chiều do máy phát tạo ra?
\choice
{Suất điện động luôn cùng pha và bằng đúng từ thông.}
{Có thể bỏ qua phương, chiều và đơn vị trong mọi trường hợp.}
{Kết quả không phụ thuộc dữ kiện và điều kiện áp dụng.}
{\True Nếu Phi = Phi0 cos(omega t + phi) thì e = omega Phi0 sin(omega t + phi).}
\loigiai{
Nếu Phi = Phi0 cos(omega t + phi) thì e = omega Phi0 sin(omega t + phi). Vì vậy phương án $D$ đúng; các phương án còn lại sai.
}
\end{ex}

% ===== Câu 146 =====
% ID: L12C3B17-32-DS
% Mức: TH
\dangbt{Nhận biết cấu tạo và nguyên lí của máy phát điện xoay chiều}
\begin{ex}
% ID: L12C3B17-32-DS
% Mức: TH
Xét Nhận biết cấu tạo và nguyên lí của máy phát điện xoay chiều (Tình huống 2). Hãy xác định tính đúng, sai của bốn phát biểu sau.
\choiceTF
{Máy phát điện xoay chiều biến điện năng thành cơ năng nhờ hiệu ứng nhiệt.}
{\True Máy phát điện xoay chiều hoạt động dựa trên hiện tượng cảm ứng điện từ khi từ thông qua cuộn dây biến thiên.}
{Kết quả không phụ thuộc dữ kiện và có thể bỏ qua điều kiện.}
{\True Khi giải cần xác định đúng dữ kiện, phương chiều, điều kiện áp dụng và đơn vị.}
\loigiai{
\begin{itemchoice}
\item (Sai) Máy phát điện xoay chiều biến điện năng thành cơ năng nhờ hiệu ứng nhiệt. (Vì): Máy phát điện xoay chiều biến điện năng thành cơ năng nhờ hiệu ứng nhiệt. b) Đúng: Máy phát điện xoay chiều hoạt động dựa trên hiện tượng cảm ứng điện từ khi từ thông qua cuộn dây biến thiên. c) Sai: Kết quả không phụ thuộc dữ kiện và có thể bỏ qua điều kiện. d) Đúng: Khi giải cần xác định đúng dữ kiện, phương chiều, điều kiện áp dụng và đơn vị. Kiến thức cốt lõi: Máy phát điện xoay chiều hoạt động dựa trên hiện tượng cảm ứng điện từ khi từ thông qua cuộn dây biến thiên
\item (Đúng) Máy phát điện xoay chiều hoạt động dựa trên hiện tượng cảm ứng điện từ khi từ thông qua cuộn dây biến thiên.
\item (Sai) Kết quả không phụ thuộc dữ kiện và có thể bỏ qua điều kiện.
\item (Đúng) Khi giải cần xác định đúng dữ kiện, phương chiều, điều kiện áp dụng và đơn vị.
\end{itemchoice}
}
\end{ex}

% ===== Câu 147 =====
% ID: L12C3B17-32-TL
% Mức: TH
\dangbt{Xác định biên độ, tần số, chu kì và pha của suất điện động xoay chiều}
\begin{bt}
% ID: L12C3B17-32-TL
% Mức: TH
Giải thích vì sao nhận định sau đúng: “Biên độ suất điện động là E0 = NBS omega và tần số f = omega/(2pi).”
\loigiai{
Cần nêu được: Biên độ suất điện động là E0 = NBS omega và tần số f = omega/(2pi). Khi vận dụng phải xác định đúng dữ kiện, phương chiều nếu có, điều kiện áp dụng, hệ đơn vị và kiểm tra tính hợp lí. Nhận định “Biên độ suất điện động không phụ thuộc tốc độ quay.” sai.
}
\end{bt}

% ===== Câu 148 =====
% ID: L12C3B17-32-TLN
% Mức: TH
\begin{ex}
% ID: L12C3B17-32-TLN
% Mức: TH
Trong bốn phát biểu về Xác định biên độ, tần số, chu kì và pha của suất điện động xoay chiều, có bao nhiêu phát biểu đúng? (Chỉ ghi một số nguyên.) (1) Biên độ suất điện động không phụ thuộc tốc độ quay. (2) Biên độ suất điện động là E0 = NBS omega và tần số f = omega/(2pi). (3) Kết quả phải phù hợp ý nghĩa vật lí. (4) Mọi đại lượng đều là vô hướng.
\shortans{1}
\loigiai{
Đối chiếu với kiến thức: Biên độ suất điện động là E0 = NBS omega và tần số f = omega/(2pi). Có 1 phát biểu đúng.
}
\end{ex}

% ===== Câu 149 =====
% ID: L12C3B17-32-TN
% Mức: NB
\dangbt{Thiết lập biểu thức suất điện động xoay chiều do máy phát tạo ra}
\begin{ex}
% ID: L12C3B17-32-TN
% Mức: NB
Trong nội dung Thiết lập biểu thức suất điện động xoay chiều do máy phát tạo ra?
\choice
{\True Nếu Phi = Phi0 cos(omega t + phi) thì e = omega Phi0 sin(omega t + phi).}
{Suất điện động luôn cùng pha và bằng đúng từ thông.}
{Có thể bỏ qua phương, chiều và đơn vị trong mọi trường hợp.}
{Kết quả không phụ thuộc dữ kiện và điều kiện áp dụng.}
\loigiai{
Nếu Phi = Phi0 cos(omega t + phi) thì e = omega Phi0 sin(omega t + phi). Vì vậy phương án $A$ đúng; các phương án còn lại sai.
}
\end{ex}

% ===== Câu 150 =====
% ID: L12C3B17-33-DS
% Mức: VD
\dangbt{Nhận biết cấu tạo và nguyên lí của máy phát điện xoay chiều}
\begin{ex}
% ID: L12C3B17-33-DS
% Mức: VD
Xét Nhận biết cấu tạo và nguyên lí của máy phát điện xoay chiều (Tình huống 3). Hãy xác định tính đúng, sai của bốn phát biểu sau.
\choiceTF
{\True Khi giải cần xác định đúng dữ kiện, phương chiều, điều kiện áp dụng và đơn vị.}
{Kết quả không phụ thuộc dữ kiện và có thể bỏ qua điều kiện.}
{\True Máy phát điện xoay chiều hoạt động dựa trên hiện tượng cảm ứng điện từ khi từ thông qua cuộn dây biến thiên.}
{Máy phát điện xoay chiều biến điện năng thành cơ năng nhờ hiệu ứng nhiệt.}
\loigiai{
\begin{itemchoice}
\item (Đúng) Khi giải cần xác định đúng dữ kiện, phương chiều, điều kiện áp dụng và đơn vị. (Vì): Khi giải cần xác định đúng dữ kiện, phương chiều, điều kiện áp dụng và đơn vị. b) Sai: Kết quả không phụ thuộc dữ kiện và có thể bỏ qua điều kiện. c) Đúng: Máy phát điện xoay chiều hoạt động dựa trên hiện tượng cảm ứng điện từ khi từ thông qua cuộn dây biến thiên. d) Sai: Máy phát điện xoay chiều biến điện năng thành cơ năng nhờ hiệu ứng nhiệt. Kiến thức cốt lõi: Máy phát điện xoay chiều hoạt động dựa trên hiện tượng cảm ứng điện từ khi từ thông qua cuộn dây biến thiên
\item (Sai) Kết quả không phụ thuộc dữ kiện và có thể bỏ qua điều kiện.
\item (Đúng) Máy phát điện xoay chiều hoạt động dựa trên hiện tượng cảm ứng điện từ khi từ thông qua cuộn dây biến thiên.
\item (Sai) Máy phát điện xoay chiều biến điện năng thành cơ năng nhờ hiệu ứng nhiệt.
\end{itemchoice}
}
\end{ex}

% ===== Câu 151 =====
% ID: L12C3B17-33-TL
% Mức: VD
\dangbt{Xác định biên độ, tần số, chu kì và pha của suất điện động xoay chiều}
\begin{bt}
% ID: L12C3B17-33-TL
% Mức: VD
Phân tích sai lầm trong nhận định: “Biên độ suất điện động không phụ thuộc tốc độ quay.”
\loigiai{
Cần nêu được: Biên độ suất điện động là E0 = NBS omega và tần số f = omega/(2pi). Khi vận dụng phải xác định đúng dữ kiện, phương chiều nếu có, điều kiện áp dụng, hệ đơn vị và kiểm tra tính hợp lí. Nhận định “Biên độ suất điện động không phụ thuộc tốc độ quay.” sai.
}
\end{bt}

% ===== Câu 152 =====
% ID: L12C3B17-33-TLN
% Mức: TH
\begin{ex}
% ID: L12C3B17-33-TLN
% Mức: TH
Trong bốn phát biểu về Xác định biên độ, tần số, chu kì và pha của suất điện động xoay chiều, có bao nhiêu phát biểu đúng? (Chỉ ghi một số nguyên.) (1) Biên độ suất điện động là E0 = NBS omega và tần số f = omega/(2pi). (2) Biên độ suất điện động không phụ thuộc tốc độ quay. (3) Cần kiểm tra điều kiện áp dụng. (4) Có thể bỏ qua đơn vị.
\shortans{2}
\loigiai{
Đối chiếu với kiến thức: Biên độ suất điện động là E0 = NBS omega và tần số f = omega/(2pi). Có 2 phát biểu đúng.
}
\end{ex}

% ===== Câu 153 =====
% ID: L12C3B17-33-TN
% Mức: NB
\dangbt{Thiết lập biểu thức suất điện động xoay chiều do máy phát tạo ra}
\begin{ex}
% ID: L12C3B17-33-TN
% Mức: NB
Tình huống 3: chọn kết luận chính xác về Thiết lập biểu thức suất điện động xoay chiều do máy phát tạo ra?
\choice
{Kết quả không phụ thuộc dữ kiện và điều kiện áp dụng.}
{\True Nếu Phi = Phi0 cos(omega t + phi) thì e = omega Phi0 sin(omega t + phi).}
{Suất điện động luôn cùng pha và bằng đúng từ thông.}
{Có thể bỏ qua phương, chiều và đơn vị trong mọi trường hợp.}
\loigiai{
Nếu Phi = Phi0 cos(omega t + phi) thì e = omega Phi0 sin(omega t + phi). Vì vậy phương án $B$ đúng; các phương án còn lại sai.
}
\end{ex}

% ===== Câu 154 =====
% ID: L12C3B17-34-DS
% Mức: VD
\dangbt{Nhận biết cấu tạo và nguyên lí của máy phát điện xoay chiều}
\begin{ex}
% ID: L12C3B17-34-DS
% Mức: VD
Xét Nhận biết cấu tạo và nguyên lí của máy phát điện xoay chiều (Tình huống 4). Hãy xác định tính đúng, sai của bốn phát biểu sau.
\choiceTF
{Kết quả không phụ thuộc dữ kiện và có thể bỏ qua điều kiện.}
{\True Khi giải cần xác định đúng dữ kiện, phương chiều, điều kiện áp dụng và đơn vị.}
{Máy phát điện xoay chiều biến điện năng thành cơ năng nhờ hiệu ứng nhiệt.}
{\True Máy phát điện xoay chiều hoạt động dựa trên hiện tượng cảm ứng điện từ khi từ thông qua cuộn dây biến thiên.}
\loigiai{
\begin{itemchoice}
\item (Sai) Kết quả không phụ thuộc dữ kiện và có thể bỏ qua điều kiện. (Vì): Kết quả không phụ thuộc dữ kiện và có thể bỏ qua điều kiện. b) Đúng: Khi giải cần xác định đúng dữ kiện, phương chiều, điều kiện áp dụng và đơn vị. c) Sai: Máy phát điện xoay chiều biến điện năng thành cơ năng nhờ hiệu ứng nhiệt. d) Đúng: Máy phát điện xoay chiều hoạt động dựa trên hiện tượng cảm ứng điện từ khi từ thông qua cuộn dây biến thiên. Kiến thức cốt lõi: Máy phát điện xoay chiều hoạt động dựa trên hiện tượng cảm ứng điện từ khi từ thông qua cuộn dây biến thiên
\item (Đúng) Khi giải cần xác định đúng dữ kiện, phương chiều, điều kiện áp dụng và đơn vị.
\item (Sai) Máy phát điện xoay chiều biến điện năng thành cơ năng nhờ hiệu ứng nhiệt.
\item (Đúng) Máy phát điện xoay chiều hoạt động dựa trên hiện tượng cảm ứng điện từ khi từ thông qua cuộn dây biến thiên.
\end{itemchoice}
}
\end{ex}

% ===== Câu 155 =====
% ID: L12C3B17-34-TL
% Mức: VD
\dangbt{Xác định biên độ, tần số, chu kì và pha của suất điện động xoay chiều}
\begin{bt}
% ID: L12C3B17-34-TL
% Mức: VD
Đề xuất các bước giải một bài toán thuộc dạng Xác định biên độ, tần số, chu kì và pha của suất điện động xoay chiều.
\loigiai{
Cần nêu được: Biên độ suất điện động là E0 = NBS omega và tần số f = omega/(2pi). Khi vận dụng phải xác định đúng dữ kiện, phương chiều nếu có, điều kiện áp dụng, hệ đơn vị và kiểm tra tính hợp lí. Nhận định “Biên độ suất điện động không phụ thuộc tốc độ quay.” sai.
}
\end{bt}

% ===== Câu 156 =====
% ID: L12C3B17-34-TLN
% Mức: VD
\begin{ex}
% ID: L12C3B17-34-TLN
% Mức: VD
Trong bốn phát biểu về Xác định biên độ, tần số, chu kì và pha của suất điện động xoay chiều, có bao nhiêu phát biểu đúng? (Chỉ ghi một số nguyên.) (1) Biên độ suất điện động không phụ thuộc tốc độ quay. (2) Biên độ suất điện động là E0 = NBS omega và tần số f = omega/(2pi). (3) Kết quả phải phù hợp ý nghĩa vật lí. (4) Mọi đại lượng đều là vô hướng.
\shortans{2}
\loigiai{
Đối chiếu với kiến thức: Biên độ suất điện động là E0 = NBS omega và tần số f = omega/(2pi). Có 2 phát biểu đúng.
}
\end{ex}

% ===== Câu 157 =====
% ID: L12C3B17-34-TN
% Mức: TH
\dangbt{Thiết lập biểu thức suất điện động xoay chiều do máy phát tạo ra}
\begin{ex}
% ID: L12C3B17-34-TN
% Mức: TH
Nội dung nào mô tả đúng nhất Thiết lập biểu thức suất điện động xoay chiều do máy phát tạo ra?
\choice
{Có thể bỏ qua phương, chiều và đơn vị trong mọi trường hợp.}
{Kết quả không phụ thuộc dữ kiện và điều kiện áp dụng.}
{\True Nếu Phi = Phi0 cos(omega t + phi) thì e = omega Phi0 sin(omega t + phi).}
{Suất điện động luôn cùng pha và bằng đúng từ thông.}
\loigiai{
Nếu Phi = Phi0 cos(omega t + phi) thì e = omega Phi0 sin(omega t + phi). Vì vậy phương án $C$ đúng; các phương án còn lại sai.
}
\end{ex}

% ===== Câu 158 =====
% ID: L12C3B17-35-DS
% Mức: VD
\dangbt{Nhận biết cấu tạo và nguyên lí của máy phát điện xoay chiều}
\begin{ex}
% ID: L12C3B17-35-DS
% Mức: VD
Xét Nhận biết cấu tạo và nguyên lí của máy phát điện xoay chiều (Tình huống 5). Hãy xác định tính đúng, sai của bốn phát biểu sau.
\choiceTF
{\True Máy phát điện xoay chiều hoạt động dựa trên hiện tượng cảm ứng điện từ khi từ thông qua cuộn dây biến thiên.}
{Kết quả không phụ thuộc dữ kiện và có thể bỏ qua điều kiện.}
{Máy phát điện xoay chiều biến điện năng thành cơ năng nhờ hiệu ứng nhiệt.}
{\True Khi giải cần xác định đúng dữ kiện, phương chiều, điều kiện áp dụng và đơn vị.}
\loigiai{
\begin{itemchoice}
\item (Đúng) Máy phát điện xoay chiều hoạt động dựa trên hiện tượng cảm ứng điện từ khi từ thông qua cuộn dây biến thiên. (Vì): Máy phát điện xoay chiều hoạt động dựa trên hiện tượng cảm ứng điện từ khi từ thông qua cuộn dây biến thiên. b) Sai: Kết quả không phụ thuộc dữ kiện và có thể bỏ qua điều kiện. c) Sai: Máy phát điện xoay chiều biến điện năng thành cơ năng nhờ hiệu ứng nhiệt. d) Đúng: Khi giải cần xác định đúng dữ kiện, phương chiều, điều kiện áp dụng và đơn vị. Kiến thức cốt lõi: Máy phát điện xoay chiều hoạt động dựa trên hiện tượng cảm ứng điện từ khi từ thông qua cuộn dây biến thiên
\item (Sai) Kết quả không phụ thuộc dữ kiện và có thể bỏ qua điều kiện.
\item (Sai) Máy phát điện xoay chiều biến điện năng thành cơ năng nhờ hiệu ứng nhiệt.
\item (Đúng) Khi giải cần xác định đúng dữ kiện, phương chiều, điều kiện áp dụng và đơn vị.
\end{itemchoice}
}
\end{ex}

% ===== Câu 159 =====
% ID: L12C3B17-35-TL
% Mức: VD
\dangbt{Xác định biên độ, tần số, chu kì và pha của suất điện động xoay chiều}
\begin{bt}
% ID: L12C3B17-35-TL
% Mức: VD
Nêu một tình huống thực tiễn liên quan đến Xác định biên độ, tần số, chu kì và pha của suất điện động xoay chiều và giải thích bằng kiến thức Vật lí.
\loigiai{
Cần nêu được: Biên độ suất điện động là E0 = NBS omega và tần số f = omega/(2pi). Khi vận dụng phải xác định đúng dữ kiện, phương chiều nếu có, điều kiện áp dụng, hệ đơn vị và kiểm tra tính hợp lí. Nhận định “Biên độ suất điện động không phụ thuộc tốc độ quay.” sai.
}
\end{bt}

% ===== Câu 160 =====
% ID: L12C3B17-35-TLN
% Mức: VD
\begin{ex}
% ID: L12C3B17-35-TLN
% Mức: VD
Trong bốn phát biểu về Xác định biên độ, tần số, chu kì và pha của suất điện động xoay chiều, có bao nhiêu phát biểu đúng? (Chỉ ghi một số nguyên.) (1) Biên độ suất điện động là E0 = NBS omega và tần số f = omega/(2pi). (2) Biên độ suất điện động không phụ thuộc tốc độ quay. (3) Cần kiểm tra điều kiện áp dụng. (4) Có thể bỏ qua đơn vị.
\shortans{2}
\loigiai{
Đối chiếu với kiến thức: Biên độ suất điện động là E0 = NBS omega và tần số f = omega/(2pi). Có 2 phát biểu đúng.
}
\end{ex}

% ===== Câu 161 =====
% ID: L12C3B17-35-TN
% Mức: TH
\dangbt{Thiết lập biểu thức suất điện động xoay chiều do máy phát tạo ra}
\begin{ex}
% ID: L12C3B17-35-TN
% Mức: TH
Khi phân tích Thiết lập biểu thức suất điện động xoay chiều do máy phát tạo ra?
\choice
{Suất điện động luôn cùng pha và bằng đúng từ thông.}
{Có thể bỏ qua phương, chiều và đơn vị trong mọi trường hợp.}
{Kết quả không phụ thuộc dữ kiện và điều kiện áp dụng.}
{\True Nếu Phi = Phi0 cos(omega t + phi) thì e = omega Phi0 sin(omega t + phi).}
\loigiai{
Nếu Phi = Phi0 cos(omega t + phi) thì e = omega Phi0 sin(omega t + phi). Vì vậy phương án $D$ đúng; các phương án còn lại sai.
}
\end{ex}

% ===== Câu 162 =====
% ID: L12C3B17-36-DS
% Mức: TH
\dangbt{Tính từ thông qua khung dây quay trong từ trường}
\begin{ex}
% ID: L12C3B17-36-DS
% Mức: TH
Xét Tính từ thông qua khung dây quay trong từ trường (Tình huống 1). Hãy xác định tính đúng, sai của bốn phát biểu sau.
\choiceTF
{\True Khung $N$ vòng quay đều có từ thông Phi = NBS cos(omega t + phi).}
{Từ thông qua khung quay luôn không đổi và bằng NBS.}
{\True Khi giải cần xác định đúng dữ kiện, phương chiều, điều kiện áp dụng và đơn vị.}
{Kết quả không phụ thuộc dữ kiện và có thể bỏ qua điều kiện.}
\loigiai{
\begin{itemchoice}
\item (Đúng) Khung $N$ vòng quay đều có từ thông Phi = NBS cos(omega t + phi). (Vì): Khung $N$ vòng quay đều có từ thông Phi = NBS cos(omega t + phi). b) Sai: Từ thông qua khung quay luôn không đổi và bằng NBS. c) Đúng: Khi giải cần xác định đúng dữ kiện, phương chiều, điều kiện áp dụng và đơn vị. d) Sai: Kết quả không phụ thuộc dữ kiện và có thể bỏ qua điều kiện. Kiến thức cốt lõi: Khung $N$ vòng quay đều có từ thông Phi = NBS cos(omega t + phi
\item (Sai) Từ thông qua khung quay luôn không đổi và bằng NBS.
\item (Đúng) Khi giải cần xác định đúng dữ kiện, phương chiều, điều kiện áp dụng và đơn vị.
\item (Sai) Kết quả không phụ thuộc dữ kiện và có thể bỏ qua điều kiện.
\end{itemchoice}
}
\end{ex}

% ===== Câu 163 =====
% ID: L12C3B17-36-TL
% Mức: VDC
\dangbt{Xác định biên độ, tần số, chu kì và pha của suất điện động xoay chiều}
\begin{bt}
% ID: L12C3B17-36-TL
% Mức: VDC
So sánh nhận định đúng “Biên độ suất điện động là E0 = NBS omega và tần số f = omega/(2pi).” với nhận định sai “Biên độ suất điện động không phụ thuộc tốc độ quay.”; chỉ rõ căn cứ Vật lí.
\loigiai{
Cần nêu được: Biên độ suất điện động là E0 = NBS omega và tần số f = omega/(2pi). Khi vận dụng phải xác định đúng dữ kiện, phương chiều nếu có, điều kiện áp dụng, hệ đơn vị và kiểm tra tính hợp lí. Nhận định “Biên độ suất điện động không phụ thuộc tốc độ quay.” sai.
}
\end{bt}

% ===== Câu 164 =====
% ID: L12C3B17-36-TLN
% Mức: VDC
\begin{ex}
% ID: L12C3B17-36-TLN
% Mức: VDC
Trong bốn phát biểu về Xác định biên độ, tần số, chu kì và pha của suất điện động xoay chiều, có bao nhiêu phát biểu đúng? (Chỉ ghi một số nguyên.) (1) Khi giải cần xác định đúng dữ kiện, phương chiều, điều kiện áp dụng và đơn vị. (2) Biên độ suất điện động không phụ thuộc tốc độ quay. (3) Phải dùng đúng định luật cho tình huống. (4) Biên độ suất điện động là E0 = NBS omega và tần số f = omega/(2pi).
\shortans{3}
\loigiai{
Đối chiếu với kiến thức: Biên độ suất điện động là E0 = NBS omega và tần số f = omega/(2pi). Có 3 phát biểu đúng.
}
\end{ex}

% ===== Câu 165 =====
% ID: L12C3B17-36-TN
% Mức: TH
\dangbt{Thiết lập biểu thức suất điện động xoay chiều do máy phát tạo ra}
\begin{ex}
% ID: L12C3B17-36-TN
% Mức: TH
Một học sinh ôn tập Thiết lập biểu thức suất điện động xoay chiều do máy phát tạo ra?
\choice
{\True Nếu Phi = Phi0 cos(omega t + phi) thì e = omega Phi0 sin(omega t + phi).}
{Suất điện động luôn cùng pha và bằng đúng từ thông.}
{Có thể bỏ qua phương, chiều và đơn vị trong mọi trường hợp.}
{Kết quả không phụ thuộc dữ kiện và điều kiện áp dụng.}
\loigiai{
Nếu Phi = Phi0 cos(omega t + phi) thì e = omega Phi0 sin(omega t + phi). Vì vậy phương án $A$ đúng; các phương án còn lại sai.
}
\end{ex}

% ===== Câu 166 =====
% ID: L12C3B17-37-DS
% Mức: TH
\dangbt{Tính từ thông qua khung dây quay trong từ trường}
\begin{ex}
% ID: L12C3B17-37-DS
% Mức: TH
Xét Tính từ thông qua khung dây quay trong từ trường (Tình huống 2). Hãy xác định tính đúng, sai của bốn phát biểu sau.
\choiceTF
{Từ thông qua khung quay luôn không đổi và bằng NBS.}
{\True Khung $N$ vòng quay đều có từ thông Phi = NBS cos(omega t + phi).}
{Kết quả không phụ thuộc dữ kiện và có thể bỏ qua điều kiện.}
{\True Khi giải cần xác định đúng dữ kiện, phương chiều, điều kiện áp dụng và đơn vị.}
\loigiai{
\begin{itemchoice}
\item (Sai) Từ thông qua khung quay luôn không đổi và bằng NBS. (Vì): Từ thông qua khung quay luôn không đổi và bằng NBS. b) Đúng: Khung $N$ vòng quay đều có từ thông Phi = NBS cos(omega t + phi). c) Sai: Kết quả không phụ thuộc dữ kiện và có thể bỏ qua điều kiện. d) Đúng: Khi giải cần xác định đúng dữ kiện, phương chiều, điều kiện áp dụng và đơn vị. Kiến thức cốt lõi: Khung $N$ vòng quay đều có từ thông Phi = NBS cos(omega t + phi
\item (Đúng) Khung $N$ vòng quay đều có từ thông Phi = NBS cos(omega t + phi).
\item (Sai) Kết quả không phụ thuộc dữ kiện và có thể bỏ qua điều kiện.
\item (Đúng) Khi giải cần xác định đúng dữ kiện, phương chiều, điều kiện áp dụng và đơn vị.
\end{itemchoice}
}
\end{ex}

% ===== Câu 167 =====
% ID: L12C3B17-37-TL
% Mức: TH
\dangbt{Nhận biết cấu tạo và nguyên lí của máy phát điện xoay chiều}
\begin{bt}
% ID: L12C3B17-37-TL
% Mức: TH
Trình bày kiến thức cốt lõi của Nhận biết cấu tạo và nguyên lí của máy phát điện xoay chiều và nêu điều kiện áp dụng.
\loigiai{
Cần nêu được: Máy phát điện xoay chiều hoạt động dựa trên hiện tượng cảm ứng điện từ khi từ thông qua cuộn dây biến thiên. Khi vận dụng phải xác định đúng dữ kiện, phương chiều nếu có, điều kiện áp dụng, hệ đơn vị và kiểm tra tính hợp lí. Nhận định “Máy phát điện xoay chiều biến điện năng thành cơ năng nhờ hiệu ứng nhiệt.” sai.
}
\end{bt}

% ===== Câu 168 =====
% ID: L12C3B17-37-TLN
% Mức: TH
\begin{ex}
% ID: L12C3B17-37-TLN
% Mức: TH
Trong bốn phát biểu về Nhận biết cấu tạo và nguyên lí của máy phát điện xoay chiều, có bao nhiêu phát biểu đúng? (Chỉ ghi một số nguyên.) (1) Máy phát điện xoay chiều hoạt động dựa trên hiện tượng cảm ứng điện từ khi từ thông qua cuộn dây biến thiên. (2) Máy phát điện xoay chiều biến điện năng thành cơ năng nhờ hiệu ứng nhiệt. (3) Cần kiểm tra điều kiện áp dụng. (4) Có thể bỏ qua đơn vị.
\shortans{2}
\loigiai{
Đối chiếu với kiến thức: Máy phát điện xoay chiều hoạt động dựa trên hiện tượng cảm ứng điện từ khi từ thông qua cuộn dây biến thiên. Có 2 phát biểu đúng.
}
\end{ex}

% ===== Câu 169 =====
% ID: L12C3B17-37-TN
% Mức: TH
\dangbt{Thiết lập biểu thức suất điện động xoay chiều do máy phát tạo ra}
\begin{ex}
% ID: L12C3B17-37-TN
% Mức: TH
Chọn phát biểu không trái với kiến thức về Thiết lập biểu thức suất điện động xoay chiều do máy phát tạo ra?
\choice
{Kết quả không phụ thuộc dữ kiện và điều kiện áp dụng.}
{\True Nếu Phi = Phi0 cos(omega t + phi) thì e = omega Phi0 sin(omega t + phi).}
{Suất điện động luôn cùng pha và bằng đúng từ thông.}
{Có thể bỏ qua phương, chiều và đơn vị trong mọi trường hợp.}
\loigiai{
Nếu Phi = Phi0 cos(omega t + phi) thì e = omega Phi0 sin(omega t + phi). Vì vậy phương án $B$ đúng; các phương án còn lại sai.
}
\end{ex}

% ===== Câu 170 =====
% ID: L12C3B17-38-DS
% Mức: VD
\dangbt{Tính từ thông qua khung dây quay trong từ trường}
\begin{ex}
% ID: L12C3B17-38-DS
% Mức: VD
Xét Tính từ thông qua khung dây quay trong từ trường (Tình huống 3). Hãy xác định tính đúng, sai của bốn phát biểu sau.
\choiceTF
{\True Khi giải cần xác định đúng dữ kiện, phương chiều, điều kiện áp dụng và đơn vị.}
{Kết quả không phụ thuộc dữ kiện và có thể bỏ qua điều kiện.}
{\True Khung $N$ vòng quay đều có từ thông Phi = NBS cos(omega t + phi).}
{Từ thông qua khung quay luôn không đổi và bằng NBS.}
\loigiai{
\begin{itemchoice}
\item (Đúng) Khi giải cần xác định đúng dữ kiện, phương chiều, điều kiện áp dụng và đơn vị. (Vì): Khi giải cần xác định đúng dữ kiện, phương chiều, điều kiện áp dụng và đơn vị. b) Sai: Kết quả không phụ thuộc dữ kiện và có thể bỏ qua điều kiện. c) Đúng: Khung $N$ vòng quay đều có từ thông Phi = NBS cos(omega t + phi). d) Sai: Từ thông qua khung quay luôn không đổi và bằng NBS. Kiến thức cốt lõi: Khung $N$ vòng quay đều có từ thông Phi = NBS cos(omega t + phi
\item (Sai) Kết quả không phụ thuộc dữ kiện và có thể bỏ qua điều kiện.
\item (Đúng) Khung $N$ vòng quay đều có từ thông Phi = NBS cos(omega t + phi).
\item (Sai) Từ thông qua khung quay luôn không đổi và bằng NBS.
\end{itemchoice}
}
\end{ex}

% ===== Câu 171 =====
% ID: L12C3B17-38-TL
% Mức: TH
\dangbt{Nhận biết cấu tạo và nguyên lí của máy phát điện xoay chiều}
\begin{bt}
% ID: L12C3B17-38-TL
% Mức: TH
Giải thích vì sao nhận định sau đúng: “Máy phát điện xoay chiều hoạt động dựa trên hiện tượng cảm ứng điện từ khi từ thông qua cuộn dây biến thiên.”
\loigiai{
Cần nêu được: Máy phát điện xoay chiều hoạt động dựa trên hiện tượng cảm ứng điện từ khi từ thông qua cuộn dây biến thiên. Khi vận dụng phải xác định đúng dữ kiện, phương chiều nếu có, điều kiện áp dụng, hệ đơn vị và kiểm tra tính hợp lí. Nhận định “Máy phát điện xoay chiều biến điện năng thành cơ năng nhờ hiệu ứng nhiệt.” sai.
}
\end{bt}

% ===== Câu 172 =====
% ID: L12C3B17-38-TLN
% Mức: TH
\begin{ex}
% ID: L12C3B17-38-TLN
% Mức: TH
Trong bốn phát biểu về Nhận biết cấu tạo và nguyên lí của máy phát điện xoay chiều, có bao nhiêu phát biểu đúng? (Chỉ ghi một số nguyên.) (1) Máy phát điện xoay chiều biến điện năng thành cơ năng nhờ hiệu ứng nhiệt. (2) Máy phát điện xoay chiều hoạt động dựa trên hiện tượng cảm ứng điện từ khi từ thông qua cuộn dây biến thiên. (3) Kết quả phải phù hợp ý nghĩa vật lí. (4) Mọi đại lượng đều là vô hướng.
\shortans{1}
\loigiai{
Đối chiếu với kiến thức: Máy phát điện xoay chiều hoạt động dựa trên hiện tượng cảm ứng điện từ khi từ thông qua cuộn dây biến thiên. Có 1 phát biểu đúng.
}
\end{ex}

% ===== Câu 173 =====
% ID: L12C3B17-38-TN
% Mức: VD
\dangbt{Thiết lập biểu thức suất điện động xoay chiều do máy phát tạo ra}
\begin{ex}
% ID: L12C3B17-38-TN
% Mức: VD
Cơ sở Vật lí đúng dùng để giải Thiết lập biểu thức suất điện động xoay chiều do máy phát tạo ra?
\choice
{Có thể bỏ qua phương, chiều và đơn vị trong mọi trường hợp.}
{Kết quả không phụ thuộc dữ kiện và điều kiện áp dụng.}
{\True Nếu Phi = Phi0 cos(omega t + phi) thì e = omega Phi0 sin(omega t + phi).}
{Suất điện động luôn cùng pha và bằng đúng từ thông.}
\loigiai{
Nếu Phi = Phi0 cos(omega t + phi) thì e = omega Phi0 sin(omega t + phi). Vì vậy phương án $C$ đúng; các phương án còn lại sai.
}
\end{ex}

% ===== Câu 174 =====
% ID: L12C3B17-39-DS
% Mức: VD
\dangbt{Tính từ thông qua khung dây quay trong từ trường}
\begin{ex}
% ID: L12C3B17-39-DS
% Mức: VD
Xét Tính từ thông qua khung dây quay trong từ trường (Tình huống 4). Hãy xác định tính đúng, sai của bốn phát biểu sau.
\choiceTF
{Kết quả không phụ thuộc dữ kiện và có thể bỏ qua điều kiện.}
{\True Khi giải cần xác định đúng dữ kiện, phương chiều, điều kiện áp dụng và đơn vị.}
{Từ thông qua khung quay luôn không đổi và bằng NBS.}
{\True Khung $N$ vòng quay đều có từ thông Phi = NBS cos(omega t + phi).}
\loigiai{
\begin{itemchoice}
\item (Sai) Kết quả không phụ thuộc dữ kiện và có thể bỏ qua điều kiện. (Vì): Kết quả không phụ thuộc dữ kiện và có thể bỏ qua điều kiện. b) Đúng: Khi giải cần xác định đúng dữ kiện, phương chiều, điều kiện áp dụng và đơn vị. c) Sai: Từ thông qua khung quay luôn không đổi và bằng NBS. d) Đúng: Khung $N$ vòng quay đều có từ thông Phi = NBS cos(omega t + phi). Kiến thức cốt lõi: Khung $N$ vòng quay đều có từ thông Phi = NBS cos(omega t + phi
\item (Đúng) Khi giải cần xác định đúng dữ kiện, phương chiều, điều kiện áp dụng và đơn vị.
\item (Sai) Từ thông qua khung quay luôn không đổi và bằng NBS.
\item (Đúng) Khung $N$ vòng quay đều có từ thông Phi = NBS cos(omega t + phi).
\end{itemchoice}
}
\end{ex}

% ===== Câu 175 =====
% ID: L12C3B17-39-TL
% Mức: VD
\dangbt{Nhận biết cấu tạo và nguyên lí của máy phát điện xoay chiều}
\begin{bt}
% ID: L12C3B17-39-TL
% Mức: VD
Phân tích sai lầm trong nhận định: “Máy phát điện xoay chiều biến điện năng thành cơ năng nhờ hiệu ứng nhiệt.”
\loigiai{
Cần nêu được: Máy phát điện xoay chiều hoạt động dựa trên hiện tượng cảm ứng điện từ khi từ thông qua cuộn dây biến thiên. Khi vận dụng phải xác định đúng dữ kiện, phương chiều nếu có, điều kiện áp dụng, hệ đơn vị và kiểm tra tính hợp lí. Nhận định “Máy phát điện xoay chiều biến điện năng thành cơ năng nhờ hiệu ứng nhiệt.” sai.
}
\end{bt}

% ===== Câu 176 =====
% ID: L12C3B17-39-TLN
% Mức: TH
\begin{ex}
% ID: L12C3B17-39-TLN
% Mức: TH
Trong bốn phát biểu về Nhận biết cấu tạo và nguyên lí của máy phát điện xoay chiều, có bao nhiêu phát biểu đúng? (Chỉ ghi một số nguyên.) (1) Máy phát điện xoay chiều hoạt động dựa trên hiện tượng cảm ứng điện từ khi từ thông qua cuộn dây biến thiên. (2) Máy phát điện xoay chiều biến điện năng thành cơ năng nhờ hiệu ứng nhiệt. (3) Cần kiểm tra điều kiện áp dụng. (4) Có thể bỏ qua đơn vị.
\shortans{2}
\loigiai{
Đối chiếu với kiến thức: Máy phát điện xoay chiều hoạt động dựa trên hiện tượng cảm ứng điện từ khi từ thông qua cuộn dây biến thiên. Có 2 phát biểu đúng.
}
\end{ex}

% ===== Câu 177 =====
% ID: L12C3B17-39-TN
% Mức: VD
\dangbt{Thiết lập biểu thức suất điện động xoay chiều do máy phát tạo ra}
\begin{ex}
% ID: L12C3B17-39-TN
% Mức: VD
Trong một bài thực hành về Thiết lập biểu thức suất điện động xoay chiều do máy phát tạo ra?
\choice
{Suất điện động luôn cùng pha và bằng đúng từ thông.}
{Có thể bỏ qua phương, chiều và đơn vị trong mọi trường hợp.}
{Kết quả không phụ thuộc dữ kiện và điều kiện áp dụng.}
{\True Nếu Phi = Phi0 cos(omega t + phi) thì e = omega Phi0 sin(omega t + phi).}
\loigiai{
Nếu Phi = Phi0 cos(omega t + phi) thì e = omega Phi0 sin(omega t + phi). Vì vậy phương án $D$ đúng; các phương án còn lại sai.
}
\end{ex}

% ===== Câu 178 =====
% ID: L12C3B17-40-DS
% Mức: VD
\dangbt{Tính từ thông qua khung dây quay trong từ trường}
\begin{ex}
% ID: L12C3B17-40-DS
% Mức: VD
Xét Tính từ thông qua khung dây quay trong từ trường (Tình huống 5). Hãy xác định tính đúng, sai của bốn phát biểu sau.
\choiceTF
{\True Khung $N$ vòng quay đều có từ thông Phi = NBS cos(omega t + phi).}
{Kết quả không phụ thuộc dữ kiện và có thể bỏ qua điều kiện.}
{Từ thông qua khung quay luôn không đổi và bằng NBS.}
{\True Khi giải cần xác định đúng dữ kiện, phương chiều, điều kiện áp dụng và đơn vị.}
\loigiai{
\begin{itemchoice}
\item (Đúng) Khung $N$ vòng quay đều có từ thông Phi = NBS cos(omega t + phi). (Vì): Khung $N$ vòng quay đều có từ thông Phi = NBS cos(omega t + phi). b) Sai: Kết quả không phụ thuộc dữ kiện và có thể bỏ qua điều kiện. c) Sai: Từ thông qua khung quay luôn không đổi và bằng NBS. d) Đúng: Khi giải cần xác định đúng dữ kiện, phương chiều, điều kiện áp dụng và đơn vị. Kiến thức cốt lõi: Khung $N$ vòng quay đều có từ thông Phi = NBS cos(omega t + phi
\item (Sai) Kết quả không phụ thuộc dữ kiện và có thể bỏ qua điều kiện.
\item (Sai) Từ thông qua khung quay luôn không đổi và bằng NBS.
\item (Đúng) Khi giải cần xác định đúng dữ kiện, phương chiều, điều kiện áp dụng và đơn vị.
\end{itemchoice}
}
\end{ex}

% ===== Câu 179 =====
% ID: L12C3B17-40-TL
% Mức: VD
\dangbt{Nhận biết cấu tạo và nguyên lí của máy phát điện xoay chiều}
\begin{bt}
% ID: L12C3B17-40-TL
% Mức: VD
Đề xuất các bước giải một bài toán thuộc dạng Nhận biết cấu tạo và nguyên lí của máy phát điện xoay chiều.
\loigiai{
Cần nêu được: Máy phát điện xoay chiều hoạt động dựa trên hiện tượng cảm ứng điện từ khi từ thông qua cuộn dây biến thiên. Khi vận dụng phải xác định đúng dữ kiện, phương chiều nếu có, điều kiện áp dụng, hệ đơn vị và kiểm tra tính hợp lí. Nhận định “Máy phát điện xoay chiều biến điện năng thành cơ năng nhờ hiệu ứng nhiệt.” sai.
}
\end{bt}

% ===== Câu 180 =====
% ID: L12C3B17-40-TLN
% Mức: VD
\begin{ex}
% ID: L12C3B17-40-TLN
% Mức: VD
Trong bốn phát biểu về Nhận biết cấu tạo và nguyên lí của máy phát điện xoay chiều, có bao nhiêu phát biểu đúng? (Chỉ ghi một số nguyên.) (1) Máy phát điện xoay chiều biến điện năng thành cơ năng nhờ hiệu ứng nhiệt. (2) Máy phát điện xoay chiều hoạt động dựa trên hiện tượng cảm ứng điện từ khi từ thông qua cuộn dây biến thiên. (3) Kết quả phải phù hợp ý nghĩa vật lí. (4) Mọi đại lượng đều là vô hướng.
\shortans{2}
\loigiai{
Đối chiếu với kiến thức: Máy phát điện xoay chiều hoạt động dựa trên hiện tượng cảm ứng điện từ khi từ thông qua cuộn dây biến thiên. Có 2 phát biểu đúng.
}
\end{ex}

% ===== Câu 181 =====
% ID: L12C3B17-40-TN
% Mức: VD
\dangbt{Thiết lập biểu thức suất điện động xoay chiều do máy phát tạo ra}
\begin{ex}
% ID: L12C3B17-40-TN
% Mức: VD
Kết luận đúng khi vận dụng Thiết lập biểu thức suất điện động xoay chiều do máy phát tạo ra?
\choice
{\True Nếu Phi = Phi0 cos(omega t + phi) thì e = omega Phi0 sin(omega t + phi).}
{Suất điện động luôn cùng pha và bằng đúng từ thông.}
{Có thể bỏ qua phương, chiều và đơn vị trong mọi trường hợp.}
{Kết quả không phụ thuộc dữ kiện và điều kiện áp dụng.}
\loigiai{
Nếu Phi = Phi0 cos(omega t + phi) thì e = omega Phi0 sin(omega t + phi). Vì vậy phương án $A$ đúng; các phương án còn lại sai.
}
\end{ex}

% ===== Câu 182 =====
% ID: L12C3B17-41-DS
% Mức: TH
\begin{ex}
% ID: L12C3B17-41-DS
% Mức: TH
Xét Thiết lập biểu thức suất điện động xoay chiều do máy phát tạo ra (Tình huống 1). Hãy xác định tính đúng, sai của bốn phát biểu sau.
\choiceTF
{\True Nếu Phi = Phi0 cos(omega t + phi) thì e = omega Phi0 sin(omega t + phi).}
{Suất điện động luôn cùng pha và bằng đúng từ thông.}
{\True Khi giải cần xác định đúng dữ kiện, phương chiều, điều kiện áp dụng và đơn vị.}
{Kết quả không phụ thuộc dữ kiện và có thể bỏ qua điều kiện.}
\loigiai{
\begin{itemchoice}
\item (Đúng) Nếu Phi = Phi0 cos(omega t + phi) thì e = omega Phi0 sin(omega t + phi). (Vì): Nếu Phi = Phi0 cos(omega t + phi) thì e = omega Phi0 sin(omega t + phi). b) Sai: Suất điện động luôn cùng pha và bằng đúng từ thông. c) Đúng: Khi giải cần xác định đúng dữ kiện, phương chiều, điều kiện áp dụng và đơn vị. d) Sai: Kết quả không phụ thuộc dữ kiện và có thể bỏ qua điều kiện. Kiến thức cốt lõi: Nếu Phi = Phi0 cos(omega t + phi) thì e = omega Phi0 sin(omega t + phi
\item (Sai) Suất điện động luôn cùng pha và bằng đúng từ thông.
\item (Đúng) Khi giải cần xác định đúng dữ kiện, phương chiều, điều kiện áp dụng và đơn vị.
\item (Sai) Kết quả không phụ thuộc dữ kiện và có thể bỏ qua điều kiện.
\end{itemchoice}
}
\end{ex}

% ===== Câu 183 =====
% ID: L12C3B17-41-TL
% Mức: VD
\dangbt{Nhận biết cấu tạo và nguyên lí của máy phát điện xoay chiều}
\begin{bt}
% ID: L12C3B17-41-TL
% Mức: VD
Nêu một tình huống thực tiễn liên quan đến Nhận biết cấu tạo và nguyên lí của máy phát điện xoay chiều và giải thích bằng kiến thức Vật lí.
\loigiai{
Cần nêu được: Máy phát điện xoay chiều hoạt động dựa trên hiện tượng cảm ứng điện từ khi từ thông qua cuộn dây biến thiên. Khi vận dụng phải xác định đúng dữ kiện, phương chiều nếu có, điều kiện áp dụng, hệ đơn vị và kiểm tra tính hợp lí. Nhận định “Máy phát điện xoay chiều biến điện năng thành cơ năng nhờ hiệu ứng nhiệt.” sai.
}
\end{bt}

% ===== Câu 184 =====
% ID: L12C3B17-41-TLN
% Mức: VD
\begin{ex}
% ID: L12C3B17-41-TLN
% Mức: VD
Trong bốn phát biểu về Nhận biết cấu tạo và nguyên lí của máy phát điện xoay chiều, có bao nhiêu phát biểu đúng? (Chỉ ghi một số nguyên.) (1) Máy phát điện xoay chiều hoạt động dựa trên hiện tượng cảm ứng điện từ khi từ thông qua cuộn dây biến thiên. (2) Máy phát điện xoay chiều biến điện năng thành cơ năng nhờ hiệu ứng nhiệt. (3) Cần kiểm tra điều kiện áp dụng. (4) Có thể bỏ qua đơn vị.
\shortans{2}
\loigiai{
Đối chiếu với kiến thức: Máy phát điện xoay chiều hoạt động dựa trên hiện tượng cảm ứng điện từ khi từ thông qua cuộn dây biến thiên. Có 2 phát biểu đúng.
}
\end{ex}

% ===== Câu 185 =====
% ID: L12C3B17-41-TN
% Mức: NB
\dangbt{Tính từ thông qua khung dây quay trong từ trường}
\begin{ex}
% ID: L12C3B17-41-TN
% Mức: NB
Phát biểu nào sau đây đúng về Tính từ thông qua khung dây quay trong từ trường?
\choice
{Từ thông qua khung quay luôn không đổi và bằng NBS.}
{Có thể bỏ qua phương, chiều và đơn vị trong mọi trường hợp.}
{Kết quả không phụ thuộc dữ kiện và điều kiện áp dụng.}
{\True Khung $N$ vòng quay đều có từ thông Phi = NBS cos(omega t + phi).}
\loigiai{
Khung $N$ vòng quay đều có từ thông Phi = NBS cos(omega t + phi). Vì vậy phương án $D$ đúng; các phương án còn lại sai.
}
\end{ex}

% ===== Câu 186 =====
% ID: L12C3B17-42-DS
% Mức: TH
\dangbt{Thiết lập biểu thức suất điện động xoay chiều do máy phát tạo ra}
\begin{ex}
% ID: L12C3B17-42-DS
% Mức: TH
Xét Thiết lập biểu thức suất điện động xoay chiều do máy phát tạo ra (Tình huống 2). Hãy xác định tính đúng, sai của bốn phát biểu sau.
\choiceTF
{Suất điện động luôn cùng pha và bằng đúng từ thông.}
{\True Nếu Phi = Phi0 cos(omega t + phi) thì e = omega Phi0 sin(omega t + phi).}
{Kết quả không phụ thuộc dữ kiện và có thể bỏ qua điều kiện.}
{\True Khi giải cần xác định đúng dữ kiện, phương chiều, điều kiện áp dụng và đơn vị.}
\loigiai{
\begin{itemchoice}
\item (Sai) Suất điện động luôn cùng pha và bằng đúng từ thông. (Vì): uất điện động luôn cùng pha và bằng đúng từ thông. b) Đúng: Nếu Phi = Phi0 cos(omega t + phi) thì e = omega Phi0 sin(omega t + phi). c) Sai: Kết quả không phụ thuộc dữ kiện và có thể bỏ qua điều kiện. d) Đúng: Khi giải cần xác định đúng dữ kiện, phương chiều, điều kiện áp dụng và đơn vị. Kiến thức cốt lõi: Nếu Phi = Phi0 cos(omega t + phi) thì e = omega Phi0 sin(omega t + phi
\item (Đúng) Nếu Phi = Phi0 cos(omega t + phi) thì e = omega Phi0 sin(omega t + phi).
\item (Sai) Kết quả không phụ thuộc dữ kiện và có thể bỏ qua điều kiện.
\item (Đúng) Khi giải cần xác định đúng dữ kiện, phương chiều, điều kiện áp dụng và đơn vị.
\end{itemchoice}
}
\end{ex}

% ===== Câu 187 =====
% ID: L12C3B17-42-TL
% Mức: VDC
\dangbt{Nhận biết cấu tạo và nguyên lí của máy phát điện xoay chiều}
\begin{bt}
% ID: L12C3B17-42-TL
% Mức: VDC
So sánh nhận định đúng “Máy phát điện xoay chiều hoạt động dựa trên hiện tượng cảm ứng điện từ khi từ thông qua cuộn dây biến thiên.” với nhận định sai “Máy phát điện xoay chiều biến điện năng thành cơ năng nhờ hiệu ứng nhiệt.”; chỉ rõ căn cứ Vật lí.
\loigiai{
Cần nêu được: Máy phát điện xoay chiều hoạt động dựa trên hiện tượng cảm ứng điện từ khi từ thông qua cuộn dây biến thiên. Khi vận dụng phải xác định đúng dữ kiện, phương chiều nếu có, điều kiện áp dụng, hệ đơn vị và kiểm tra tính hợp lí. Nhận định “Máy phát điện xoay chiều biến điện năng thành cơ năng nhờ hiệu ứng nhiệt.” sai.
}
\end{bt}

% ===== Câu 188 =====
% ID: L12C3B17-42-TLN
% Mức: VDC
\begin{ex}
% ID: L12C3B17-42-TLN
% Mức: VDC
Trong bốn phát biểu về Nhận biết cấu tạo và nguyên lí của máy phát điện xoay chiều, có bao nhiêu phát biểu đúng? (Chỉ ghi một số nguyên.) (1) Khi giải cần xác định đúng dữ kiện, phương chiều, điều kiện áp dụng và đơn vị. (2) Máy phát điện xoay chiều biến điện năng thành cơ năng nhờ hiệu ứng nhiệt. (3) Phải dùng đúng định luật cho tình huống. (4) Máy phát điện xoay chiều hoạt động dựa trên hiện tượng cảm ứng điện từ khi từ thông qua cuộn dây biến thiên.
\shortans{3}
\loigiai{
Đối chiếu với kiến thức: Máy phát điện xoay chiều hoạt động dựa trên hiện tượng cảm ứng điện từ khi từ thông qua cuộn dây biến thiên. Có 3 phát biểu đúng.
}
\end{ex}

% ===== Câu 189 =====
% ID: L12C3B17-42-TN
% Mức: NB
\dangbt{Tính từ thông qua khung dây quay trong từ trường}
\begin{ex}
% ID: L12C3B17-42-TN
% Mức: NB
Trong nội dung Tính từ thông qua khung dây quay trong từ trường?
\choice
{\True Khung $N$ vòng quay đều có từ thông Phi = NBS cos(omega t + phi).}
{Từ thông qua khung quay luôn không đổi và bằng NBS.}
{Có thể bỏ qua phương, chiều và đơn vị trong mọi trường hợp.}
{Kết quả không phụ thuộc dữ kiện và điều kiện áp dụng.}
\loigiai{
Khung $N$ vòng quay đều có từ thông Phi = NBS cos(omega t + phi). Vì vậy phương án $A$ đúng; các phương án còn lại sai.
}
\end{ex}

% ===== Câu 190 =====
% ID: L12C3B17-43-DS
% Mức: VD
\dangbt{Thiết lập biểu thức suất điện động xoay chiều do máy phát tạo ra}
\begin{ex}
% ID: L12C3B17-43-DS
% Mức: VD
Xét Thiết lập biểu thức suất điện động xoay chiều do máy phát tạo ra (Tình huống 3). Hãy xác định tính đúng, sai của bốn phát biểu sau.
\choiceTF
{\True Khi giải cần xác định đúng dữ kiện, phương chiều, điều kiện áp dụng và đơn vị.}
{Kết quả không phụ thuộc dữ kiện và có thể bỏ qua điều kiện.}
{\True Nếu Phi = Phi0 cos(omega t + phi) thì e = omega Phi0 sin(omega t + phi).}
{Suất điện động luôn cùng pha và bằng đúng từ thông.}
\loigiai{
\begin{itemchoice}
\item (Đúng) Khi giải cần xác định đúng dữ kiện, phương chiều, điều kiện áp dụng và đơn vị. (Vì): Khi giải cần xác định đúng dữ kiện, phương chiều, điều kiện áp dụng và đơn vị. b) Sai: Kết quả không phụ thuộc dữ kiện và có thể bỏ qua điều kiện. c) Đúng: Nếu Phi = Phi0 cos(omega t + phi) thì e = omega Phi0 sin(omega t + phi). d) Sai: Suất điện động luôn cùng pha và bằng đúng từ thông. Kiến thức cốt lõi: Nếu Phi = Phi0 cos(omega t + phi) thì e = omega Phi0 sin(omega t + phi
\item (Sai) Kết quả không phụ thuộc dữ kiện và có thể bỏ qua điều kiện.
\item (Đúng) Nếu Phi = Phi0 cos(omega t + phi) thì e = omega Phi0 sin(omega t + phi).
\item (Sai) Suất điện động luôn cùng pha và bằng đúng từ thông.
\end{itemchoice}
}
\end{ex}

% ===== Câu 191 =====
% ID: L12C3B17-43-TL
% Mức: TH
\dangbt{Tính từ thông qua khung dây quay trong từ trường}
\begin{bt}
% ID: L12C3B17-43-TL
% Mức: TH
Trình bày kiến thức cốt lõi của Tính từ thông qua khung dây quay trong từ trường và nêu điều kiện áp dụng.
\loigiai{
Cần nêu được: Khung $N$ vòng quay đều có từ thông Phi = NBS cos(omega t + phi). Khi vận dụng phải xác định đúng dữ kiện, phương chiều nếu có, điều kiện áp dụng, hệ đơn vị và kiểm tra tính hợp lí. Nhận định “Từ thông qua khung quay luôn không đổi và bằng NBS.” sai.
}
\end{bt}

% ===== Câu 192 =====
% ID: L12C3B17-43-TLN
% Mức: TH
\begin{ex}
% ID: L12C3B17-43-TLN
% Mức: TH
Trong bốn phát biểu về Tính từ thông qua khung dây quay trong từ trường, có bao nhiêu phát biểu đúng? (Chỉ ghi một số nguyên.) (1) Khung $N$ vòng quay đều có từ thông Phi = NBS cos(omega t + phi). (2) Từ thông qua khung quay luôn không đổi và bằng NBS. (3) Cần kiểm tra điều kiện áp dụng. (4) Có thể bỏ qua đơn vị.
\shortans{2}
\loigiai{
Đối chiếu với kiến thức: Khung $N$ vòng quay đều có từ thông Phi = NBS cos(omega t + phi). Có 2 phát biểu đúng.
}
\end{ex}

% ===== Câu 193 =====
% ID: L12C3B17-43-TN
% Mức: NB
\begin{ex}
% ID: L12C3B17-43-TN
% Mức: NB
Tình huống 3: chọn kết luận chính xác về Tính từ thông qua khung dây quay trong từ trường?
\choice
{Kết quả không phụ thuộc dữ kiện và điều kiện áp dụng.}
{\True Khung $N$ vòng quay đều có từ thông Phi = NBS cos(omega t + phi).}
{Từ thông qua khung quay luôn không đổi và bằng NBS.}
{Có thể bỏ qua phương, chiều và đơn vị trong mọi trường hợp.}
\loigiai{
Khung $N$ vòng quay đều có từ thông Phi = NBS cos(omega t + phi). Vì vậy phương án $B$ đúng; các phương án còn lại sai.
}
\end{ex}

% ===== Câu 194 =====
% ID: L12C3B17-44-DS
% Mức: VD
\dangbt{Thiết lập biểu thức suất điện động xoay chiều do máy phát tạo ra}
\begin{ex}
% ID: L12C3B17-44-DS
% Mức: VD
Xét Thiết lập biểu thức suất điện động xoay chiều do máy phát tạo ra (Tình huống 4). Hãy xác định tính đúng, sai của bốn phát biểu sau.
\choiceTF
{Kết quả không phụ thuộc dữ kiện và có thể bỏ qua điều kiện.}
{\True Khi giải cần xác định đúng dữ kiện, phương chiều, điều kiện áp dụng và đơn vị.}
{Suất điện động luôn cùng pha và bằng đúng từ thông.}
{\True Nếu Phi = Phi0 cos(omega t + phi) thì e = omega Phi0 sin(omega t + phi).}
\loigiai{
\begin{itemchoice}
\item (Sai) Kết quả không phụ thuộc dữ kiện và có thể bỏ qua điều kiện. (Vì): Kết quả không phụ thuộc dữ kiện và có thể bỏ qua điều kiện. b) Đúng: Khi giải cần xác định đúng dữ kiện, phương chiều, điều kiện áp dụng và đơn vị. c) Sai: Suất điện động luôn cùng pha và bằng đúng từ thông. d) Đúng: Nếu Phi = Phi0 cos(omega t + phi) thì e = omega Phi0 sin(omega t + phi). Kiến thức cốt lõi: Nếu Phi = Phi0 cos(omega t + phi) thì e = omega Phi0 sin(omega t + phi
\item (Đúng) Khi giải cần xác định đúng dữ kiện, phương chiều, điều kiện áp dụng và đơn vị.
\item (Sai) Suất điện động luôn cùng pha và bằng đúng từ thông.
\item (Đúng) Nếu Phi = Phi0 cos(omega t + phi) thì e = omega Phi0 sin(omega t + phi).
\end{itemchoice}
}
\end{ex}

% ===== Câu 195 =====
% ID: L12C3B17-44-TL
% Mức: TH
\dangbt{Tính từ thông qua khung dây quay trong từ trường}
\begin{bt}
% ID: L12C3B17-44-TL
% Mức: TH
Giải thích vì sao nhận định sau đúng: “Khung $N$ vòng quay đều có từ thông Phi = NBS cos(omega t + phi).”
\loigiai{
Cần nêu được: Khung $N$ vòng quay đều có từ thông Phi = NBS cos(omega t + phi). Khi vận dụng phải xác định đúng dữ kiện, phương chiều nếu có, điều kiện áp dụng, hệ đơn vị và kiểm tra tính hợp lí. Nhận định “Từ thông qua khung quay luôn không đổi và bằng NBS.” sai.
}
\end{bt}

% ===== Câu 196 =====
% ID: L12C3B17-44-TLN
% Mức: TH
\begin{ex}
% ID: L12C3B17-44-TLN
% Mức: TH
Trong bốn phát biểu về Tính từ thông qua khung dây quay trong từ trường, có bao nhiêu phát biểu đúng? (Chỉ ghi một số nguyên.) (1) Từ thông qua khung quay luôn không đổi và bằng NBS. (2) Khung $N$ vòng quay đều có từ thông Phi = NBS cos(omega t + phi). (3) Kết quả phải phù hợp ý nghĩa vật lí. (4) Mọi đại lượng đều là vô hướng.
\shortans{1}
\loigiai{
Đối chiếu với kiến thức: Khung $N$ vòng quay đều có từ thông Phi = NBS cos(omega t + phi). Có 1 phát biểu đúng.
}
\end{ex}

% ===== Câu 197 =====
% ID: L12C3B17-44-TN
% Mức: TH
\begin{ex}
% ID: L12C3B17-44-TN
% Mức: TH
Nội dung nào mô tả đúng nhất Tính từ thông qua khung dây quay trong từ trường?
\choice
{Có thể bỏ qua phương, chiều và đơn vị trong mọi trường hợp.}
{Kết quả không phụ thuộc dữ kiện và điều kiện áp dụng.}
{\True Khung $N$ vòng quay đều có từ thông Phi = NBS cos(omega t + phi).}
{Từ thông qua khung quay luôn không đổi và bằng NBS.}
\loigiai{
Khung $N$ vòng quay đều có từ thông Phi = NBS cos(omega t + phi). Vì vậy phương án $C$ đúng; các phương án còn lại sai.
}
\end{ex}

% ===== Câu 198 =====
% ID: L12C3B17-45-DS
% Mức: VD
\dangbt{Thiết lập biểu thức suất điện động xoay chiều do máy phát tạo ra}
\begin{ex}
% ID: L12C3B17-45-DS
% Mức: VD
Xét Thiết lập biểu thức suất điện động xoay chiều do máy phát tạo ra (Tình huống 5). Hãy xác định tính đúng, sai của bốn phát biểu sau.
\choiceTF
{\True Nếu Phi = Phi0 cos(omega t + phi) thì e = omega Phi0 sin(omega t + phi).}
{Kết quả không phụ thuộc dữ kiện và có thể bỏ qua điều kiện.}
{Suất điện động luôn cùng pha và bằng đúng từ thông.}
{\True Khi giải cần xác định đúng dữ kiện, phương chiều, điều kiện áp dụng và đơn vị.}
\loigiai{
\begin{itemchoice}
\item (Đúng) Nếu Phi = Phi0 cos(omega t + phi) thì e = omega Phi0 sin(omega t + phi). (Vì): Nếu Phi = Phi0 cos(omega t + phi) thì e = omega Phi0 sin(omega t + phi). b) Sai: Kết quả không phụ thuộc dữ kiện và có thể bỏ qua điều kiện. c) Sai: Suất điện động luôn cùng pha và bằng đúng từ thông. d) Đúng: Khi giải cần xác định đúng dữ kiện, phương chiều, điều kiện áp dụng và đơn vị. Kiến thức cốt lõi: Nếu Phi = Phi0 cos(omega t + phi) thì e = omega Phi0 sin(omega t + phi
\item (Sai) Kết quả không phụ thuộc dữ kiện và có thể bỏ qua điều kiện.
\item (Sai) Suất điện động luôn cùng pha và bằng đúng từ thông.
\item (Đúng) Khi giải cần xác định đúng dữ kiện, phương chiều, điều kiện áp dụng và đơn vị.
\end{itemchoice}
}
\end{ex}

% ===== Câu 199 =====
% ID: L12C3B17-45-TL
% Mức: VD
\dangbt{Tính từ thông qua khung dây quay trong từ trường}
\begin{bt}
% ID: L12C3B17-45-TL
% Mức: VD
Phân tích sai lầm trong nhận định: “Từ thông qua khung quay luôn không đổi và bằng NBS.”
\loigiai{
Cần nêu được: Khung $N$ vòng quay đều có từ thông Phi = NBS cos(omega t + phi). Khi vận dụng phải xác định đúng dữ kiện, phương chiều nếu có, điều kiện áp dụng, hệ đơn vị và kiểm tra tính hợp lí. Nhận định “Từ thông qua khung quay luôn không đổi và bằng NBS.” sai.
}
\end{bt}

% ===== Câu 200 =====
% ID: L12C3B17-45-TLN
% Mức: TH
\begin{ex}
% ID: L12C3B17-45-TLN
% Mức: TH
Trong bốn phát biểu về Tính từ thông qua khung dây quay trong từ trường, có bao nhiêu phát biểu đúng? (Chỉ ghi một số nguyên.) (1) Khung $N$ vòng quay đều có từ thông Phi = NBS cos(omega t + phi). (2) Từ thông qua khung quay luôn không đổi và bằng NBS. (3) Cần kiểm tra điều kiện áp dụng. (4) Có thể bỏ qua đơn vị.
\shortans{2}
\loigiai{
Đối chiếu với kiến thức: Khung $N$ vòng quay đều có từ thông Phi = NBS cos(omega t + phi). Có 2 phát biểu đúng.
}
\end{ex}

% ===== Câu 201 =====
% ID: L12C3B17-45-TN
% Mức: TH
\begin{ex}
% ID: L12C3B17-45-TN
% Mức: TH
Khi phân tích Tính từ thông qua khung dây quay trong từ trường?
\choice
{Từ thông qua khung quay luôn không đổi và bằng NBS.}
{Có thể bỏ qua phương, chiều và đơn vị trong mọi trường hợp.}
{Kết quả không phụ thuộc dữ kiện và điều kiện áp dụng.}
{\True Khung $N$ vòng quay đều có từ thông Phi = NBS cos(omega t + phi).}
\loigiai{
Khung $N$ vòng quay đều có từ thông Phi = NBS cos(omega t + phi). Vì vậy phương án $D$ đúng; các phương án còn lại sai.
}
\end{ex}

% ===== Câu 202 =====
% ID: L12C3B17-46-DS
% Mức: TH
\dangbt{Xác định biên độ, tần số, chu kì và pha của suất điện động xoay chiều}
\begin{ex}
% ID: L12C3B17-46-DS
% Mức: TH
Xét Xác định biên độ, tần số, chu kì và pha của suất điện động xoay chiều (Tình huống 1). Hãy xác định tính đúng, sai của bốn phát biểu sau.
\choiceTF
{\True Biên độ suất điện động là E0 = NBS omega và tần số f = omega/(2pi).}
{Biên độ suất điện động không phụ thuộc tốc độ quay.}
{\True Khi giải cần xác định đúng dữ kiện, phương chiều, điều kiện áp dụng và đơn vị.}
{Kết quả không phụ thuộc dữ kiện và có thể bỏ qua điều kiện.}
\loigiai{
\begin{itemchoice}
\item (Đúng) Biên độ suất điện động là E0 = NBS omega và tần số f = omega/(2pi). (Vì): Biên độ suất điện động là E0 = NBS omega và tần số f = omega/(2pi). b) Sai: Biên độ suất điện động không phụ thuộc tốc độ quay. c) Đúng: Khi giải cần xác định đúng dữ kiện, phương chiều, điều kiện áp dụng và đơn vị. d) Sai: Kết quả không phụ thuộc dữ kiện và có thể bỏ qua điều kiện. Kiến thức cốt lõi: Biên độ suất điện động là E0 = NBS omega và tần số f = omega/(2pi
\item (Sai) Biên độ suất điện động không phụ thuộc tốc độ quay.
\item (Đúng) Khi giải cần xác định đúng dữ kiện, phương chiều, điều kiện áp dụng và đơn vị.
\item (Sai) Kết quả không phụ thuộc dữ kiện và có thể bỏ qua điều kiện.
\end{itemchoice}
}
\end{ex}

% ===== Câu 203 =====
% ID: L12C3B17-46-TL
% Mức: VD
\dangbt{Tính từ thông qua khung dây quay trong từ trường}
\begin{bt}
% ID: L12C3B17-46-TL
% Mức: VD
Đề xuất các bước giải một bài toán thuộc dạng Tính từ thông qua khung dây quay trong từ trường.
\loigiai{
Cần nêu được: Khung $N$ vòng quay đều có từ thông Phi = NBS cos(omega t + phi). Khi vận dụng phải xác định đúng dữ kiện, phương chiều nếu có, điều kiện áp dụng, hệ đơn vị và kiểm tra tính hợp lí. Nhận định “Từ thông qua khung quay luôn không đổi và bằng NBS.” sai.
}
\end{bt}

% ===== Câu 204 =====
% ID: L12C3B17-46-TLN
% Mức: VD
\begin{ex}
% ID: L12C3B17-46-TLN
% Mức: VD
Trong bốn phát biểu về Tính từ thông qua khung dây quay trong từ trường, có bao nhiêu phát biểu đúng? (Chỉ ghi một số nguyên.) (1) Từ thông qua khung quay luôn không đổi và bằng NBS. (2) Khung $N$ vòng quay đều có từ thông Phi = NBS cos(omega t + phi). (3) Kết quả phải phù hợp ý nghĩa vật lí. (4) Mọi đại lượng đều là vô hướng.
\shortans{2}
\loigiai{
Đối chiếu với kiến thức: Khung $N$ vòng quay đều có từ thông Phi = NBS cos(omega t + phi). Có 2 phát biểu đúng.
}
\end{ex}

% ===== Câu 205 =====
% ID: L12C3B17-46-TN
% Mức: TH
\begin{ex}
% ID: L12C3B17-46-TN
% Mức: TH
Một học sinh ôn tập Tính từ thông qua khung dây quay trong từ trường?
\choice
{\True Khung $N$ vòng quay đều có từ thông Phi = NBS cos(omega t + phi).}
{Từ thông qua khung quay luôn không đổi và bằng NBS.}
{Có thể bỏ qua phương, chiều và đơn vị trong mọi trường hợp.}
{Kết quả không phụ thuộc dữ kiện và điều kiện áp dụng.}
\loigiai{
Khung $N$ vòng quay đều có từ thông Phi = NBS cos(omega t + phi). Vì vậy phương án $A$ đúng; các phương án còn lại sai.
}
\end{ex}

% ===== Câu 206 =====
% ID: L12C3B17-47-DS
% Mức: TH
\dangbt{Xác định biên độ, tần số, chu kì và pha của suất điện động xoay chiều}
\begin{ex}
% ID: L12C3B17-47-DS
% Mức: TH
Xét Xác định biên độ, tần số, chu kì và pha của suất điện động xoay chiều (Tình huống 2). Hãy xác định tính đúng, sai của bốn phát biểu sau.
\choiceTF
{Biên độ suất điện động không phụ thuộc tốc độ quay.}
{\True Biên độ suất điện động là E0 = NBS omega và tần số f = omega/(2pi).}
{Kết quả không phụ thuộc dữ kiện và có thể bỏ qua điều kiện.}
{\True Khi giải cần xác định đúng dữ kiện, phương chiều, điều kiện áp dụng và đơn vị.}
\loigiai{
\begin{itemchoice}
\item (Sai) Biên độ suất điện động không phụ thuộc tốc độ quay. (Vì): Biên độ suất điện động không phụ thuộc tốc độ quay. b) Đúng: Biên độ suất điện động là E0 = NBS omega và tần số f = omega/(2pi). c) Sai: Kết quả không phụ thuộc dữ kiện và có thể bỏ qua điều kiện. d) Đúng: Khi giải cần xác định đúng dữ kiện, phương chiều, điều kiện áp dụng và đơn vị. Kiến thức cốt lõi: Biên độ suất điện động là E0 = NBS omega và tần số f = omega/(2pi
\item (Đúng) Biên độ suất điện động là E0 = NBS omega và tần số f = omega/(2pi).
\item (Sai) Kết quả không phụ thuộc dữ kiện và có thể bỏ qua điều kiện.
\item (Đúng) Khi giải cần xác định đúng dữ kiện, phương chiều, điều kiện áp dụng và đơn vị.
\end{itemchoice}
}
\end{ex}

% ===== Câu 207 =====
% ID: L12C3B17-47-TL
% Mức: VD
\dangbt{Tính từ thông qua khung dây quay trong từ trường}
\begin{bt}
% ID: L12C3B17-47-TL
% Mức: VD
Nêu một tình huống thực tiễn liên quan đến Tính từ thông qua khung dây quay trong từ trường và giải thích bằng kiến thức Vật lí.
\loigiai{
Cần nêu được: Khung $N$ vòng quay đều có từ thông Phi = NBS cos(omega t + phi). Khi vận dụng phải xác định đúng dữ kiện, phương chiều nếu có, điều kiện áp dụng, hệ đơn vị và kiểm tra tính hợp lí. Nhận định “Từ thông qua khung quay luôn không đổi và bằng NBS.” sai.
}
\end{bt}

% ===== Câu 208 =====
% ID: L12C3B17-47-TLN
% Mức: VD
\begin{ex}
% ID: L12C3B17-47-TLN
% Mức: VD
Trong bốn phát biểu về Tính từ thông qua khung dây quay trong từ trường, có bao nhiêu phát biểu đúng? (Chỉ ghi một số nguyên.) (1) Khung $N$ vòng quay đều có từ thông Phi = NBS cos(omega t + phi). (2) Từ thông qua khung quay luôn không đổi và bằng NBS. (3) Cần kiểm tra điều kiện áp dụng. (4) Có thể bỏ qua đơn vị.
\shortans{2}
\loigiai{
Đối chiếu với kiến thức: Khung $N$ vòng quay đều có từ thông Phi = NBS cos(omega t + phi). Có 2 phát biểu đúng.
}
\end{ex}

% ===== Câu 209 =====
% ID: L12C3B17-47-TN
% Mức: TH
\begin{ex}
% ID: L12C3B17-47-TN
% Mức: TH
Chọn phát biểu không trái với kiến thức về Tính từ thông qua khung dây quay trong từ trường?
\choice
{Kết quả không phụ thuộc dữ kiện và điều kiện áp dụng.}
{\True Khung $N$ vòng quay đều có từ thông Phi = NBS cos(omega t + phi).}
{Từ thông qua khung quay luôn không đổi và bằng NBS.}
{Có thể bỏ qua phương, chiều và đơn vị trong mọi trường hợp.}
\loigiai{
Khung $N$ vòng quay đều có từ thông Phi = NBS cos(omega t + phi). Vì vậy phương án $B$ đúng; các phương án còn lại sai.
}
\end{ex}

% ===== Câu 210 =====
% ID: L12C3B17-48-DS
% Mức: VD
\dangbt{Xác định biên độ, tần số, chu kì và pha của suất điện động xoay chiều}
\begin{ex}
% ID: L12C3B17-48-DS
% Mức: VD
Xét Xác định biên độ, tần số, chu kì và pha của suất điện động xoay chiều (Tình huống 3). Hãy xác định tính đúng, sai của bốn phát biểu sau.
\choiceTF
{\True Khi giải cần xác định đúng dữ kiện, phương chiều, điều kiện áp dụng và đơn vị.}
{Kết quả không phụ thuộc dữ kiện và có thể bỏ qua điều kiện.}
{\True Biên độ suất điện động là E0 = NBS omega và tần số f = omega/(2pi).}
{Biên độ suất điện động không phụ thuộc tốc độ quay.}
\loigiai{
\begin{itemchoice}
\item (Đúng) Khi giải cần xác định đúng dữ kiện, phương chiều, điều kiện áp dụng và đơn vị. (Vì): Khi giải cần xác định đúng dữ kiện, phương chiều, điều kiện áp dụng và đơn vị. b) Sai: Kết quả không phụ thuộc dữ kiện và có thể bỏ qua điều kiện. c) Đúng: Biên độ suất điện động là E0 = NBS omega và tần số f = omega/(2pi). d) Sai: Biên độ suất điện động không phụ thuộc tốc độ quay. Kiến thức cốt lõi: Biên độ suất điện động là E0 = NBS omega và tần số f = omega/(2pi
\item (Sai) Kết quả không phụ thuộc dữ kiện và có thể bỏ qua điều kiện.
\item (Đúng) Biên độ suất điện động là E0 = NBS omega và tần số f = omega/(2pi).
\item (Sai) Biên độ suất điện động không phụ thuộc tốc độ quay.
\end{itemchoice}
}
\end{ex}

% ===== Câu 211 =====
% ID: L12C3B17-48-TL
% Mức: VDC
\dangbt{Tính từ thông qua khung dây quay trong từ trường}
\begin{bt}
% ID: L12C3B17-48-TL
% Mức: VDC
So sánh nhận định đúng “Khung $N$ vòng quay đều có từ thông Phi = NBS cos(omega t + phi).” với nhận định sai “Từ thông qua khung quay luôn không đổi và bằng NBS.”; chỉ rõ căn cứ Vật lí.
\loigiai{
Cần nêu được: Khung $N$ vòng quay đều có từ thông Phi = NBS cos(omega t + phi). Khi vận dụng phải xác định đúng dữ kiện, phương chiều nếu có, điều kiện áp dụng, hệ đơn vị và kiểm tra tính hợp lí. Nhận định “Từ thông qua khung quay luôn không đổi và bằng NBS.” sai.
}
\end{bt}

% ===== Câu 212 =====
% ID: L12C3B17-48-TLN
% Mức: VDC
\begin{ex}
% ID: L12C3B17-48-TLN
% Mức: VDC
Trong bốn phát biểu về Tính từ thông qua khung dây quay trong từ trường, có bao nhiêu phát biểu đúng? (Chỉ ghi một số nguyên.) (1) Khi giải cần xác định đúng dữ kiện, phương chiều, điều kiện áp dụng và đơn vị. (2) Từ thông qua khung quay luôn không đổi và bằng NBS. (3) Phải dùng đúng định luật cho tình huống. (4) Khung $N$ vòng quay đều có từ thông Phi = NBS cos(omega t + phi).
\shortans{3}
\loigiai{
Đối chiếu với kiến thức: Khung $N$ vòng quay đều có từ thông Phi = NBS cos(omega t + phi). Có 3 phát biểu đúng.
}
\end{ex}

% ===== Câu 213 =====
% ID: L12C3B17-48-TN
% Mức: VD
\begin{ex}
% ID: L12C3B17-48-TN
% Mức: VD
Cơ sở Vật lí đúng dùng để giải Tính từ thông qua khung dây quay trong từ trường?
\choice
{Có thể bỏ qua phương, chiều và đơn vị trong mọi trường hợp.}
{Kết quả không phụ thuộc dữ kiện và điều kiện áp dụng.}
{\True Khung $N$ vòng quay đều có từ thông Phi = NBS cos(omega t + phi).}
{Từ thông qua khung quay luôn không đổi và bằng NBS.}
\loigiai{
Khung $N$ vòng quay đều có từ thông Phi = NBS cos(omega t + phi). Vì vậy phương án $C$ đúng; các phương án còn lại sai.
}
\end{ex}

% ===== Câu 214 =====
% ID: L12C3B17-49-DS
% Mức: VD
\dangbt{Xác định biên độ, tần số, chu kì và pha của suất điện động xoay chiều}
\begin{ex}
% ID: L12C3B17-49-DS
% Mức: VD
Xét Xác định biên độ, tần số, chu kì và pha của suất điện động xoay chiều (Tình huống 4). Hãy xác định tính đúng, sai của bốn phát biểu sau.
\choiceTF
{Kết quả không phụ thuộc dữ kiện và có thể bỏ qua điều kiện.}
{\True Khi giải cần xác định đúng dữ kiện, phương chiều, điều kiện áp dụng và đơn vị.}
{Biên độ suất điện động không phụ thuộc tốc độ quay.}
{\True Biên độ suất điện động là E0 = NBS omega và tần số f = omega/(2pi).}
\loigiai{
\begin{itemchoice}
\item (Sai) Kết quả không phụ thuộc dữ kiện và có thể bỏ qua điều kiện. (Vì): Kết quả không phụ thuộc dữ kiện và có thể bỏ qua điều kiện. b) Đúng: Khi giải cần xác định đúng dữ kiện, phương chiều, điều kiện áp dụng và đơn vị. c) Sai: Biên độ suất điện động không phụ thuộc tốc độ quay. d) Đúng: Biên độ suất điện động là E0 = NBS omega và tần số f = omega/(2pi). Kiến thức cốt lõi: Biên độ suất điện động là E0 = NBS omega và tần số f = omega/(2pi
\item (Đúng) Khi giải cần xác định đúng dữ kiện, phương chiều, điều kiện áp dụng và đơn vị.
\item (Sai) Biên độ suất điện động không phụ thuộc tốc độ quay.
\item (Đúng) Biên độ suất điện động là E0 = NBS omega và tần số f = omega/(2pi).
\end{itemchoice}
}
\end{ex}

% ===== Câu 215 =====
% ID: L12C3B17-49-TL
% Mức: TH
\dangbt{Thiết lập biểu thức suất điện động xoay chiều do máy phát tạo ra}
\begin{bt}
% ID: L12C3B17-49-TL
% Mức: TH
Trình bày kiến thức cốt lõi của Thiết lập biểu thức suất điện động xoay chiều do máy phát tạo ra và nêu điều kiện áp dụng.
\loigiai{
Cần nêu được: Nếu Phi = Phi0 cos(omega t + phi) thì e = omega Phi0 sin(omega t + phi). Khi vận dụng phải xác định đúng dữ kiện, phương chiều nếu có, điều kiện áp dụng, hệ đơn vị và kiểm tra tính hợp lí. Nhận định “Suất điện động luôn cùng pha và bằng đúng từ thông.” sai.
}
\end{bt}

% ===== Câu 216 =====
% ID: L12C3B17-49-TLN
% Mức: TH
\begin{ex}
% ID: L12C3B17-49-TLN
% Mức: TH
Trong bốn phát biểu về Thiết lập biểu thức suất điện động xoay chiều do máy phát tạo ra, có bao nhiêu phát biểu đúng? (Chỉ ghi một số nguyên.) (1) Nếu Phi = Phi0 cos(omega t + phi) thì e = omega Phi0 sin(omega t + phi). (2) Suất điện động luôn cùng pha và bằng đúng từ thông. (3) Cần kiểm tra điều kiện áp dụng. (4) Có thể bỏ qua đơn vị.
\shortans{2}
\loigiai{
Đối chiếu với kiến thức: Nếu Phi = Phi0 cos(omega t + phi) thì e = omega Phi0 sin(omega t + phi). Có 2 phát biểu đúng.
}
\end{ex}

% ===== Câu 217 =====
% ID: L12C3B17-49-TN
% Mức: VD
\dangbt{Tính từ thông qua khung dây quay trong từ trường}
\begin{ex}
% ID: L12C3B17-49-TN
% Mức: VD
Trong một bài thực hành về Tính từ thông qua khung dây quay trong từ trường?
\choice
{Từ thông qua khung quay luôn không đổi và bằng NBS.}
{Có thể bỏ qua phương, chiều và đơn vị trong mọi trường hợp.}
{Kết quả không phụ thuộc dữ kiện và điều kiện áp dụng.}
{\True Khung $N$ vòng quay đều có từ thông Phi = NBS cos(omega t + phi).}
\loigiai{
Khung $N$ vòng quay đều có từ thông Phi = NBS cos(omega t + phi). Vì vậy phương án $D$ đúng; các phương án còn lại sai.
}
\end{ex}

% ===== Câu 218 =====
% ID: L12C3B17-50-DS
% Mức: VD
\dangbt{Xác định biên độ, tần số, chu kì và pha của suất điện động xoay chiều}
\begin{ex}
% ID: L12C3B17-50-DS
% Mức: VD
Xét Xác định biên độ, tần số, chu kì và pha của suất điện động xoay chiều (Tình huống 5). Hãy xác định tính đúng, sai của bốn phát biểu sau.
\choiceTF
{\True Biên độ suất điện động là E0 = NBS omega và tần số f = omega/(2pi).}
{Kết quả không phụ thuộc dữ kiện và có thể bỏ qua điều kiện.}
{Biên độ suất điện động không phụ thuộc tốc độ quay.}
{\True Khi giải cần xác định đúng dữ kiện, phương chiều, điều kiện áp dụng và đơn vị.}
\loigiai{
\begin{itemchoice}
\item (Đúng) Biên độ suất điện động là E0 = NBS omega và tần số f = omega/(2pi). (Vì): Biên độ suất điện động là E0 = NBS omega và tần số f = omega/(2pi). b) Sai: Kết quả không phụ thuộc dữ kiện và có thể bỏ qua điều kiện. c) Sai: Biên độ suất điện động không phụ thuộc tốc độ quay. d) Đúng: Khi giải cần xác định đúng dữ kiện, phương chiều, điều kiện áp dụng và đơn vị. Kiến thức cốt lõi: Biên độ suất điện động là E0 = NBS omega và tần số f = omega/(2pi
\item (Sai) Kết quả không phụ thuộc dữ kiện và có thể bỏ qua điều kiện.
\item (Sai) Biên độ suất điện động không phụ thuộc tốc độ quay.
\item (Đúng) Khi giải cần xác định đúng dữ kiện, phương chiều, điều kiện áp dụng và đơn vị.
\end{itemchoice}
}
\end{ex}

% ===== Câu 219 =====
% ID: L12C3B17-50-TL
% Mức: TH
\dangbt{Thiết lập biểu thức suất điện động xoay chiều do máy phát tạo ra}
\begin{bt}
% ID: L12C3B17-50-TL
% Mức: TH
Giải thích vì sao nhận định sau đúng: “Nếu Phi = Phi0 cos(omega t + phi) thì e = omega Phi0 sin(omega t + phi).”
\loigiai{
Cần nêu được: Nếu Phi = Phi0 cos(omega t + phi) thì e = omega Phi0 sin(omega t + phi). Khi vận dụng phải xác định đúng dữ kiện, phương chiều nếu có, điều kiện áp dụng, hệ đơn vị và kiểm tra tính hợp lí. Nhận định “Suất điện động luôn cùng pha và bằng đúng từ thông.” sai.
}
\end{bt}

% ===== Câu 220 =====
% ID: L12C3B17-50-TLN
% Mức: TH
\begin{ex}
% ID: L12C3B17-50-TLN
% Mức: TH
Trong bốn phát biểu về Thiết lập biểu thức suất điện động xoay chiều do máy phát tạo ra, có bao nhiêu phát biểu đúng? (Chỉ ghi một số nguyên.) (1) Suất điện động luôn cùng pha và bằng đúng từ thông. (2) Nếu Phi = Phi0 cos(omega t + phi) thì e = omega Phi0 sin(omega t + phi). (3) Kết quả phải phù hợp ý nghĩa vật lí. (4) Mọi đại lượng đều là vô hướng.
\shortans{1}
\loigiai{
Đối chiếu với kiến thức: Nếu Phi = Phi0 cos(omega t + phi) thì e = omega Phi0 sin(omega t + phi). Có 1 phát biểu đúng.
}
\end{ex}

% ===== Câu 221 =====
% ID: L12C3B17-50-TN
% Mức: VD
\dangbt{Tính từ thông qua khung dây quay trong từ trường}
\begin{ex}
% ID: L12C3B17-50-TN
% Mức: VD
Kết luận đúng khi vận dụng Tính từ thông qua khung dây quay trong từ trường?
\choice
{\True Khung $N$ vòng quay đều có từ thông Phi = NBS cos(omega t + phi).}
{Từ thông qua khung quay luôn không đổi và bằng NBS.}
{Có thể bỏ qua phương, chiều và đơn vị trong mọi trường hợp.}
{Kết quả không phụ thuộc dữ kiện và điều kiện áp dụng.}
\loigiai{
Khung $N$ vòng quay đều có từ thông Phi = NBS cos(omega t + phi). Vì vậy phương án $A$ đúng; các phương án còn lại sai.
}
\end{ex}

% ===== Câu 222 =====
% ID: L12C3B17-51-DS
% Mức: TH
\dangbt{Khai thác đồ thị từ thông và suất điện động theo thời gian}
\begin{ex}
% ID: L12C3B17-51-DS
% Mức: TH
Xét Khai thác đồ thị từ thông và suất điện động theo thời gian (Tình huống 1). Hãy xác định tính đúng, sai của bốn phát biểu sau.
\choiceTF
{\True Suất điện động lệch pha pi/2 so với từ thông; khi từ thông cực đại thì suất điện động bằng 0.}
{Từ thông và suất điện động luôn đạt cực đại đồng thời.}
{\True Khi giải cần xác định đúng dữ kiện, phương chiều, điều kiện áp dụng và đơn vị.}
{Kết quả không phụ thuộc dữ kiện và có thể bỏ qua điều kiện.}
\loigiai{
\begin{itemchoice}
\item (Đúng) Suất điện động lệch pha pi/2 so với từ thông; khi từ thông cực đại thì suất điện động bằng 0. (Vì): Suất điện động lệch pha pi/2 so với từ thông; khi từ thông cực đại thì suất điện động bằng 0. b) Sai: Từ thông và suất điện động luôn đạt cực đại đồng thời. c) Đúng: Khi giải cần xác định đúng dữ kiện, phương chiều, điều kiện áp dụng và đơn vị. d) Sai: Kết quả không phụ thuộc dữ kiện và có thể bỏ qua điều kiện. Kiến thức cốt lõi: Suất điện động lệch pha pi/2 so với từ thông; khi từ thông cực đại thì suất điện động bằng 0
\item (Sai) Từ thông và suất điện động luôn đạt cực đại đồng thời.
\item (Đúng) Khi giải cần xác định đúng dữ kiện, phương chiều, điều kiện áp dụng và đơn vị.
\item (Sai) Kết quả không phụ thuộc dữ kiện và có thể bỏ qua điều kiện.
\end{itemchoice}
}
\end{ex}

% ===== Câu 223 =====
% ID: L12C3B17-51-TL
% Mức: VD
\dangbt{Thiết lập biểu thức suất điện động xoay chiều do máy phát tạo ra}
\begin{bt}
% ID: L12C3B17-51-TL
% Mức: VD
Phân tích sai lầm trong nhận định: “Suất điện động luôn cùng pha và bằng đúng từ thông.”
\loigiai{
Cần nêu được: Nếu Phi = Phi0 cos(omega t + phi) thì e = omega Phi0 sin(omega t + phi). Khi vận dụng phải xác định đúng dữ kiện, phương chiều nếu có, điều kiện áp dụng, hệ đơn vị và kiểm tra tính hợp lí. Nhận định “Suất điện động luôn cùng pha và bằng đúng từ thông.” sai.
}
\end{bt}

% ===== Câu 224 =====
% ID: L12C3B17-51-TLN
% Mức: TH
\begin{ex}
% ID: L12C3B17-51-TLN
% Mức: TH
Trong bốn phát biểu về Thiết lập biểu thức suất điện động xoay chiều do máy phát tạo ra, có bao nhiêu phát biểu đúng? (Chỉ ghi một số nguyên.) (1) Nếu Phi = Phi0 cos(omega t + phi) thì e = omega Phi0 sin(omega t + phi). (2) Suất điện động luôn cùng pha và bằng đúng từ thông. (3) Cần kiểm tra điều kiện áp dụng. (4) Có thể bỏ qua đơn vị.
\shortans{2}
\loigiai{
Đối chiếu với kiến thức: Nếu Phi = Phi0 cos(omega t + phi) thì e = omega Phi0 sin(omega t + phi). Có 2 phát biểu đúng.
}
\end{ex}

% ===== Câu 225 =====
% ID: L12C3B17-51-TN
% Mức: NB
\dangbt{Xác định biên độ, tần số, chu kì và pha của suất điện động xoay chiều}
\begin{ex}
% ID: L12C3B17-51-TN
% Mức: NB
Phát biểu nào sau đây đúng về Xác định biên độ, tần số, chu kì và pha của suất điện động xoay chiều?
\choice
{Biên độ suất điện động không phụ thuộc tốc độ quay.}
{Có thể bỏ qua phương, chiều và đơn vị trong mọi trường hợp.}
{Kết quả không phụ thuộc dữ kiện và điều kiện áp dụng.}
{\True Biên độ suất điện động là E0 = NBS omega và tần số f = omega/(2pi).}
\loigiai{
Biên độ suất điện động là E0 = NBS omega và tần số f = omega/(2pi). Vì vậy phương án $D$ đúng; các phương án còn lại sai.
}
\end{ex}

% ===== Câu 226 =====
% ID: L12C3B17-52-DS
% Mức: TH
\dangbt{Khai thác đồ thị từ thông và suất điện động theo thời gian}
\begin{ex}
% ID: L12C3B17-52-DS
% Mức: TH
Xét Khai thác đồ thị từ thông và suất điện động theo thời gian (Tình huống 2). Hãy xác định tính đúng, sai của bốn phát biểu sau.
\choiceTF
{Từ thông và suất điện động luôn đạt cực đại đồng thời.}
{\True Suất điện động lệch pha pi/2 so với từ thông; khi từ thông cực đại thì suất điện động bằng 0.}
{Kết quả không phụ thuộc dữ kiện và có thể bỏ qua điều kiện.}
{\True Khi giải cần xác định đúng dữ kiện, phương chiều, điều kiện áp dụng và đơn vị.}
\loigiai{
\begin{itemchoice}
\item (Sai) Từ thông và suất điện động luôn đạt cực đại đồng thời. (Vì): Từ thông và suất điện động luôn đạt cực đại đồng thời. b) Đúng: Suất điện động lệch pha pi/2 so với từ thông; khi từ thông cực đại thì suất điện động bằng 0. c) Sai: Kết quả không phụ thuộc dữ kiện và có thể bỏ qua điều kiện. d) Đúng: Khi giải cần xác định đúng dữ kiện, phương chiều, điều kiện áp dụng và đơn vị. Kiến thức cốt lõi: Suất điện động lệch pha pi/2 so với từ thông; khi từ thông cực đại thì suất điện động bằng 0
\item (Đúng) Suất điện động lệch pha pi/2 so với từ thông; khi từ thông cực đại thì suất điện động bằng 0.
\item (Sai) Kết quả không phụ thuộc dữ kiện và có thể bỏ qua điều kiện.
\item (Đúng) Khi giải cần xác định đúng dữ kiện, phương chiều, điều kiện áp dụng và đơn vị.
\end{itemchoice}
}
\end{ex}

% ===== Câu 227 =====
% ID: L12C3B17-52-TL
% Mức: VD
\dangbt{Thiết lập biểu thức suất điện động xoay chiều do máy phát tạo ra}
\begin{bt}
% ID: L12C3B17-52-TL
% Mức: VD
Đề xuất các bước giải một bài toán thuộc dạng Thiết lập biểu thức suất điện động xoay chiều do máy phát tạo ra.
\loigiai{
Cần nêu được: Nếu Phi = Phi0 cos(omega t + phi) thì e = omega Phi0 sin(omega t + phi). Khi vận dụng phải xác định đúng dữ kiện, phương chiều nếu có, điều kiện áp dụng, hệ đơn vị và kiểm tra tính hợp lí. Nhận định “Suất điện động luôn cùng pha và bằng đúng từ thông.” sai.
}
\end{bt}

% ===== Câu 228 =====
% ID: L12C3B17-52-TLN
% Mức: VD
\begin{ex}
% ID: L12C3B17-52-TLN
% Mức: VD
Trong bốn phát biểu về Thiết lập biểu thức suất điện động xoay chiều do máy phát tạo ra, có bao nhiêu phát biểu đúng? (Chỉ ghi một số nguyên.) (1) Suất điện động luôn cùng pha và bằng đúng từ thông. (2) Nếu Phi = Phi0 cos(omega t + phi) thì e = omega Phi0 sin(omega t + phi). (3) Kết quả phải phù hợp ý nghĩa vật lí. (4) Mọi đại lượng đều là vô hướng.
\shortans{2}
\loigiai{
Đối chiếu với kiến thức: Nếu Phi = Phi0 cos(omega t + phi) thì e = omega Phi0 sin(omega t + phi). Có 2 phát biểu đúng.
}
\end{ex}

% ===== Câu 229 =====
% ID: L12C3B17-52-TN
% Mức: NB
\dangbt{Xác định biên độ, tần số, chu kì và pha của suất điện động xoay chiều}
\begin{ex}
% ID: L12C3B17-52-TN
% Mức: NB
Trong nội dung Xác định biên độ, tần số, chu kì và pha của suất điện động xoay chiều?
\choice
{\True Biên độ suất điện động là E0 = NBS omega và tần số f = omega/(2pi).}
{Biên độ suất điện động không phụ thuộc tốc độ quay.}
{Có thể bỏ qua phương, chiều và đơn vị trong mọi trường hợp.}
{Kết quả không phụ thuộc dữ kiện và điều kiện áp dụng.}
\loigiai{
Biên độ suất điện động là E0 = NBS omega và tần số f = omega/(2pi). Vì vậy phương án $A$ đúng; các phương án còn lại sai.
}
\end{ex}

% ===== Câu 230 =====
% ID: L12C3B17-53-DS
% Mức: VD
\dangbt{Khai thác đồ thị từ thông và suất điện động theo thời gian}
\begin{ex}
% ID: L12C3B17-53-DS
% Mức: VD
Xét Khai thác đồ thị từ thông và suất điện động theo thời gian (Tình huống 3). Hãy xác định tính đúng, sai của bốn phát biểu sau.
\choiceTF
{\True Khi giải cần xác định đúng dữ kiện, phương chiều, điều kiện áp dụng và đơn vị.}
{Kết quả không phụ thuộc dữ kiện và có thể bỏ qua điều kiện.}
{\True Suất điện động lệch pha pi/2 so với từ thông; khi từ thông cực đại thì suất điện động bằng 0.}
{Từ thông và suất điện động luôn đạt cực đại đồng thời.}
\loigiai{
\begin{itemchoice}
\item (Đúng) Khi giải cần xác định đúng dữ kiện, phương chiều, điều kiện áp dụng và đơn vị. (Vì): Khi giải cần xác định đúng dữ kiện, phương chiều, điều kiện áp dụng và đơn vị. b) Sai: Kết quả không phụ thuộc dữ kiện và có thể bỏ qua điều kiện. c) Đúng: Suất điện động lệch pha pi/2 so với từ thông; khi từ thông cực đại thì suất điện động bằng 0. d) Sai: Từ thông và suất điện động luôn đạt cực đại đồng thời. Kiến thức cốt lõi: Suất điện động lệch pha pi/2 so với từ thông; khi từ thông cực đại thì suất điện động bằng 0
\item (Sai) Kết quả không phụ thuộc dữ kiện và có thể bỏ qua điều kiện.
\item (Đúng) Suất điện động lệch pha pi/2 so với từ thông; khi từ thông cực đại thì suất điện động bằng 0.
\item (Sai) Từ thông và suất điện động luôn đạt cực đại đồng thời.
\end{itemchoice}
}
\end{ex}

% ===== Câu 231 =====
% ID: L12C3B17-53-TL
% Mức: VD
\dangbt{Thiết lập biểu thức suất điện động xoay chiều do máy phát tạo ra}
\begin{bt}
% ID: L12C3B17-53-TL
% Mức: VD
Nêu một tình huống thực tiễn liên quan đến Thiết lập biểu thức suất điện động xoay chiều do máy phát tạo ra và giải thích bằng kiến thức Vật lí.
\loigiai{
Cần nêu được: Nếu Phi = Phi0 cos(omega t + phi) thì e = omega Phi0 sin(omega t + phi). Khi vận dụng phải xác định đúng dữ kiện, phương chiều nếu có, điều kiện áp dụng, hệ đơn vị và kiểm tra tính hợp lí. Nhận định “Suất điện động luôn cùng pha và bằng đúng từ thông.” sai.
}
\end{bt}

% ===== Câu 232 =====
% ID: L12C3B17-53-TLN
% Mức: VD
\begin{ex}
% ID: L12C3B17-53-TLN
% Mức: VD
Trong bốn phát biểu về Thiết lập biểu thức suất điện động xoay chiều do máy phát tạo ra, có bao nhiêu phát biểu đúng? (Chỉ ghi một số nguyên.) (1) Nếu Phi = Phi0 cos(omega t + phi) thì e = omega Phi0 sin(omega t + phi). (2) Suất điện động luôn cùng pha và bằng đúng từ thông. (3) Cần kiểm tra điều kiện áp dụng. (4) Có thể bỏ qua đơn vị.
\shortans{2}
\loigiai{
Đối chiếu với kiến thức: Nếu Phi = Phi0 cos(omega t + phi) thì e = omega Phi0 sin(omega t + phi). Có 2 phát biểu đúng.
}
\end{ex}

% ===== Câu 233 =====
% ID: L12C3B17-53-TN
% Mức: NB
\dangbt{Xác định biên độ, tần số, chu kì và pha của suất điện động xoay chiều}
\begin{ex}
% ID: L12C3B17-53-TN
% Mức: NB
Tình huống 3: chọn kết luận chính xác về Xác định biên độ, tần số, chu kì và pha của suất điện động xoay chiều?
\choice
{Kết quả không phụ thuộc dữ kiện và điều kiện áp dụng.}
{\True Biên độ suất điện động là E0 = NBS omega và tần số f = omega/(2pi).}
{Biên độ suất điện động không phụ thuộc tốc độ quay.}
{Có thể bỏ qua phương, chiều và đơn vị trong mọi trường hợp.}
\loigiai{
Biên độ suất điện động là E0 = NBS omega và tần số f = omega/(2pi). Vì vậy phương án $B$ đúng; các phương án còn lại sai.
}
\end{ex}

% ===== Câu 234 =====
% ID: L12C3B17-54-DS
% Mức: VD
\dangbt{Khai thác đồ thị từ thông và suất điện động theo thời gian}
\begin{ex}
% ID: L12C3B17-54-DS
% Mức: VD
Xét Khai thác đồ thị từ thông và suất điện động theo thời gian (Tình huống 4). Hãy xác định tính đúng, sai của bốn phát biểu sau.
\choiceTF
{Kết quả không phụ thuộc dữ kiện và có thể bỏ qua điều kiện.}
{\True Khi giải cần xác định đúng dữ kiện, phương chiều, điều kiện áp dụng và đơn vị.}
{Từ thông và suất điện động luôn đạt cực đại đồng thời.}
{\True Suất điện động lệch pha pi/2 so với từ thông; khi từ thông cực đại thì suất điện động bằng 0.}
\loigiai{
\begin{itemchoice}
\item (Sai) Kết quả không phụ thuộc dữ kiện và có thể bỏ qua điều kiện. (Vì): Kết quả không phụ thuộc dữ kiện và có thể bỏ qua điều kiện. b) Đúng: Khi giải cần xác định đúng dữ kiện, phương chiều, điều kiện áp dụng và đơn vị. c) Sai: Từ thông và suất điện động luôn đạt cực đại đồng thời. d) Đúng: Suất điện động lệch pha pi/2 so với từ thông; khi từ thông cực đại thì suất điện động bằng 0. Kiến thức cốt lõi: Suất điện động lệch pha pi/2 so với từ thông; khi từ thông cực đại thì suất điện động bằng 0
\item (Đúng) Khi giải cần xác định đúng dữ kiện, phương chiều, điều kiện áp dụng và đơn vị.
\item (Sai) Từ thông và suất điện động luôn đạt cực đại đồng thời.
\item (Đúng) Suất điện động lệch pha pi/2 so với từ thông; khi từ thông cực đại thì suất điện động bằng 0.
\end{itemchoice}
}
\end{ex}

% ===== Câu 235 =====
% ID: L12C3B17-54-TL
% Mức: VDC
\dangbt{Thiết lập biểu thức suất điện động xoay chiều do máy phát tạo ra}
\begin{bt}
% ID: L12C3B17-54-TL
% Mức: VDC
So sánh nhận định đúng “Nếu Phi = Phi0 cos(omega t + phi) thì e = omega Phi0 sin(omega t + phi).” với nhận định sai “Suất điện động luôn cùng pha và bằng đúng từ thông.”; chỉ rõ căn cứ Vật lí.
\loigiai{
Cần nêu được: Nếu Phi = Phi0 cos(omega t + phi) thì e = omega Phi0 sin(omega t + phi). Khi vận dụng phải xác định đúng dữ kiện, phương chiều nếu có, điều kiện áp dụng, hệ đơn vị và kiểm tra tính hợp lí. Nhận định “Suất điện động luôn cùng pha và bằng đúng từ thông.” sai.
}
\end{bt}

% ===== Câu 236 =====
% ID: L12C3B17-54-TLN
% Mức: VDC
\begin{ex}
% ID: L12C3B17-54-TLN
% Mức: VDC
Trong bốn phát biểu về Thiết lập biểu thức suất điện động xoay chiều do máy phát tạo ra, có bao nhiêu phát biểu đúng? (Chỉ ghi một số nguyên.) (1) Khi giải cần xác định đúng dữ kiện, phương chiều, điều kiện áp dụng và đơn vị. (2) Suất điện động luôn cùng pha và bằng đúng từ thông. (3) Phải dùng đúng định luật cho tình huống. (4) Nếu Phi = Phi0 cos(omega t + phi) thì e = omega Phi0 sin(omega t + phi).
\shortans{3}
\loigiai{
Đối chiếu với kiến thức: Nếu Phi = Phi0 cos(omega t + phi) thì e = omega Phi0 sin(omega t + phi). Có 3 phát biểu đúng.
}
\end{ex}

% ===== Câu 237 =====
% ID: L12C3B17-54-TN
% Mức: TH
\dangbt{Xác định biên độ, tần số, chu kì và pha của suất điện động xoay chiều}
\begin{ex}
% ID: L12C3B17-54-TN
% Mức: TH
Nội dung nào mô tả đúng nhất Xác định biên độ, tần số, chu kì và pha của suất điện động xoay chiều?
\choice
{Có thể bỏ qua phương, chiều và đơn vị trong mọi trường hợp.}
{Kết quả không phụ thuộc dữ kiện và điều kiện áp dụng.}
{\True Biên độ suất điện động là E0 = NBS omega và tần số f = omega/(2pi).}
{Biên độ suất điện động không phụ thuộc tốc độ quay.}
\loigiai{
Biên độ suất điện động là E0 = NBS omega và tần số f = omega/(2pi). Vì vậy phương án $C$ đúng; các phương án còn lại sai.
}
\end{ex}

% ===== Câu 238 =====
% ID: L12C3B17-55-DS
% Mức: VD
\dangbt{Khai thác đồ thị từ thông và suất điện động theo thời gian}
\begin{ex}
% ID: L12C3B17-55-DS
% Mức: VD
Xét Khai thác đồ thị từ thông và suất điện động theo thời gian (Tình huống 5). Hãy xác định tính đúng, sai của bốn phát biểu sau.
\choiceTF
{\True Suất điện động lệch pha pi/2 so với từ thông; khi từ thông cực đại thì suất điện động bằng 0.}
{Kết quả không phụ thuộc dữ kiện và có thể bỏ qua điều kiện.}
{Từ thông và suất điện động luôn đạt cực đại đồng thời.}
{\True Khi giải cần xác định đúng dữ kiện, phương chiều, điều kiện áp dụng và đơn vị.}
\loigiai{
\begin{itemchoice}
\item (Đúng) Suất điện động lệch pha pi/2 so với từ thông; khi từ thông cực đại thì suất điện động bằng 0. (Vì): Suất điện động lệch pha pi/2 so với từ thông; khi từ thông cực đại thì suất điện động bằng 0. b) Sai: Kết quả không phụ thuộc dữ kiện và có thể bỏ qua điều kiện. c) Sai: Từ thông và suất điện động luôn đạt cực đại đồng thời. d) Đúng: Khi giải cần xác định đúng dữ kiện, phương chiều, điều kiện áp dụng và đơn vị. Kiến thức cốt lõi: Suất điện động lệch pha pi/2 so với từ thông; khi từ thông cực đại thì suất điện động bằng 0
\item (Sai) Kết quả không phụ thuộc dữ kiện và có thể bỏ qua điều kiện.
\item (Sai) Từ thông và suất điện động luôn đạt cực đại đồng thời.
\item (Đúng) Khi giải cần xác định đúng dữ kiện, phương chiều, điều kiện áp dụng và đơn vị.
\end{itemchoice}
}
\end{ex}

% ===== Câu 239 =====
% ID: L12C3B17-55-TL
% Mức: TH
\dangbt{Xác định biên độ, tần số, chu kì và pha của suất điện động xoay chiều}
\begin{bt}
% ID: L12C3B17-55-TL
% Mức: TH
Trình bày kiến thức cốt lõi của Xác định biên độ, tần số, chu kì và pha của suất điện động xoay chiều và nêu điều kiện áp dụng.
\loigiai{
Cần nêu được: Biên độ suất điện động là E0 = NBS omega và tần số f = omega/(2pi). Khi vận dụng phải xác định đúng dữ kiện, phương chiều nếu có, điều kiện áp dụng, hệ đơn vị và kiểm tra tính hợp lí. Nhận định “Biên độ suất điện động không phụ thuộc tốc độ quay.” sai.
}
\end{bt}

% ===== Câu 240 =====
% ID: L12C3B17-55-TLN
% Mức: TH
\begin{ex}
% ID: L12C3B17-55-TLN
% Mức: TH
Trong bốn phát biểu về Xác định biên độ, tần số, chu kì và pha của suất điện động xoay chiều, có bao nhiêu phát biểu đúng? (Chỉ ghi một số nguyên.) (1) Biên độ suất điện động là E0 = NBS omega và tần số f = omega/(2pi). (2) Biên độ suất điện động không phụ thuộc tốc độ quay. (3) Cần kiểm tra điều kiện áp dụng. (4) Có thể bỏ qua đơn vị.
\shortans{2}
\loigiai{
Đối chiếu với kiến thức: Biên độ suất điện động là E0 = NBS omega và tần số f = omega/(2pi). Có 2 phát biểu đúng.
}
\end{ex}

% ===== Câu 241 =====
% ID: L12C3B17-55-TN
% Mức: TH
\begin{ex}
% ID: L12C3B17-55-TN
% Mức: TH
Khi phân tích Xác định biên độ, tần số, chu kì và pha của suất điện động xoay chiều?
\choice
{Biên độ suất điện động không phụ thuộc tốc độ quay.}
{Có thể bỏ qua phương, chiều và đơn vị trong mọi trường hợp.}
{Kết quả không phụ thuộc dữ kiện và điều kiện áp dụng.}
{\True Biên độ suất điện động là E0 = NBS omega và tần số f = omega/(2pi).}
\loigiai{
Biên độ suất điện động là E0 = NBS omega và tần số f = omega/(2pi). Vì vậy phương án $D$ đúng; các phương án còn lại sai.
}
\end{ex}

% ===== Câu 242 =====
% ID: L12C3B17-56-DS
% Mức: TH
\dangbt{Bài toán thực tiễn về vận hành và hiệu suất máy phát điện xoay chiều}
\begin{ex}
% ID: L12C3B17-56-DS
% Mức: TH
Xét Bài toán thực tiễn về vận hành và hiệu suất máy phát điện xoay chiều (Tình huống 1). Hãy xác định tính đúng, sai của bốn phát biểu sau.
\choiceTF
{\True Hiệu suất máy phát được tính bằng công suất điện có ích chia cho công suất cơ đưa vào.}
{Hiệu suất máy phát bằng công suất hao phí chia công suất có ích.}
{\True Khi giải cần xác định đúng dữ kiện, phương chiều, điều kiện áp dụng và đơn vị.}
{Kết quả không phụ thuộc dữ kiện và có thể bỏ qua điều kiện.}
\loigiai{
\begin{itemchoice}
\item (Đúng) Hiệu suất máy phát được tính bằng công suất điện có ích chia cho công suất cơ đưa vào. (Vì): Hiệu suất máy phát được tính bằng công suất điện có ích chia cho công suất cơ đưa vào. b) Sai: Hiệu suất máy phát bằng công suất hao phí chia công suất có ích. c) Đúng: Khi giải cần xác định đúng dữ kiện, phương chiều, điều kiện áp dụng và đơn vị. d) Sai: Kết quả không phụ thuộc dữ kiện và có thể bỏ qua điều kiện. Kiến thức cốt lõi: Hiệu suất máy phát được tính bằng công suất điện có ích chia cho công suất cơ đưa vào
\item (Sai) Hiệu suất máy phát bằng công suất hao phí chia công suất có ích.
\item (Đúng) Khi giải cần xác định đúng dữ kiện, phương chiều, điều kiện áp dụng và đơn vị.
\item (Sai) Kết quả không phụ thuộc dữ kiện và có thể bỏ qua điều kiện.
\end{itemchoice}
}
\end{ex}

% ===== Câu 243 =====
% ID: L12C3B17-56-TL
% Mức: TH
\dangbt{Xác định biên độ, tần số, chu kì và pha của suất điện động xoay chiều}
\begin{bt}
% ID: L12C3B17-56-TL
% Mức: TH
Giải thích vì sao nhận định sau đúng: “Biên độ suất điện động là E0 = NBS omega và tần số f = omega/(2pi).”
\loigiai{
Cần nêu được: Biên độ suất điện động là E0 = NBS omega và tần số f = omega/(2pi). Khi vận dụng phải xác định đúng dữ kiện, phương chiều nếu có, điều kiện áp dụng, hệ đơn vị và kiểm tra tính hợp lí. Nhận định “Biên độ suất điện động không phụ thuộc tốc độ quay.” sai.
}
\end{bt}

% ===== Câu 244 =====
% ID: L12C3B17-56-TLN
% Mức: TH
\begin{ex}
% ID: L12C3B17-56-TLN
% Mức: TH
Trong bốn phát biểu về Xác định biên độ, tần số, chu kì và pha của suất điện động xoay chiều, có bao nhiêu phát biểu đúng? (Chỉ ghi một số nguyên.) (1) Biên độ suất điện động không phụ thuộc tốc độ quay. (2) Biên độ suất điện động là E0 = NBS omega và tần số f = omega/(2pi). (3) Kết quả phải phù hợp ý nghĩa vật lí. (4) Mọi đại lượng đều là vô hướng.
\shortans{1}
\loigiai{
Đối chiếu với kiến thức: Biên độ suất điện động là E0 = NBS omega và tần số f = omega/(2pi). Có 1 phát biểu đúng.
}
\end{ex}

% ===== Câu 245 =====
% ID: L12C3B17-56-TN
% Mức: TH
\begin{ex}
% ID: L12C3B17-56-TN
% Mức: TH
Một học sinh ôn tập Xác định biên độ, tần số, chu kì và pha của suất điện động xoay chiều?
\choice
{\True Biên độ suất điện động là E0 = NBS omega và tần số f = omega/(2pi).}
{Biên độ suất điện động không phụ thuộc tốc độ quay.}
{Có thể bỏ qua phương, chiều và đơn vị trong mọi trường hợp.}
{Kết quả không phụ thuộc dữ kiện và điều kiện áp dụng.}
\loigiai{
Biên độ suất điện động là E0 = NBS omega và tần số f = omega/(2pi). Vì vậy phương án $A$ đúng; các phương án còn lại sai.
}
\end{ex}

% ===== Câu 246 =====
% ID: L12C3B17-57-DS
% Mức: TH
\dangbt{Bài toán thực tiễn về vận hành và hiệu suất máy phát điện xoay chiều}
\begin{ex}
% ID: L12C3B17-57-DS
% Mức: TH
Xét Bài toán thực tiễn về vận hành và hiệu suất máy phát điện xoay chiều (Tình huống 2). Hãy xác định tính đúng, sai của bốn phát biểu sau.
\choiceTF
{Hiệu suất máy phát bằng công suất hao phí chia công suất có ích.}
{\True Hiệu suất máy phát được tính bằng công suất điện có ích chia cho công suất cơ đưa vào.}
{Kết quả không phụ thuộc dữ kiện và có thể bỏ qua điều kiện.}
{\True Khi giải cần xác định đúng dữ kiện, phương chiều, điều kiện áp dụng và đơn vị.}
\loigiai{
\begin{itemchoice}
\item (Sai) Hiệu suất máy phát bằng công suất hao phí chia công suất có ích. (Vì): Hiệu suất máy phát bằng công suất hao phí chia công suất có ích. b) Đúng: Hiệu suất máy phát được tính bằng công suất điện có ích chia cho công suất cơ đưa vào. c) Sai: Kết quả không phụ thuộc dữ kiện và có thể bỏ qua điều kiện. d) Đúng: Khi giải cần xác định đúng dữ kiện, phương chiều, điều kiện áp dụng và đơn vị. Kiến thức cốt lõi: Hiệu suất máy phát được tính bằng công suất điện có ích chia cho công suất cơ đưa vào
\item (Đúng) Hiệu suất máy phát được tính bằng công suất điện có ích chia cho công suất cơ đưa vào.
\item (Sai) Kết quả không phụ thuộc dữ kiện và có thể bỏ qua điều kiện.
\item (Đúng) Khi giải cần xác định đúng dữ kiện, phương chiều, điều kiện áp dụng và đơn vị.
\end{itemchoice}
}
\end{ex}

% ===== Câu 247 =====
% ID: L12C3B17-57-TL
% Mức: VD
\dangbt{Xác định biên độ, tần số, chu kì và pha của suất điện động xoay chiều}
\begin{bt}
% ID: L12C3B17-57-TL
% Mức: VD
Phân tích sai lầm trong nhận định: “Biên độ suất điện động không phụ thuộc tốc độ quay.”
\loigiai{
Cần nêu được: Biên độ suất điện động là E0 = NBS omega và tần số f = omega/(2pi). Khi vận dụng phải xác định đúng dữ kiện, phương chiều nếu có, điều kiện áp dụng, hệ đơn vị và kiểm tra tính hợp lí. Nhận định “Biên độ suất điện động không phụ thuộc tốc độ quay.” sai.
}
\end{bt}

% ===== Câu 248 =====
% ID: L12C3B17-57-TLN
% Mức: TH
\begin{ex}
% ID: L12C3B17-57-TLN
% Mức: TH
Trong bốn phát biểu về Xác định biên độ, tần số, chu kì và pha của suất điện động xoay chiều, có bao nhiêu phát biểu đúng? (Chỉ ghi một số nguyên.) (1) Biên độ suất điện động là E0 = NBS omega và tần số f = omega/(2pi). (2) Biên độ suất điện động không phụ thuộc tốc độ quay. (3) Cần kiểm tra điều kiện áp dụng. (4) Có thể bỏ qua đơn vị.
\shortans{2}
\loigiai{
Đối chiếu với kiến thức: Biên độ suất điện động là E0 = NBS omega và tần số f = omega/(2pi). Có 2 phát biểu đúng.
}
\end{ex}

% ===== Câu 249 =====
% ID: L12C3B17-57-TN
% Mức: TH
\begin{ex}
% ID: L12C3B17-57-TN
% Mức: TH
Chọn phát biểu không trái với kiến thức về Xác định biên độ, tần số, chu kì và pha của suất điện động xoay chiều?
\choice
{Kết quả không phụ thuộc dữ kiện và điều kiện áp dụng.}
{\True Biên độ suất điện động là E0 = NBS omega và tần số f = omega/(2pi).}
{Biên độ suất điện động không phụ thuộc tốc độ quay.}
{Có thể bỏ qua phương, chiều và đơn vị trong mọi trường hợp.}
\loigiai{
Biên độ suất điện động là E0 = NBS omega và tần số f = omega/(2pi). Vì vậy phương án $B$ đúng; các phương án còn lại sai.
}
\end{ex}

% ===== Câu 250 =====
% ID: L12C3B17-58-DS
% Mức: VD
\dangbt{Bài toán thực tiễn về vận hành và hiệu suất máy phát điện xoay chiều}
\begin{ex}
% ID: L12C3B17-58-DS
% Mức: VD
Xét Bài toán thực tiễn về vận hành và hiệu suất máy phát điện xoay chiều (Tình huống 3). Hãy xác định tính đúng, sai của bốn phát biểu sau.
\choiceTF
{\True Khi giải cần xác định đúng dữ kiện, phương chiều, điều kiện áp dụng và đơn vị.}
{Kết quả không phụ thuộc dữ kiện và có thể bỏ qua điều kiện.}
{\True Hiệu suất máy phát được tính bằng công suất điện có ích chia cho công suất cơ đưa vào.}
{Hiệu suất máy phát bằng công suất hao phí chia công suất có ích.}
\loigiai{
\begin{itemchoice}
\item (Đúng) Khi giải cần xác định đúng dữ kiện, phương chiều, điều kiện áp dụng và đơn vị. (Vì): Khi giải cần xác định đúng dữ kiện, phương chiều, điều kiện áp dụng và đơn vị. b) Sai: Kết quả không phụ thuộc dữ kiện và có thể bỏ qua điều kiện. c) Đúng: Hiệu suất máy phát được tính bằng công suất điện có ích chia cho công suất cơ đưa vào. d) Sai: Hiệu suất máy phát bằng công suất hao phí chia công suất có ích. Kiến thức cốt lõi: Hiệu suất máy phát được tính bằng công suất điện có ích chia cho công suất cơ đưa vào
\item (Sai) Kết quả không phụ thuộc dữ kiện và có thể bỏ qua điều kiện.
\item (Đúng) Hiệu suất máy phát được tính bằng công suất điện có ích chia cho công suất cơ đưa vào.
\item (Sai) Hiệu suất máy phát bằng công suất hao phí chia công suất có ích.
\end{itemchoice}
}
\end{ex}

% ===== Câu 251 =====
% ID: L12C3B17-58-TL
% Mức: VD
\dangbt{Xác định biên độ, tần số, chu kì và pha của suất điện động xoay chiều}
\begin{bt}
% ID: L12C3B17-58-TL
% Mức: VD
Đề xuất các bước giải một bài toán thuộc dạng Xác định biên độ, tần số, chu kì và pha của suất điện động xoay chiều.
\loigiai{
Cần nêu được: Biên độ suất điện động là E0 = NBS omega và tần số f = omega/(2pi). Khi vận dụng phải xác định đúng dữ kiện, phương chiều nếu có, điều kiện áp dụng, hệ đơn vị và kiểm tra tính hợp lí. Nhận định “Biên độ suất điện động không phụ thuộc tốc độ quay.” sai.
}
\end{bt}

% ===== Câu 252 =====
% ID: L12C3B17-58-TLN
% Mức: VD
\begin{ex}
% ID: L12C3B17-58-TLN
% Mức: VD
Trong bốn phát biểu về Xác định biên độ, tần số, chu kì và pha của suất điện động xoay chiều, có bao nhiêu phát biểu đúng? (Chỉ ghi một số nguyên.) (1) Biên độ suất điện động không phụ thuộc tốc độ quay. (2) Biên độ suất điện động là E0 = NBS omega và tần số f = omega/(2pi). (3) Kết quả phải phù hợp ý nghĩa vật lí. (4) Mọi đại lượng đều là vô hướng.
\shortans{2}
\loigiai{
Đối chiếu với kiến thức: Biên độ suất điện động là E0 = NBS omega và tần số f = omega/(2pi). Có 2 phát biểu đúng.
}
\end{ex}

% ===== Câu 253 =====
% ID: L12C3B17-58-TN
% Mức: VD
\begin{ex}
% ID: L12C3B17-58-TN
% Mức: VD
Cơ sở Vật lí đúng dùng để giải Xác định biên độ, tần số, chu kì và pha của suất điện động xoay chiều?
\choice
{Có thể bỏ qua phương, chiều và đơn vị trong mọi trường hợp.}
{Kết quả không phụ thuộc dữ kiện và điều kiện áp dụng.}
{\True Biên độ suất điện động là E0 = NBS omega và tần số f = omega/(2pi).}
{Biên độ suất điện động không phụ thuộc tốc độ quay.}
\loigiai{
Biên độ suất điện động là E0 = NBS omega và tần số f = omega/(2pi). Vì vậy phương án $C$ đúng; các phương án còn lại sai.
}
\end{ex}

% ===== Câu 254 =====
% ID: L12C3B17-59-DS
% Mức: VD
\dangbt{Bài toán thực tiễn về vận hành và hiệu suất máy phát điện xoay chiều}
\begin{ex}
% ID: L12C3B17-59-DS
% Mức: VD
Xét Bài toán thực tiễn về vận hành và hiệu suất máy phát điện xoay chiều (Tình huống 4). Hãy xác định tính đúng, sai của bốn phát biểu sau.
\choiceTF
{Kết quả không phụ thuộc dữ kiện và có thể bỏ qua điều kiện.}
{\True Khi giải cần xác định đúng dữ kiện, phương chiều, điều kiện áp dụng và đơn vị.}
{Hiệu suất máy phát bằng công suất hao phí chia công suất có ích.}
{\True Hiệu suất máy phát được tính bằng công suất điện có ích chia cho công suất cơ đưa vào.}
\loigiai{
\begin{itemchoice}
\item (Sai) Kết quả không phụ thuộc dữ kiện và có thể bỏ qua điều kiện. (Vì): Kết quả không phụ thuộc dữ kiện và có thể bỏ qua điều kiện. b) Đúng: Khi giải cần xác định đúng dữ kiện, phương chiều, điều kiện áp dụng và đơn vị. c) Sai: Hiệu suất máy phát bằng công suất hao phí chia công suất có ích. d) Đúng: Hiệu suất máy phát được tính bằng công suất điện có ích chia cho công suất cơ đưa vào. Kiến thức cốt lõi: Hiệu suất máy phát được tính bằng công suất điện có ích chia cho công suất cơ đưa vào
\item (Đúng) Khi giải cần xác định đúng dữ kiện, phương chiều, điều kiện áp dụng và đơn vị.
\item (Sai) Hiệu suất máy phát bằng công suất hao phí chia công suất có ích.
\item (Đúng) Hiệu suất máy phát được tính bằng công suất điện có ích chia cho công suất cơ đưa vào.
\end{itemchoice}
}
\end{ex}

% ===== Câu 255 =====
% ID: L12C3B17-59-TL
% Mức: VD
\dangbt{Xác định biên độ, tần số, chu kì và pha của suất điện động xoay chiều}
\begin{bt}
% ID: L12C3B17-59-TL
% Mức: VD
Nêu một tình huống thực tiễn liên quan đến Xác định biên độ, tần số, chu kì và pha của suất điện động xoay chiều và giải thích bằng kiến thức Vật lí.
\loigiai{
Cần nêu được: Biên độ suất điện động là E0 = NBS omega và tần số f = omega/(2pi). Khi vận dụng phải xác định đúng dữ kiện, phương chiều nếu có, điều kiện áp dụng, hệ đơn vị và kiểm tra tính hợp lí. Nhận định “Biên độ suất điện động không phụ thuộc tốc độ quay.” sai.
}
\end{bt}

% ===== Câu 256 =====
% ID: L12C3B17-59-TLN
% Mức: VD
\begin{ex}
% ID: L12C3B17-59-TLN
% Mức: VD
Trong bốn phát biểu về Xác định biên độ, tần số, chu kì và pha của suất điện động xoay chiều, có bao nhiêu phát biểu đúng? (Chỉ ghi một số nguyên.) (1) Biên độ suất điện động là E0 = NBS omega và tần số f = omega/(2pi). (2) Biên độ suất điện động không phụ thuộc tốc độ quay. (3) Cần kiểm tra điều kiện áp dụng. (4) Có thể bỏ qua đơn vị.
\shortans{2}
\loigiai{
Đối chiếu với kiến thức: Biên độ suất điện động là E0 = NBS omega và tần số f = omega/(2pi). Có 2 phát biểu đúng.
}
\end{ex}

% ===== Câu 257 =====
% ID: L12C3B17-59-TN
% Mức: VD
\begin{ex}
% ID: L12C3B17-59-TN
% Mức: VD
Trong một bài thực hành về Xác định biên độ, tần số, chu kì và pha của suất điện động xoay chiều?
\choice
{Biên độ suất điện động không phụ thuộc tốc độ quay.}
{Có thể bỏ qua phương, chiều và đơn vị trong mọi trường hợp.}
{Kết quả không phụ thuộc dữ kiện và điều kiện áp dụng.}
{\True Biên độ suất điện động là E0 = NBS omega và tần số f = omega/(2pi).}
\loigiai{
Biên độ suất điện động là E0 = NBS omega và tần số f = omega/(2pi). Vì vậy phương án $D$ đúng; các phương án còn lại sai.
}
\end{ex}

% ===== Câu 258 =====
% ID: L12C3B17-60-DS
% Mức: VD
\dangbt{Bài toán thực tiễn về vận hành và hiệu suất máy phát điện xoay chiều}
\begin{ex}
% ID: L12C3B17-60-DS
% Mức: VD
Xét Bài toán thực tiễn về vận hành và hiệu suất máy phát điện xoay chiều (Tình huống 5). Hãy xác định tính đúng, sai của bốn phát biểu sau.
\choiceTF
{\True Hiệu suất máy phát được tính bằng công suất điện có ích chia cho công suất cơ đưa vào.}
{Kết quả không phụ thuộc dữ kiện và có thể bỏ qua điều kiện.}
{Hiệu suất máy phát bằng công suất hao phí chia công suất có ích.}
{\True Khi giải cần xác định đúng dữ kiện, phương chiều, điều kiện áp dụng và đơn vị.}
\loigiai{
\begin{itemchoice}
\item (Đúng) Hiệu suất máy phát được tính bằng công suất điện có ích chia cho công suất cơ đưa vào. (Vì): Hiệu suất máy phát được tính bằng công suất điện có ích chia cho công suất cơ đưa vào. b) Sai: Kết quả không phụ thuộc dữ kiện và có thể bỏ qua điều kiện. c) Sai: Hiệu suất máy phát bằng công suất hao phí chia công suất có ích. d) Đúng: Khi giải cần xác định đúng dữ kiện, phương chiều, điều kiện áp dụng và đơn vị. Kiến thức cốt lõi: Hiệu suất máy phát được tính bằng công suất điện có ích chia cho công suất cơ đưa vào
\item (Sai) Kết quả không phụ thuộc dữ kiện và có thể bỏ qua điều kiện.
\item (Sai) Hiệu suất máy phát bằng công suất hao phí chia công suất có ích.
\item (Đúng) Khi giải cần xác định đúng dữ kiện, phương chiều, điều kiện áp dụng và đơn vị.
\end{itemchoice}
}
\end{ex}

% ===== Câu 259 =====
% ID: L12C3B17-60-TL
% Mức: VDC
\dangbt{Xác định biên độ, tần số, chu kì và pha của suất điện động xoay chiều}
\begin{bt}
% ID: L12C3B17-60-TL
% Mức: VDC
So sánh nhận định đúng “Biên độ suất điện động là E0 = NBS omega và tần số f = omega/(2pi).” với nhận định sai “Biên độ suất điện động không phụ thuộc tốc độ quay.”; chỉ rõ căn cứ Vật lí.
\loigiai{
Cần nêu được: Biên độ suất điện động là E0 = NBS omega và tần số f = omega/(2pi). Khi vận dụng phải xác định đúng dữ kiện, phương chiều nếu có, điều kiện áp dụng, hệ đơn vị và kiểm tra tính hợp lí. Nhận định “Biên độ suất điện động không phụ thuộc tốc độ quay.” sai.
}
\end{bt}

% ===== Câu 260 =====
% ID: L12C3B17-60-TLN
% Mức: VDC
\begin{ex}
% ID: L12C3B17-60-TLN
% Mức: VDC
Trong bốn phát biểu về Xác định biên độ, tần số, chu kì và pha của suất điện động xoay chiều, có bao nhiêu phát biểu đúng? (Chỉ ghi một số nguyên.) (1) Khi giải cần xác định đúng dữ kiện, phương chiều, điều kiện áp dụng và đơn vị. (2) Biên độ suất điện động không phụ thuộc tốc độ quay. (3) Phải dùng đúng định luật cho tình huống. (4) Biên độ suất điện động là E0 = NBS omega và tần số f = omega/(2pi).
\shortans{3}
\loigiai{
Đối chiếu với kiến thức: Biên độ suất điện động là E0 = NBS omega và tần số f = omega/(2pi). Có 3 phát biểu đúng.
}
\end{ex}

% ===== Câu 261 =====
% ID: L12C3B17-60-TN
% Mức: VD
\begin{ex}
% ID: L12C3B17-60-TN
% Mức: VD
Kết luận đúng khi vận dụng Xác định biên độ, tần số, chu kì và pha của suất điện động xoay chiều?
\choice
{\True Biên độ suất điện động là E0 = NBS omega và tần số f = omega/(2pi).}
{Biên độ suất điện động không phụ thuộc tốc độ quay.}
{Có thể bỏ qua phương, chiều và đơn vị trong mọi trường hợp.}
{Kết quả không phụ thuộc dữ kiện và điều kiện áp dụng.}
\loigiai{
Biên độ suất điện động là E0 = NBS omega và tần số f = omega/(2pi). Vì vậy phương án $A$ đúng; các phương án còn lại sai.
}
\end{ex}

% ===== Câu 262 =====
% ID: L12C3B17-61-TL
% Mức: TH
\dangbt{Khai thác đồ thị từ thông và suất điện động theo thời gian}
\begin{bt}
% ID: L12C3B17-61-TL
% Mức: TH
Trình bày kiến thức cốt lõi của Khai thác đồ thị từ thông và suất điện động theo thời gian và nêu điều kiện áp dụng.
\loigiai{
Cần nêu được: Suất điện động lệch pha pi/2 so với từ thông; khi từ thông cực đại thì suất điện động bằng 0. Khi vận dụng phải xác định đúng dữ kiện, phương chiều nếu có, điều kiện áp dụng, hệ đơn vị và kiểm tra tính hợp lí. Nhận định “Từ thông và suất điện động luôn đạt cực đại đồng thời.” sai.
}
\end{bt}

% ===== Câu 263 =====
% ID: L12C3B17-61-TLN
% Mức: TH
\begin{ex}
% ID: L12C3B17-61-TLN
% Mức: TH
Trong bốn phát biểu về Khai thác đồ thị từ thông và suất điện động theo thời gian, có bao nhiêu phát biểu đúng? (Chỉ ghi một số nguyên.) (1) Suất điện động lệch pha pi/2 so với từ thông; khi từ thông cực đại thì suất điện động bằng 0. (2) Từ thông và suất điện động luôn đạt cực đại đồng thời. (3) Cần kiểm tra điều kiện áp dụng. (4) Có thể bỏ qua đơn vị.
\shortans{2}
\loigiai{
Đối chiếu với kiến thức: Suất điện động lệch pha pi/2 so với từ thông; khi từ thông cực đại thì suất điện động bằng 0. Có 2 phát biểu đúng.
}
\end{ex}

% ===== Câu 264 =====
% ID: L12C3B17-61-TN
% Mức: NB
\dangbt{Nhận biết cấu tạo và nguyên lí của máy phát điện xoay chiều}
\begin{ex}
% ID: L12C3B17-61-TN
% Mức: NB
Phát biểu nào sau đây đúng về Nhận biết cấu tạo và nguyên lí của máy phát điện xoay chiều?
\choice
{Máy phát điện xoay chiều biến điện năng thành cơ năng nhờ hiệu ứng nhiệt.}
{Có thể bỏ qua phương, chiều và đơn vị trong mọi trường hợp.}
{Kết quả không phụ thuộc dữ kiện và điều kiện áp dụng.}
{\True Máy phát điện xoay chiều hoạt động dựa trên hiện tượng cảm ứng điện từ khi từ thông qua cuộn dây biến thiên.}
\loigiai{
Máy phát điện xoay chiều hoạt động dựa trên hiện tượng cảm ứng điện từ khi từ thông qua cuộn dây biến thiên. Vì vậy phương án $D$ đúng; các phương án còn lại sai.
}
\end{ex}

% ===== Câu 265 =====
% ID: L12C3B17-62-TL
% Mức: TH
\dangbt{Khai thác đồ thị từ thông và suất điện động theo thời gian}
\begin{bt}
% ID: L12C3B17-62-TL
% Mức: TH
Giải thích vì sao nhận định sau đúng: “Suất điện động lệch pha pi/2 so với từ thông; khi từ thông cực đại thì suất điện động bằng 0.”
\loigiai{
Cần nêu được: Suất điện động lệch pha pi/2 so với từ thông; khi từ thông cực đại thì suất điện động bằng 0. Khi vận dụng phải xác định đúng dữ kiện, phương chiều nếu có, điều kiện áp dụng, hệ đơn vị và kiểm tra tính hợp lí. Nhận định “Từ thông và suất điện động luôn đạt cực đại đồng thời.” sai.
}
\end{bt}

% ===== Câu 266 =====
% ID: L12C3B17-62-TLN
% Mức: TH
\begin{ex}
% ID: L12C3B17-62-TLN
% Mức: TH
Trong bốn phát biểu về Khai thác đồ thị từ thông và suất điện động theo thời gian, có bao nhiêu phát biểu đúng? (Chỉ ghi một số nguyên.) (1) Từ thông và suất điện động luôn đạt cực đại đồng thời. (2) Suất điện động lệch pha pi/2 so với từ thông; khi từ thông cực đại thì suất điện động bằng 0. (3) Kết quả phải phù hợp ý nghĩa vật lí. (4) Mọi đại lượng đều là vô hướng.
\shortans{1}
\loigiai{
Đối chiếu với kiến thức: Suất điện động lệch pha pi/2 so với từ thông; khi từ thông cực đại thì suất điện động bằng 0. Có 1 phát biểu đúng.
}
\end{ex}

% ===== Câu 267 =====
% ID: L12C3B17-62-TN
% Mức: NB
\dangbt{Nhận biết cấu tạo và nguyên lí của máy phát điện xoay chiều}
\begin{ex}
% ID: L12C3B17-62-TN
% Mức: NB
Trong nội dung Nhận biết cấu tạo và nguyên lí của máy phát điện xoay chiều?
\choice
{\True Máy phát điện xoay chiều hoạt động dựa trên hiện tượng cảm ứng điện từ khi từ thông qua cuộn dây biến thiên.}
{Máy phát điện xoay chiều biến điện năng thành cơ năng nhờ hiệu ứng nhiệt.}
{Có thể bỏ qua phương, chiều và đơn vị trong mọi trường hợp.}
{Kết quả không phụ thuộc dữ kiện và điều kiện áp dụng.}
\loigiai{
Máy phát điện xoay chiều hoạt động dựa trên hiện tượng cảm ứng điện từ khi từ thông qua cuộn dây biến thiên. Vì vậy phương án $A$ đúng; các phương án còn lại sai.
}
\end{ex}

% ===== Câu 268 =====
% ID: L12C3B17-63-TL
% Mức: VD
\dangbt{Khai thác đồ thị từ thông và suất điện động theo thời gian}
\begin{bt}
% ID: L12C3B17-63-TL
% Mức: VD
Phân tích sai lầm trong nhận định: “Từ thông và suất điện động luôn đạt cực đại đồng thời.”
\loigiai{
Cần nêu được: Suất điện động lệch pha pi/2 so với từ thông; khi từ thông cực đại thì suất điện động bằng 0. Khi vận dụng phải xác định đúng dữ kiện, phương chiều nếu có, điều kiện áp dụng, hệ đơn vị và kiểm tra tính hợp lí. Nhận định “Từ thông và suất điện động luôn đạt cực đại đồng thời.” sai.
}
\end{bt}

% ===== Câu 269 =====
% ID: L12C3B17-63-TLN
% Mức: TH
\begin{ex}
% ID: L12C3B17-63-TLN
% Mức: TH
Trong bốn phát biểu về Khai thác đồ thị từ thông và suất điện động theo thời gian, có bao nhiêu phát biểu đúng? (Chỉ ghi một số nguyên.) (1) Suất điện động lệch pha pi/2 so với từ thông; khi từ thông cực đại thì suất điện động bằng 0. (2) Từ thông và suất điện động luôn đạt cực đại đồng thời. (3) Cần kiểm tra điều kiện áp dụng. (4) Có thể bỏ qua đơn vị.
\shortans{2}
\loigiai{
Đối chiếu với kiến thức: Suất điện động lệch pha pi/2 so với từ thông; khi từ thông cực đại thì suất điện động bằng 0. Có 2 phát biểu đúng.
}
\end{ex}

% ===== Câu 270 =====
% ID: L12C3B17-63-TN
% Mức: NB
\dangbt{Nhận biết cấu tạo và nguyên lí của máy phát điện xoay chiều}
\begin{ex}
% ID: L12C3B17-63-TN
% Mức: NB
Tình huống 3: chọn kết luận chính xác về Nhận biết cấu tạo và nguyên lí của máy phát điện xoay chiều?
\choice
{Kết quả không phụ thuộc dữ kiện và điều kiện áp dụng.}
{\True Máy phát điện xoay chiều hoạt động dựa trên hiện tượng cảm ứng điện từ khi từ thông qua cuộn dây biến thiên.}
{Máy phát điện xoay chiều biến điện năng thành cơ năng nhờ hiệu ứng nhiệt.}
{Có thể bỏ qua phương, chiều và đơn vị trong mọi trường hợp.}
\loigiai{
Máy phát điện xoay chiều hoạt động dựa trên hiện tượng cảm ứng điện từ khi từ thông qua cuộn dây biến thiên. Vì vậy phương án $B$ đúng; các phương án còn lại sai.
}
\end{ex}

% ===== Câu 271 =====
% ID: L12C3B17-64-TL
% Mức: VD
\dangbt{Khai thác đồ thị từ thông và suất điện động theo thời gian}
\begin{bt}
% ID: L12C3B17-64-TL
% Mức: VD
Đề xuất các bước giải một bài toán thuộc dạng Khai thác đồ thị từ thông và suất điện động theo thời gian.
\loigiai{
Cần nêu được: Suất điện động lệch pha pi/2 so với từ thông; khi từ thông cực đại thì suất điện động bằng 0. Khi vận dụng phải xác định đúng dữ kiện, phương chiều nếu có, điều kiện áp dụng, hệ đơn vị và kiểm tra tính hợp lí. Nhận định “Từ thông và suất điện động luôn đạt cực đại đồng thời.” sai.
}
\end{bt}

% ===== Câu 272 =====
% ID: L12C3B17-64-TLN
% Mức: VD
\begin{ex}
% ID: L12C3B17-64-TLN
% Mức: VD
Trong bốn phát biểu về Khai thác đồ thị từ thông và suất điện động theo thời gian, có bao nhiêu phát biểu đúng? (Chỉ ghi một số nguyên.) (1) Từ thông và suất điện động luôn đạt cực đại đồng thời. (2) Suất điện động lệch pha pi/2 so với từ thông; khi từ thông cực đại thì suất điện động bằng 0. (3) Kết quả phải phù hợp ý nghĩa vật lí. (4) Mọi đại lượng đều là vô hướng.
\shortans{2}
\loigiai{
Đối chiếu với kiến thức: Suất điện động lệch pha pi/2 so với từ thông; khi từ thông cực đại thì suất điện động bằng 0. Có 2 phát biểu đúng.
}
\end{ex}

% ===== Câu 273 =====
% ID: L12C3B17-64-TN
% Mức: TH
\dangbt{Nhận biết cấu tạo và nguyên lí của máy phát điện xoay chiều}
\begin{ex}
% ID: L12C3B17-64-TN
% Mức: TH
Nội dung nào mô tả đúng nhất Nhận biết cấu tạo và nguyên lí của máy phát điện xoay chiều?
\choice
{Có thể bỏ qua phương, chiều và đơn vị trong mọi trường hợp.}
{Kết quả không phụ thuộc dữ kiện và điều kiện áp dụng.}
{\True Máy phát điện xoay chiều hoạt động dựa trên hiện tượng cảm ứng điện từ khi từ thông qua cuộn dây biến thiên.}
{Máy phát điện xoay chiều biến điện năng thành cơ năng nhờ hiệu ứng nhiệt.}
\loigiai{
Máy phát điện xoay chiều hoạt động dựa trên hiện tượng cảm ứng điện từ khi từ thông qua cuộn dây biến thiên. Vì vậy phương án $C$ đúng; các phương án còn lại sai.
}
\end{ex}

% ===== Câu 274 =====
% ID: L12C3B17-65-TL
% Mức: VD
\dangbt{Khai thác đồ thị từ thông và suất điện động theo thời gian}
\begin{bt}
% ID: L12C3B17-65-TL
% Mức: VD
Nêu một tình huống thực tiễn liên quan đến Khai thác đồ thị từ thông và suất điện động theo thời gian và giải thích bằng kiến thức Vật lí.
\loigiai{
Cần nêu được: Suất điện động lệch pha pi/2 so với từ thông; khi từ thông cực đại thì suất điện động bằng 0. Khi vận dụng phải xác định đúng dữ kiện, phương chiều nếu có, điều kiện áp dụng, hệ đơn vị và kiểm tra tính hợp lí. Nhận định “Từ thông và suất điện động luôn đạt cực đại đồng thời.” sai.
}
\end{bt}

% ===== Câu 275 =====
% ID: L12C3B17-65-TLN
% Mức: VD
\begin{ex}
% ID: L12C3B17-65-TLN
% Mức: VD
Trong bốn phát biểu về Khai thác đồ thị từ thông và suất điện động theo thời gian, có bao nhiêu phát biểu đúng? (Chỉ ghi một số nguyên.) (1) Suất điện động lệch pha pi/2 so với từ thông; khi từ thông cực đại thì suất điện động bằng 0. (2) Từ thông và suất điện động luôn đạt cực đại đồng thời. (3) Cần kiểm tra điều kiện áp dụng. (4) Có thể bỏ qua đơn vị.
\shortans{2}
\loigiai{
Đối chiếu với kiến thức: Suất điện động lệch pha pi/2 so với từ thông; khi từ thông cực đại thì suất điện động bằng 0. Có 2 phát biểu đúng.
}
\end{ex}

% ===== Câu 276 =====
% ID: L12C3B17-65-TN
% Mức: TH
\dangbt{Nhận biết cấu tạo và nguyên lí của máy phát điện xoay chiều}
\begin{ex}
% ID: L12C3B17-65-TN
% Mức: TH
Khi phân tích Nhận biết cấu tạo và nguyên lí của máy phát điện xoay chiều?
\choice
{Máy phát điện xoay chiều biến điện năng thành cơ năng nhờ hiệu ứng nhiệt.}
{Có thể bỏ qua phương, chiều và đơn vị trong mọi trường hợp.}
{Kết quả không phụ thuộc dữ kiện và điều kiện áp dụng.}
{\True Máy phát điện xoay chiều hoạt động dựa trên hiện tượng cảm ứng điện từ khi từ thông qua cuộn dây biến thiên.}
\loigiai{
Máy phát điện xoay chiều hoạt động dựa trên hiện tượng cảm ứng điện từ khi từ thông qua cuộn dây biến thiên. Vì vậy phương án $D$ đúng; các phương án còn lại sai.
}
\end{ex}

% ===== Câu 277 =====
% ID: L12C3B17-66-TL
% Mức: VDC
\dangbt{Khai thác đồ thị từ thông và suất điện động theo thời gian}
\begin{bt}
% ID: L12C3B17-66-TL
% Mức: VDC
So sánh nhận định đúng “Suất điện động lệch pha pi/2 so với từ thông; khi từ thông cực đại thì suất điện động bằng 0.” với nhận định sai “Từ thông và suất điện động luôn đạt cực đại đồng thời.”; chỉ rõ căn cứ Vật lí.
\loigiai{
Cần nêu được: Suất điện động lệch pha pi/2 so với từ thông; khi từ thông cực đại thì suất điện động bằng 0. Khi vận dụng phải xác định đúng dữ kiện, phương chiều nếu có, điều kiện áp dụng, hệ đơn vị và kiểm tra tính hợp lí. Nhận định “Từ thông và suất điện động luôn đạt cực đại đồng thời.” sai.
}
\end{bt}

% ===== Câu 278 =====
% ID: L12C3B17-66-TLN
% Mức: VDC
\begin{ex}
% ID: L12C3B17-66-TLN
% Mức: VDC
Trong bốn phát biểu về Khai thác đồ thị từ thông và suất điện động theo thời gian, có bao nhiêu phát biểu đúng? (Chỉ ghi một số nguyên.) (1) Khi giải cần xác định đúng dữ kiện, phương chiều, điều kiện áp dụng và đơn vị. (2) Từ thông và suất điện động luôn đạt cực đại đồng thời. (3) Phải dùng đúng định luật cho tình huống. (4) Suất điện động lệch pha pi/2 so với từ thông; khi từ thông cực đại thì suất điện động bằng 0.
\shortans{3}
\loigiai{
Đối chiếu với kiến thức: Suất điện động lệch pha pi/2 so với từ thông; khi từ thông cực đại thì suất điện động bằng 0. Có 3 phát biểu đúng.
}
\end{ex}

% ===== Câu 279 =====
% ID: L12C3B17-66-TN
% Mức: TH
\dangbt{Nhận biết cấu tạo và nguyên lí của máy phát điện xoay chiều}
\begin{ex}
% ID: L12C3B17-66-TN
% Mức: TH
Một học sinh ôn tập Nhận biết cấu tạo và nguyên lí của máy phát điện xoay chiều?
\choice
{\True Máy phát điện xoay chiều hoạt động dựa trên hiện tượng cảm ứng điện từ khi từ thông qua cuộn dây biến thiên.}
{Máy phát điện xoay chiều biến điện năng thành cơ năng nhờ hiệu ứng nhiệt.}
{Có thể bỏ qua phương, chiều và đơn vị trong mọi trường hợp.}
{Kết quả không phụ thuộc dữ kiện và điều kiện áp dụng.}
\loigiai{
Máy phát điện xoay chiều hoạt động dựa trên hiện tượng cảm ứng điện từ khi từ thông qua cuộn dây biến thiên. Vì vậy phương án $A$ đúng; các phương án còn lại sai.
}
\end{ex}

% ===== Câu 280 =====
% ID: L12C3B17-67-TL
% Mức: TH
\dangbt{Bài toán thực tiễn về vận hành và hiệu suất máy phát điện xoay chiều}
\begin{bt}
% ID: L12C3B17-67-TL
% Mức: TH
Trình bày kiến thức cốt lõi của Bài toán thực tiễn về vận hành và hiệu suất máy phát điện xoay chiều và nêu điều kiện áp dụng.
\loigiai{
Cần nêu được: Hiệu suất máy phát được tính bằng công suất điện có ích chia cho công suất cơ đưa vào. Khi vận dụng phải xác định đúng dữ kiện, phương chiều nếu có, điều kiện áp dụng, hệ đơn vị và kiểm tra tính hợp lí. Nhận định “Hiệu suất máy phát bằng công suất hao phí chia công suất có ích.” sai.
}
\end{bt}

% ===== Câu 281 =====
% ID: L12C3B17-67-TLN
% Mức: TH
\begin{ex}
% ID: L12C3B17-67-TLN
% Mức: TH
Trong bốn phát biểu về Bài toán thực tiễn về vận hành và hiệu suất máy phát điện xoay chiều, có bao nhiêu phát biểu đúng? (Chỉ ghi một số nguyên.) (1) Hiệu suất máy phát được tính bằng công suất điện có ích chia cho công suất cơ đưa vào. (2) Hiệu suất máy phát bằng công suất hao phí chia công suất có ích. (3) Cần kiểm tra điều kiện áp dụng. (4) Có thể bỏ qua đơn vị.
\shortans{2}
\loigiai{
Đối chiếu với kiến thức: Hiệu suất máy phát được tính bằng công suất điện có ích chia cho công suất cơ đưa vào. Có 2 phát biểu đúng.
}
\end{ex}

% ===== Câu 282 =====
% ID: L12C3B17-67-TN
% Mức: TH
\dangbt{Nhận biết cấu tạo và nguyên lí của máy phát điện xoay chiều}
\begin{ex}
% ID: L12C3B17-67-TN
% Mức: TH
Chọn phát biểu không trái với kiến thức về Nhận biết cấu tạo và nguyên lí của máy phát điện xoay chiều?
\choice
{Kết quả không phụ thuộc dữ kiện và điều kiện áp dụng.}
{\True Máy phát điện xoay chiều hoạt động dựa trên hiện tượng cảm ứng điện từ khi từ thông qua cuộn dây biến thiên.}
{Máy phát điện xoay chiều biến điện năng thành cơ năng nhờ hiệu ứng nhiệt.}
{Có thể bỏ qua phương, chiều và đơn vị trong mọi trường hợp.}
\loigiai{
Máy phát điện xoay chiều hoạt động dựa trên hiện tượng cảm ứng điện từ khi từ thông qua cuộn dây biến thiên. Vì vậy phương án $B$ đúng; các phương án còn lại sai.
}
\end{ex}

% ===== Câu 283 =====
% ID: L12C3B17-68-TL
% Mức: TH
\dangbt{Bài toán thực tiễn về vận hành và hiệu suất máy phát điện xoay chiều}
\begin{bt}
% ID: L12C3B17-68-TL
% Mức: TH
Giải thích vì sao nhận định sau đúng: “Hiệu suất máy phát được tính bằng công suất điện có ích chia cho công suất cơ đưa vào.”
\loigiai{
Cần nêu được: Hiệu suất máy phát được tính bằng công suất điện có ích chia cho công suất cơ đưa vào. Khi vận dụng phải xác định đúng dữ kiện, phương chiều nếu có, điều kiện áp dụng, hệ đơn vị và kiểm tra tính hợp lí. Nhận định “Hiệu suất máy phát bằng công suất hao phí chia công suất có ích.” sai.
}
\end{bt}

% ===== Câu 284 =====
% ID: L12C3B17-68-TLN
% Mức: TH
\begin{ex}
% ID: L12C3B17-68-TLN
% Mức: TH
Trong bốn phát biểu về Bài toán thực tiễn về vận hành và hiệu suất máy phát điện xoay chiều, có bao nhiêu phát biểu đúng? (Chỉ ghi một số nguyên.) (1) Hiệu suất máy phát bằng công suất hao phí chia công suất có ích. (2) Hiệu suất máy phát được tính bằng công suất điện có ích chia cho công suất cơ đưa vào. (3) Kết quả phải phù hợp ý nghĩa vật lí. (4) Mọi đại lượng đều là vô hướng.
\shortans{1}
\loigiai{
Đối chiếu với kiến thức: Hiệu suất máy phát được tính bằng công suất điện có ích chia cho công suất cơ đưa vào. Có 1 phát biểu đúng.
}
\end{ex}

% ===== Câu 285 =====
% ID: L12C3B17-68-TN
% Mức: VD
\dangbt{Nhận biết cấu tạo và nguyên lí của máy phát điện xoay chiều}
\begin{ex}
% ID: L12C3B17-68-TN
% Mức: VD
Cơ sở Vật lí đúng dùng để giải Nhận biết cấu tạo và nguyên lí của máy phát điện xoay chiều?
\choice
{Có thể bỏ qua phương, chiều và đơn vị trong mọi trường hợp.}
{Kết quả không phụ thuộc dữ kiện và điều kiện áp dụng.}
{\True Máy phát điện xoay chiều hoạt động dựa trên hiện tượng cảm ứng điện từ khi từ thông qua cuộn dây biến thiên.}
{Máy phát điện xoay chiều biến điện năng thành cơ năng nhờ hiệu ứng nhiệt.}
\loigiai{
Máy phát điện xoay chiều hoạt động dựa trên hiện tượng cảm ứng điện từ khi từ thông qua cuộn dây biến thiên. Vì vậy phương án $C$ đúng; các phương án còn lại sai.
}
\end{ex}

% ===== Câu 286 =====
% ID: L12C3B17-69-TL
% Mức: VD
\dangbt{Bài toán thực tiễn về vận hành và hiệu suất máy phát điện xoay chiều}
\begin{bt}
% ID: L12C3B17-69-TL
% Mức: VD
Phân tích sai lầm trong nhận định: “Hiệu suất máy phát bằng công suất hao phí chia công suất có ích.”
\loigiai{
Cần nêu được: Hiệu suất máy phát được tính bằng công suất điện có ích chia cho công suất cơ đưa vào. Khi vận dụng phải xác định đúng dữ kiện, phương chiều nếu có, điều kiện áp dụng, hệ đơn vị và kiểm tra tính hợp lí. Nhận định “Hiệu suất máy phát bằng công suất hao phí chia công suất có ích.” sai.
}
\end{bt}

% ===== Câu 287 =====
% ID: L12C3B17-69-TLN
% Mức: TH
\begin{ex}
% ID: L12C3B17-69-TLN
% Mức: TH
Trong bốn phát biểu về Bài toán thực tiễn về vận hành và hiệu suất máy phát điện xoay chiều, có bao nhiêu phát biểu đúng? (Chỉ ghi một số nguyên.) (1) Hiệu suất máy phát được tính bằng công suất điện có ích chia cho công suất cơ đưa vào. (2) Hiệu suất máy phát bằng công suất hao phí chia công suất có ích. (3) Cần kiểm tra điều kiện áp dụng. (4) Có thể bỏ qua đơn vị.
\shortans{2}
\loigiai{
Đối chiếu với kiến thức: Hiệu suất máy phát được tính bằng công suất điện có ích chia cho công suất cơ đưa vào. Có 2 phát biểu đúng.
}
\end{ex}

% ===== Câu 288 =====
% ID: L12C3B17-69-TN
% Mức: VD
\dangbt{Nhận biết cấu tạo và nguyên lí của máy phát điện xoay chiều}
\begin{ex}
% ID: L12C3B17-69-TN
% Mức: VD
Trong một bài thực hành về Nhận biết cấu tạo và nguyên lí của máy phát điện xoay chiều?
\choice
{Máy phát điện xoay chiều biến điện năng thành cơ năng nhờ hiệu ứng nhiệt.}
{Có thể bỏ qua phương, chiều và đơn vị trong mọi trường hợp.}
{Kết quả không phụ thuộc dữ kiện và điều kiện áp dụng.}
{\True Máy phát điện xoay chiều hoạt động dựa trên hiện tượng cảm ứng điện từ khi từ thông qua cuộn dây biến thiên.}
\loigiai{
Máy phát điện xoay chiều hoạt động dựa trên hiện tượng cảm ứng điện từ khi từ thông qua cuộn dây biến thiên. Vì vậy phương án $D$ đúng; các phương án còn lại sai.
}
\end{ex}

% ===== Câu 289 =====
% ID: L12C3B17-70-TL
% Mức: VD
\dangbt{Bài toán thực tiễn về vận hành và hiệu suất máy phát điện xoay chiều}
\begin{bt}
% ID: L12C3B17-70-TL
% Mức: VD
Đề xuất các bước giải một bài toán thuộc dạng Bài toán thực tiễn về vận hành và hiệu suất máy phát điện xoay chiều.
\loigiai{
Cần nêu được: Hiệu suất máy phát được tính bằng công suất điện có ích chia cho công suất cơ đưa vào. Khi vận dụng phải xác định đúng dữ kiện, phương chiều nếu có, điều kiện áp dụng, hệ đơn vị và kiểm tra tính hợp lí. Nhận định “Hiệu suất máy phát bằng công suất hao phí chia công suất có ích.” sai.
}
\end{bt}

% ===== Câu 290 =====
% ID: L12C3B17-70-TLN
% Mức: VD
\begin{ex}
% ID: L12C3B17-70-TLN
% Mức: VD
Trong bốn phát biểu về Bài toán thực tiễn về vận hành và hiệu suất máy phát điện xoay chiều, có bao nhiêu phát biểu đúng? (Chỉ ghi một số nguyên.) (1) Hiệu suất máy phát bằng công suất hao phí chia công suất có ích. (2) Hiệu suất máy phát được tính bằng công suất điện có ích chia cho công suất cơ đưa vào. (3) Kết quả phải phù hợp ý nghĩa vật lí. (4) Mọi đại lượng đều là vô hướng.
\shortans{2}
\loigiai{
Đối chiếu với kiến thức: Hiệu suất máy phát được tính bằng công suất điện có ích chia cho công suất cơ đưa vào. Có 2 phát biểu đúng.
}
\end{ex}

% ===== Câu 291 =====
% ID: L12C3B17-70-TN
% Mức: VD
\dangbt{Nhận biết cấu tạo và nguyên lí của máy phát điện xoay chiều}
\begin{ex}
% ID: L12C3B17-70-TN
% Mức: VD
Kết luận đúng khi vận dụng Nhận biết cấu tạo và nguyên lí của máy phát điện xoay chiều?
\choice
{\True Máy phát điện xoay chiều hoạt động dựa trên hiện tượng cảm ứng điện từ khi từ thông qua cuộn dây biến thiên.}
{Máy phát điện xoay chiều biến điện năng thành cơ năng nhờ hiệu ứng nhiệt.}
{Có thể bỏ qua phương, chiều và đơn vị trong mọi trường hợp.}
{Kết quả không phụ thuộc dữ kiện và điều kiện áp dụng.}
\loigiai{
Máy phát điện xoay chiều hoạt động dựa trên hiện tượng cảm ứng điện từ khi từ thông qua cuộn dây biến thiên. Vì vậy phương án $A$ đúng; các phương án còn lại sai.
}
\end{ex}

% ===== Câu 292 =====
% ID: L12C3B17-71-TL
% Mức: VD
\dangbt{Bài toán thực tiễn về vận hành và hiệu suất máy phát điện xoay chiều}
\begin{bt}
% ID: L12C3B17-71-TL
% Mức: VD
Nêu một tình huống thực tiễn liên quan đến Bài toán thực tiễn về vận hành và hiệu suất máy phát điện xoay chiều và giải thích bằng kiến thức Vật lí.
\loigiai{
Cần nêu được: Hiệu suất máy phát được tính bằng công suất điện có ích chia cho công suất cơ đưa vào. Khi vận dụng phải xác định đúng dữ kiện, phương chiều nếu có, điều kiện áp dụng, hệ đơn vị và kiểm tra tính hợp lí. Nhận định “Hiệu suất máy phát bằng công suất hao phí chia công suất có ích.” sai.
}
\end{bt}

% ===== Câu 293 =====
% ID: L12C3B17-71-TLN
% Mức: VD
\begin{ex}
% ID: L12C3B17-71-TLN
% Mức: VD
Trong bốn phát biểu về Bài toán thực tiễn về vận hành và hiệu suất máy phát điện xoay chiều, có bao nhiêu phát biểu đúng? (Chỉ ghi một số nguyên.) (1) Hiệu suất máy phát được tính bằng công suất điện có ích chia cho công suất cơ đưa vào. (2) Hiệu suất máy phát bằng công suất hao phí chia công suất có ích. (3) Cần kiểm tra điều kiện áp dụng. (4) Có thể bỏ qua đơn vị.
\shortans{2}
\loigiai{
Đối chiếu với kiến thức: Hiệu suất máy phát được tính bằng công suất điện có ích chia cho công suất cơ đưa vào. Có 2 phát biểu đúng.
}
\end{ex}

% ===== Câu 294 =====
% ID: L12C3B17-71-TN
% Mức: NB
\dangbt{Tính từ thông qua khung dây quay trong từ trường}
\begin{ex}
% ID: L12C3B17-71-TN
% Mức: NB
Phát biểu nào sau đây đúng về Tính từ thông qua khung dây quay trong từ trường?
\choice
{Từ thông qua khung quay luôn không đổi và bằng NBS.}
{Có thể bỏ qua phương, chiều và đơn vị trong mọi trường hợp.}
{Kết quả không phụ thuộc dữ kiện và điều kiện áp dụng.}
{\True Khung $N$ vòng quay đều có từ thông Phi = NBS cos(omega t + phi).}
\loigiai{
Khung $N$ vòng quay đều có từ thông Phi = NBS cos(omega t + phi). Vì vậy phương án $D$ đúng; các phương án còn lại sai.
}
\end{ex}

% ===== Câu 295 =====
% ID: L12C3B17-72-TL
% Mức: VDC
\dangbt{Bài toán thực tiễn về vận hành và hiệu suất máy phát điện xoay chiều}
\begin{bt}
% ID: L12C3B17-72-TL
% Mức: VDC
So sánh nhận định đúng “Hiệu suất máy phát được tính bằng công suất điện có ích chia cho công suất cơ đưa vào.” với nhận định sai “Hiệu suất máy phát bằng công suất hao phí chia công suất có ích.”; chỉ rõ căn cứ Vật lí.
\loigiai{
Cần nêu được: Hiệu suất máy phát được tính bằng công suất điện có ích chia cho công suất cơ đưa vào. Khi vận dụng phải xác định đúng dữ kiện, phương chiều nếu có, điều kiện áp dụng, hệ đơn vị và kiểm tra tính hợp lí. Nhận định “Hiệu suất máy phát bằng công suất hao phí chia công suất có ích.” sai.
}
\end{bt}

% ===== Câu 296 =====
% ID: L12C3B17-72-TLN
% Mức: VDC
\begin{ex}
% ID: L12C3B17-72-TLN
% Mức: VDC
Trong bốn phát biểu về Bài toán thực tiễn về vận hành và hiệu suất máy phát điện xoay chiều, có bao nhiêu phát biểu đúng? (Chỉ ghi một số nguyên.) (1) Khi giải cần xác định đúng dữ kiện, phương chiều, điều kiện áp dụng và đơn vị. (2) Hiệu suất máy phát bằng công suất hao phí chia công suất có ích. (3) Phải dùng đúng định luật cho tình huống. (4) Hiệu suất máy phát được tính bằng công suất điện có ích chia cho công suất cơ đưa vào.
\shortans{3}
\loigiai{
Đối chiếu với kiến thức: Hiệu suất máy phát được tính bằng công suất điện có ích chia cho công suất cơ đưa vào. Có 3 phát biểu đúng.
}
\end{ex}

% ===== Câu 297 =====
% ID: L12C3B17-72-TN
% Mức: NB
\dangbt{Tính từ thông qua khung dây quay trong từ trường}
\begin{ex}
% ID: L12C3B17-72-TN
% Mức: NB
Trong nội dung Tính từ thông qua khung dây quay trong từ trường?
\choice
{\True Khung $N$ vòng quay đều có từ thông Phi = NBS cos(omega t + phi).}
{Từ thông qua khung quay luôn không đổi và bằng NBS.}
{Có thể bỏ qua phương, chiều và đơn vị trong mọi trường hợp.}
{Kết quả không phụ thuộc dữ kiện và điều kiện áp dụng.}
\loigiai{
Khung $N$ vòng quay đều có từ thông Phi = NBS cos(omega t + phi). Vì vậy phương án $A$ đúng; các phương án còn lại sai.
}
\end{ex}

% ===== Câu 298 =====
% ID: L12C3B17-73-TN
% Mức: NB
\begin{ex}
% ID: L12C3B17-73-TN
% Mức: NB
Tình huống 3: chọn kết luận chính xác về Tính từ thông qua khung dây quay trong từ trường?
\choice
{Kết quả không phụ thuộc dữ kiện và điều kiện áp dụng.}
{\True Khung $N$ vòng quay đều có từ thông Phi = NBS cos(omega t + phi).}
{Từ thông qua khung quay luôn không đổi và bằng NBS.}
{Có thể bỏ qua phương, chiều và đơn vị trong mọi trường hợp.}
\loigiai{
Khung $N$ vòng quay đều có từ thông Phi = NBS cos(omega t + phi). Vì vậy phương án $B$ đúng; các phương án còn lại sai.
}
\end{ex}

% ===== Câu 299 =====
% ID: L12C3B17-74-TN
% Mức: TH
\begin{ex}
% ID: L12C3B17-74-TN
% Mức: TH
Nội dung nào mô tả đúng nhất Tính từ thông qua khung dây quay trong từ trường?
\choice
{Có thể bỏ qua phương, chiều và đơn vị trong mọi trường hợp.}
{Kết quả không phụ thuộc dữ kiện và điều kiện áp dụng.}
{\True Khung $N$ vòng quay đều có từ thông Phi = NBS cos(omega t + phi).}
{Từ thông qua khung quay luôn không đổi và bằng NBS.}
\loigiai{
Khung $N$ vòng quay đều có từ thông Phi = NBS cos(omega t + phi). Vì vậy phương án $C$ đúng; các phương án còn lại sai.
}
\end{ex}

% ===== Câu 300 =====
% ID: L12C3B17-75-TN
% Mức: TH
\begin{ex}
% ID: L12C3B17-75-TN
% Mức: TH
Khi phân tích Tính từ thông qua khung dây quay trong từ trường?
\choice
{Từ thông qua khung quay luôn không đổi và bằng NBS.}
{Có thể bỏ qua phương, chiều và đơn vị trong mọi trường hợp.}
{Kết quả không phụ thuộc dữ kiện và điều kiện áp dụng.}
{\True Khung $N$ vòng quay đều có từ thông Phi = NBS cos(omega t + phi).}
\loigiai{
Khung $N$ vòng quay đều có từ thông Phi = NBS cos(omega t + phi). Vì vậy phương án $D$ đúng; các phương án còn lại sai.
}
\end{ex}

% ===== Câu 301 =====
% ID: L12C3B17-76-TN
% Mức: TH
\begin{ex}
% ID: L12C3B17-76-TN
% Mức: TH
Một học sinh ôn tập Tính từ thông qua khung dây quay trong từ trường?
\choice
{\True Khung $N$ vòng quay đều có từ thông Phi = NBS cos(omega t + phi).}
{Từ thông qua khung quay luôn không đổi và bằng NBS.}
{Có thể bỏ qua phương, chiều và đơn vị trong mọi trường hợp.}
{Kết quả không phụ thuộc dữ kiện và điều kiện áp dụng.}
\loigiai{
Khung $N$ vòng quay đều có từ thông Phi = NBS cos(omega t + phi). Vì vậy phương án $A$ đúng; các phương án còn lại sai.
}
\end{ex}

% ===== Câu 302 =====
% ID: L12C3B17-77-TN
% Mức: TH
\begin{ex}
% ID: L12C3B17-77-TN
% Mức: TH
Chọn phát biểu không trái với kiến thức về Tính từ thông qua khung dây quay trong từ trường?
\choice
{Kết quả không phụ thuộc dữ kiện và điều kiện áp dụng.}
{\True Khung $N$ vòng quay đều có từ thông Phi = NBS cos(omega t + phi).}
{Từ thông qua khung quay luôn không đổi và bằng NBS.}
{Có thể bỏ qua phương, chiều và đơn vị trong mọi trường hợp.}
\loigiai{
Khung $N$ vòng quay đều có từ thông Phi = NBS cos(omega t + phi). Vì vậy phương án $B$ đúng; các phương án còn lại sai.
}
\end{ex}

% ===== Câu 303 =====
% ID: L12C3B17-78-TN
% Mức: VD
\begin{ex}
% ID: L12C3B17-78-TN
% Mức: VD
Cơ sở Vật lí đúng dùng để giải Tính từ thông qua khung dây quay trong từ trường?
\choice
{Có thể bỏ qua phương, chiều và đơn vị trong mọi trường hợp.}
{Kết quả không phụ thuộc dữ kiện và điều kiện áp dụng.}
{\True Khung $N$ vòng quay đều có từ thông Phi = NBS cos(omega t + phi).}
{Từ thông qua khung quay luôn không đổi và bằng NBS.}
\loigiai{
Khung $N$ vòng quay đều có từ thông Phi = NBS cos(omega t + phi). Vì vậy phương án $C$ đúng; các phương án còn lại sai.
}
\end{ex}

% ===== Câu 304 =====
% ID: L12C3B17-79-TN
% Mức: VD
\begin{ex}
% ID: L12C3B17-79-TN
% Mức: VD
Trong một bài thực hành về Tính từ thông qua khung dây quay trong từ trường?
\choice
{Từ thông qua khung quay luôn không đổi và bằng NBS.}
{Có thể bỏ qua phương, chiều và đơn vị trong mọi trường hợp.}
{Kết quả không phụ thuộc dữ kiện và điều kiện áp dụng.}
{\True Khung $N$ vòng quay đều có từ thông Phi = NBS cos(omega t + phi).}
\loigiai{
Khung $N$ vòng quay đều có từ thông Phi = NBS cos(omega t + phi). Vì vậy phương án $D$ đúng; các phương án còn lại sai.
}
\end{ex}

% ===== Câu 305 =====
% ID: L12C3B17-80-TN
% Mức: VD
\begin{ex}
% ID: L12C3B17-80-TN
% Mức: VD
Kết luận đúng khi vận dụng Tính từ thông qua khung dây quay trong từ trường?
\choice
{\True Khung $N$ vòng quay đều có từ thông Phi = NBS cos(omega t + phi).}
{Từ thông qua khung quay luôn không đổi và bằng NBS.}
{Có thể bỏ qua phương, chiều và đơn vị trong mọi trường hợp.}
{Kết quả không phụ thuộc dữ kiện và điều kiện áp dụng.}
\loigiai{
Khung $N$ vòng quay đều có từ thông Phi = NBS cos(omega t + phi). Vì vậy phương án $A$ đúng; các phương án còn lại sai.
}
\end{ex}

% ===== Câu 306 =====
% ID: L12C3B17-81-TN
% Mức: NB
\dangbt{Thiết lập biểu thức suất điện động xoay chiều do máy phát tạo ra}
\begin{ex}
% ID: L12C3B17-81-TN
% Mức: NB
Phát biểu nào sau đây đúng về Thiết lập biểu thức suất điện động xoay chiều do máy phát tạo ra?
\choice
{Suất điện động luôn cùng pha và bằng đúng từ thông.}
{Có thể bỏ qua phương, chiều và đơn vị trong mọi trường hợp.}
{Kết quả không phụ thuộc dữ kiện và điều kiện áp dụng.}
{\True Nếu Phi = Phi0 cos(omega t + phi) thì e = omega Phi0 sin(omega t + phi).}
\loigiai{
Nếu Phi = Phi0 cos(omega t + phi) thì e = omega Phi0 sin(omega t + phi). Vì vậy phương án $D$ đúng; các phương án còn lại sai.
}
\end{ex}

% ===== Câu 307 =====
% ID: L12C3B17-82-TN
% Mức: NB
\begin{ex}
% ID: L12C3B17-82-TN
% Mức: NB
Trong nội dung Thiết lập biểu thức suất điện động xoay chiều do máy phát tạo ra?
\choice
{\True Nếu Phi = Phi0 cos(omega t + phi) thì e = omega Phi0 sin(omega t + phi).}
{Suất điện động luôn cùng pha và bằng đúng từ thông.}
{Có thể bỏ qua phương, chiều và đơn vị trong mọi trường hợp.}
{Kết quả không phụ thuộc dữ kiện và điều kiện áp dụng.}
\loigiai{
Nếu Phi = Phi0 cos(omega t + phi) thì e = omega Phi0 sin(omega t + phi). Vì vậy phương án $A$ đúng; các phương án còn lại sai.
}
\end{ex}

% ===== Câu 308 =====
% ID: L12C3B17-83-TN
% Mức: NB
\begin{ex}
% ID: L12C3B17-83-TN
% Mức: NB
Tình huống 3: chọn kết luận chính xác về Thiết lập biểu thức suất điện động xoay chiều do máy phát tạo ra?
\choice
{Kết quả không phụ thuộc dữ kiện và điều kiện áp dụng.}
{\True Nếu Phi = Phi0 cos(omega t + phi) thì e = omega Phi0 sin(omega t + phi).}
{Suất điện động luôn cùng pha và bằng đúng từ thông.}
{Có thể bỏ qua phương, chiều và đơn vị trong mọi trường hợp.}
\loigiai{
Nếu Phi = Phi0 cos(omega t + phi) thì e = omega Phi0 sin(omega t + phi). Vì vậy phương án $B$ đúng; các phương án còn lại sai.
}
\end{ex}

% ===== Câu 309 =====
% ID: L12C3B17-84-TN
% Mức: TH
\begin{ex}
% ID: L12C3B17-84-TN
% Mức: TH
Nội dung nào mô tả đúng nhất Thiết lập biểu thức suất điện động xoay chiều do máy phát tạo ra?
\choice
{Có thể bỏ qua phương, chiều và đơn vị trong mọi trường hợp.}
{Kết quả không phụ thuộc dữ kiện và điều kiện áp dụng.}
{\True Nếu Phi = Phi0 cos(omega t + phi) thì e = omega Phi0 sin(omega t + phi).}
{Suất điện động luôn cùng pha và bằng đúng từ thông.}
\loigiai{
Nếu Phi = Phi0 cos(omega t + phi) thì e = omega Phi0 sin(omega t + phi). Vì vậy phương án $C$ đúng; các phương án còn lại sai.
}
\end{ex}

% ===== Câu 310 =====
% ID: L12C3B17-85-TN
% Mức: TH
\begin{ex}
% ID: L12C3B17-85-TN
% Mức: TH
Khi phân tích Thiết lập biểu thức suất điện động xoay chiều do máy phát tạo ra?
\choice
{Suất điện động luôn cùng pha và bằng đúng từ thông.}
{Có thể bỏ qua phương, chiều và đơn vị trong mọi trường hợp.}
{Kết quả không phụ thuộc dữ kiện và điều kiện áp dụng.}
{\True Nếu Phi = Phi0 cos(omega t + phi) thì e = omega Phi0 sin(omega t + phi).}
\loigiai{
Nếu Phi = Phi0 cos(omega t + phi) thì e = omega Phi0 sin(omega t + phi). Vì vậy phương án $D$ đúng; các phương án còn lại sai.
}
\end{ex}

% ===== Câu 311 =====
% ID: L12C3B17-86-TN
% Mức: TH
\begin{ex}
% ID: L12C3B17-86-TN
% Mức: TH
Một học sinh ôn tập Thiết lập biểu thức suất điện động xoay chiều do máy phát tạo ra?
\choice
{\True Nếu Phi = Phi0 cos(omega t + phi) thì e = omega Phi0 sin(omega t + phi).}
{Suất điện động luôn cùng pha và bằng đúng từ thông.}
{Có thể bỏ qua phương, chiều và đơn vị trong mọi trường hợp.}
{Kết quả không phụ thuộc dữ kiện và điều kiện áp dụng.}
\loigiai{
Nếu Phi = Phi0 cos(omega t + phi) thì e = omega Phi0 sin(omega t + phi). Vì vậy phương án $A$ đúng; các phương án còn lại sai.
}
\end{ex}

% ===== Câu 312 =====
% ID: L12C3B17-87-TN
% Mức: TH
\begin{ex}
% ID: L12C3B17-87-TN
% Mức: TH
Chọn phát biểu không trái với kiến thức về Thiết lập biểu thức suất điện động xoay chiều do máy phát tạo ra?
\choice
{Kết quả không phụ thuộc dữ kiện và điều kiện áp dụng.}
{\True Nếu Phi = Phi0 cos(omega t + phi) thì e = omega Phi0 sin(omega t + phi).}
{Suất điện động luôn cùng pha và bằng đúng từ thông.}
{Có thể bỏ qua phương, chiều và đơn vị trong mọi trường hợp.}
\loigiai{
Nếu Phi = Phi0 cos(omega t + phi) thì e = omega Phi0 sin(omega t + phi). Vì vậy phương án $B$ đúng; các phương án còn lại sai.
}
\end{ex}

% ===== Câu 313 =====
% ID: L12C3B17-88-TN
% Mức: VD
\begin{ex}
% ID: L12C3B17-88-TN
% Mức: VD
Cơ sở Vật lí đúng dùng để giải Thiết lập biểu thức suất điện động xoay chiều do máy phát tạo ra?
\choice
{Có thể bỏ qua phương, chiều và đơn vị trong mọi trường hợp.}
{Kết quả không phụ thuộc dữ kiện và điều kiện áp dụng.}
{\True Nếu Phi = Phi0 cos(omega t + phi) thì e = omega Phi0 sin(omega t + phi).}
{Suất điện động luôn cùng pha và bằng đúng từ thông.}
\loigiai{
Nếu Phi = Phi0 cos(omega t + phi) thì e = omega Phi0 sin(omega t + phi). Vì vậy phương án $C$ đúng; các phương án còn lại sai.
}
\end{ex}

% ===== Câu 314 =====
% ID: L12C3B17-89-TN
% Mức: VD
\begin{ex}
% ID: L12C3B17-89-TN
% Mức: VD
Trong một bài thực hành về Thiết lập biểu thức suất điện động xoay chiều do máy phát tạo ra?
\choice
{Suất điện động luôn cùng pha và bằng đúng từ thông.}
{Có thể bỏ qua phương, chiều và đơn vị trong mọi trường hợp.}
{Kết quả không phụ thuộc dữ kiện và điều kiện áp dụng.}
{\True Nếu Phi = Phi0 cos(omega t + phi) thì e = omega Phi0 sin(omega t + phi).}
\loigiai{
Nếu Phi = Phi0 cos(omega t + phi) thì e = omega Phi0 sin(omega t + phi). Vì vậy phương án $D$ đúng; các phương án còn lại sai.
}
\end{ex}

% ===== Câu 315 =====
% ID: L12C3B17-90-TN
% Mức: VD
\begin{ex}
% ID: L12C3B17-90-TN
% Mức: VD
Kết luận đúng khi vận dụng Thiết lập biểu thức suất điện động xoay chiều do máy phát tạo ra?
\choice
{\True Nếu Phi = Phi0 cos(omega t + phi) thì e = omega Phi0 sin(omega t + phi).}
{Suất điện động luôn cùng pha và bằng đúng từ thông.}
{Có thể bỏ qua phương, chiều và đơn vị trong mọi trường hợp.}
{Kết quả không phụ thuộc dữ kiện và điều kiện áp dụng.}
\loigiai{
Nếu Phi = Phi0 cos(omega t + phi) thì e = omega Phi0 sin(omega t + phi). Vì vậy phương án $A$ đúng; các phương án còn lại sai.
}
\end{ex}

% ===== Câu 316 =====
% ID: L12C3B17-91-TN
% Mức: NB
\dangbt{Xác định biên độ, tần số, chu kì và pha của suất điện động xoay chiều}
\begin{ex}
% ID: L12C3B17-91-TN
% Mức: NB
Phát biểu nào sau đây đúng về Xác định biên độ, tần số, chu kì và pha của suất điện động xoay chiều?
\choice
{Biên độ suất điện động không phụ thuộc tốc độ quay.}
{Có thể bỏ qua phương, chiều và đơn vị trong mọi trường hợp.}
{Kết quả không phụ thuộc dữ kiện và điều kiện áp dụng.}
{\True Biên độ suất điện động là E0 = NBS omega và tần số f = omega/(2pi).}
\loigiai{
Biên độ suất điện động là E0 = NBS omega và tần số f = omega/(2pi). Vì vậy phương án $D$ đúng; các phương án còn lại sai.
}
\end{ex}

% ===== Câu 317 =====
% ID: L12C3B17-92-TN
% Mức: NB
\begin{ex}
% ID: L12C3B17-92-TN
% Mức: NB
Trong nội dung Xác định biên độ, tần số, chu kì và pha của suất điện động xoay chiều?
\choice
{\True Biên độ suất điện động là E0 = NBS omega và tần số f = omega/(2pi).}
{Biên độ suất điện động không phụ thuộc tốc độ quay.}
{Có thể bỏ qua phương, chiều và đơn vị trong mọi trường hợp.}
{Kết quả không phụ thuộc dữ kiện và điều kiện áp dụng.}
\loigiai{
Biên độ suất điện động là E0 = NBS omega và tần số f = omega/(2pi). Vì vậy phương án $A$ đúng; các phương án còn lại sai.
}
\end{ex}

% ===== Câu 318 =====
% ID: L12C3B17-93-TN
% Mức: NB
\begin{ex}
% ID: L12C3B17-93-TN
% Mức: NB
Tình huống 3: chọn kết luận chính xác về Xác định biên độ, tần số, chu kì và pha của suất điện động xoay chiều?
\choice
{Kết quả không phụ thuộc dữ kiện và điều kiện áp dụng.}
{\True Biên độ suất điện động là E0 = NBS omega và tần số f = omega/(2pi).}
{Biên độ suất điện động không phụ thuộc tốc độ quay.}
{Có thể bỏ qua phương, chiều và đơn vị trong mọi trường hợp.}
\loigiai{
Biên độ suất điện động là E0 = NBS omega và tần số f = omega/(2pi). Vì vậy phương án $B$ đúng; các phương án còn lại sai.
}
\end{ex}

% ===== Câu 319 =====
% ID: L12C3B17-94-TN
% Mức: TH
\begin{ex}
% ID: L12C3B17-94-TN
% Mức: TH
Nội dung nào mô tả đúng nhất Xác định biên độ, tần số, chu kì và pha của suất điện động xoay chiều?
\choice
{Có thể bỏ qua phương, chiều và đơn vị trong mọi trường hợp.}
{Kết quả không phụ thuộc dữ kiện và điều kiện áp dụng.}
{\True Biên độ suất điện động là E0 = NBS omega và tần số f = omega/(2pi).}
{Biên độ suất điện động không phụ thuộc tốc độ quay.}
\loigiai{
Biên độ suất điện động là E0 = NBS omega và tần số f = omega/(2pi). Vì vậy phương án $C$ đúng; các phương án còn lại sai.
}
\end{ex}

% ===== Câu 320 =====
% ID: L12C3B17-95-TN
% Mức: TH
\begin{ex}
% ID: L12C3B17-95-TN
% Mức: TH
Khi phân tích Xác định biên độ, tần số, chu kì và pha của suất điện động xoay chiều?
\choice
{Biên độ suất điện động không phụ thuộc tốc độ quay.}
{Có thể bỏ qua phương, chiều và đơn vị trong mọi trường hợp.}
{Kết quả không phụ thuộc dữ kiện và điều kiện áp dụng.}
{\True Biên độ suất điện động là E0 = NBS omega và tần số f = omega/(2pi).}
\loigiai{
Biên độ suất điện động là E0 = NBS omega và tần số f = omega/(2pi). Vì vậy phương án $D$ đúng; các phương án còn lại sai.
}
\end{ex}

% ===== Câu 321 =====
% ID: L12C3B17-96-TN
% Mức: TH
\begin{ex}
% ID: L12C3B17-96-TN
% Mức: TH
Một học sinh ôn tập Xác định biên độ, tần số, chu kì và pha của suất điện động xoay chiều?
\choice
{\True Biên độ suất điện động là E0 = NBS omega và tần số f = omega/(2pi).}
{Biên độ suất điện động không phụ thuộc tốc độ quay.}
{Có thể bỏ qua phương, chiều và đơn vị trong mọi trường hợp.}
{Kết quả không phụ thuộc dữ kiện và điều kiện áp dụng.}
\loigiai{
Biên độ suất điện động là E0 = NBS omega và tần số f = omega/(2pi). Vì vậy phương án $A$ đúng; các phương án còn lại sai.
}
\end{ex}

% ===== Câu 322 =====
% ID: L12C3B17-97-TN
% Mức: TH
\begin{ex}
% ID: L12C3B17-97-TN
% Mức: TH
Chọn phát biểu không trái với kiến thức về Xác định biên độ, tần số, chu kì và pha của suất điện động xoay chiều?
\choice
{Kết quả không phụ thuộc dữ kiện và điều kiện áp dụng.}
{\True Biên độ suất điện động là E0 = NBS omega và tần số f = omega/(2pi).}
{Biên độ suất điện động không phụ thuộc tốc độ quay.}
{Có thể bỏ qua phương, chiều và đơn vị trong mọi trường hợp.}
\loigiai{
Biên độ suất điện động là E0 = NBS omega và tần số f = omega/(2pi). Vì vậy phương án $B$ đúng; các phương án còn lại sai.
}
\end{ex}

% ===== Câu 323 =====
% ID: L12C3B17-98-TN
% Mức: VD
\begin{ex}
% ID: L12C3B17-98-TN
% Mức: VD
Cơ sở Vật lí đúng dùng để giải Xác định biên độ, tần số, chu kì và pha của suất điện động xoay chiều?
\choice
{Có thể bỏ qua phương, chiều và đơn vị trong mọi trường hợp.}
{Kết quả không phụ thuộc dữ kiện và điều kiện áp dụng.}
{\True Biên độ suất điện động là E0 = NBS omega và tần số f = omega/(2pi).}
{Biên độ suất điện động không phụ thuộc tốc độ quay.}
\loigiai{
Biên độ suất điện động là E0 = NBS omega và tần số f = omega/(2pi). Vì vậy phương án $C$ đúng; các phương án còn lại sai.
}
\end{ex}

% ===== Câu 324 =====
% ID: L12C3B17-99-TN
% Mức: VD
\begin{ex}
% ID: L12C3B17-99-TN
% Mức: VD
Trong một bài thực hành về Xác định biên độ, tần số, chu kì và pha của suất điện động xoay chiều?
\choice
{Biên độ suất điện động không phụ thuộc tốc độ quay.}
{Có thể bỏ qua phương, chiều và đơn vị trong mọi trường hợp.}
{Kết quả không phụ thuộc dữ kiện và điều kiện áp dụng.}
{\True Biên độ suất điện động là E0 = NBS omega và tần số f = omega/(2pi).}
\loigiai{
Biên độ suất điện động là E0 = NBS omega và tần số f = omega/(2pi). Vì vậy phương án $D$ đúng; các phương án còn lại sai.
}
\end{ex}
